\documentclass[11pt,a4paper,twoside]{book}
\usepackage[T1]{fontenc}
\usepackage{lmodern}
\usepackage[utf8]{inputenc} 
\usepackage{amsmath, amsthm, amsfonts, amssymb, xfrac}
\usepackage{arydshln} %Linea dotted in array
\usepackage[italian]{babel}
\usepackage{titling} %Funzioni per modificare maketitle
\usepackage{titlesec} %Ridefinizione semplice di titoli 
\usepackage{float} %Comando H per posizione di float
\usepackage[left=2.5cm, right=2cm, top=2cm, bottom=2cm]{geometry} %Gestione margini pagina
\usepackage{perpage} %Permette di resettare una numerazione per pagina
\usepackage{calligra}
\usepackage{fancyhdr} %Lo uso per avere titolo e autore running
\usepackage[hidelinks]{hyperref} %lo uso per avere l'indice cliccabiile
\usepackage{graphicx, wrapfig} %Immagini
\usepackage[shortlabels]{enumitem} %Enumerate con diversi label, l'argomento mi permette di usare (a) al posto di label=(\alph*) etc
\usepackage{ulem}

\MakePerPage{footnote} %Resetta la numerazione delle footnote per pagina

%Autore-------------------------------------------
\author{Marco Accerenzi}
\date{}
\title{Appunti di fisica matematica}


%Titleformat-----------------------------------------------
\renewcommand{\thesection}{\arabic{section}}
\renewcommand{\thesubsection}{\thesection.\arabic{subsection}}
\titleformat{\section}{\huge\bfseries}{{\Large \thesection}}{0.5em}{}
\titleformat{\subsection}{\Large\bfseries}{{\large \thesubsection}}{0.5em}{}
\titleformat{\paragraph}[runin]{\normalsize\bfseries}{{}}{0}{}

%Pagestyle pagine normali--------------------------------
\pagestyle{fancy} %Serve per poter modificare le pagine con fancyhdr facilmente
\renewcommand{\sectionmark}[1]{ \markright{#1}{} }

\lhead[]{\small Marco Accerenzi}  %Autore a sinistra nelle pagine dispari
\chead[\small \rightmark]{} %Capitolo al centro nelle pagine pari
\rhead[]{\small Appunti di fisica matematica} %Titolo a destra nelle pagine dispari
\fancyfoot[LE,RO]{\large\thepage} %Numero pagina grande a sinistra in basso nelle pagine dispari
\fancyfoot[C]{}
\renewcommand{\headrulewidth}{0.25pt}
\renewcommand{\footrulewidth}{0.3pt}

%Modifica pagestyle pagine speciali-------------------------------------------
\fancypagestyle{plain}{%Ridefinisco il pagestyle della prima pagina dei chapter
\fancyhf{} %Toglie tutti i footer ed header
\fancyfoot[LO,RE]{\thepage} 
\renewcommand{\headrulewidth}{0pt}
\renewcommand{\footrulewidth}{0pt}
}

%Tcolorbox per i teoremi, esempi ed esercizi---------------------------
\usepackage[skins, breakable, theorems]{tcolorbox}

%Box per proposizione
\newtcbtheorem[number within=section]{Proposizione}{Proposizione}{
	enhanced,
	boxrule=0pt,frame hidden,
	borderline west={2.5pt}{0pt}{black},
	colback=yellow!1!white,
	sharp corners,
	coltitle=black,
	attach title to upper={\noindent},
	fonttitle=\bfseries\large,
}{thm}

%Box per Definizione
\newtcbtheorem[use counter from=Proposizione]{Definizione}{Definizione}{
enhanced,
	boxrule=0pt,frame hidden,
	borderline west={2.5pt}{0pt}{black},
	colback=yellow!1!white,
	sharp corners,
	coltitle=black,
	attach title to upper={\noindent},
	fonttitle=\bfseries\large,
}{defi}

%Box per Lemma
\newtcbtheorem[use counter from=Proposizione]{Lemma}{Lemma}{
enhanced,
	boxrule=0pt,frame hidden,
	borderline west={2.5pt}{0pt}{black},
	colback=yellow!1!white,
	sharp corners,
	coltitle=black,
	attach title to upper={\noindent},
	fonttitle=\bfseries\large,
}{lemm}

%Box per Corollario
\newtcbtheorem[use counter from=Proposizione]{Corollario}{Corollario}{
enhanced,
	boxrule=0pt,frame hidden,
	borderline west={2.5pt}{0pt}{black},
	colback=yellow!1!white,
	sharp corners,
	coltitle=black,
	attach title to upper={\noindent},
	fonttitle=\bfseries\large,
}{Cor}



%Ambienti per proposizione, definizione, lemma e corollario
\newenvironment{proposizione}[1][\empty]{\begin{Proposizione}{
\ifx\empty#1\relax\else\mdseries\itshape#1\\\fi
}{}}{\end{Proposizione}}  
\newenvironment{definizione}[1][\empty]{\begin{Definizione}{
\ifx\empty#1\relax\else\mdseries\itshape#1\\\fi
}{}}{\end{Definizione}}
\newenvironment{lemma}[1][\empty]{\begin{Lemma}{
\ifx\empty#1\relax\else\mdseries\itshape#1\\\fi
}{}}{\end{Lemma}}
\newenvironment{corollario}[1][\empty]{\begin{Corollario}{
\ifx\empty#1\relax\else\mdseries\itshape#1\\\fi
}{}}{\end{Corollario}}

%Box per pezzi generici importanti-------------------------
\newenvironment{importante}[0]{ 
  \begin{center}\begin{tcolorbox}[enhanced,
  before skip=2mm,after skip=3mm,
  boxrule=0.4pt,left=5mm,right=2mm,top=1mm,bottom=1mm,
  colback=yellow!50,
  colframe=yellow!20!black,
  sharp corners,arc=3mm,
  underlay={
  	\path[fill=red!50!black,draw=none] (interior.south west) rectangle 			node[white]{\Huge\bfseries !} ([xshift=4mm]interior.north west);
  },
  drop fuzzy shadow]
  }{
 \normalsize \end{tcolorbox}\end{center}
}

%Box per esempi----------------------
\newenvironment{esempio}[1][Esempio]{ 
  \begin{center}\begin{tcolorbox}[
  title={#1:\\},
  width=0.9\linewidth,
  flush right,
  enhanced jigsaw, 
  colback=white, 
  frame hidden, 
  sharpish corners, arc=3mm,
  borderline west = {1pt}{0pt}{gray},
  borderline south = {1pt}{0pt}{gray},
  coltitle=black, 
  fonttitle=\small\bfseries, 
  attach title to upper, 
  breakable,
  top=2.5pt, bottom=0.3pt]\small
  }{
 \normalsize \end{tcolorbox}\end{center}
}

%Ambiente proof con riga a sinistra---------------------------
\newenvironment{dimo}[1][\proofname]{
	\begin{tcolorbox}[
	breakable, 
	enhanced, 
	interior hidden, 
	frame hidden, 
	borderline west = {1pt}{0pt}{gray}, 
	top=0pt, bottom=0pt, left=0pt]
\begin{proof}[#1:~]
}
{
\end{proof}
\end{tcolorbox}
}

%Stili teorema per soluzione esercizi--------------------------------------
\newtheoremstyle{esss}{\topsep}{\topsep}{\normalfont}{0pt}{\bfseries}{}{ }{\thmname{#1}\thmnumber{ #2}\thmnote{ #3}}
\theoremstyle{esss}

\newtheorem{esercizio}{Esercizio}[section]

\newtheoremstyle{dotless}{\topsep}{\topsep}{\normalfont}{0pt}{\bfseries}{}{ }{\thmname{#1}\thmnumber{ #2}\thmnote{ #3}}
\theoremstyle{dotless}

\newtheorem*{ris}{}
\newenvironment{risoluzione}[1][]{  
  \begin{ris}[\bfseries Risoluzione #1]$ $\\
  }{
  \end{ris}
}

%Miste------------------------------------------------------------------
\numberwithin{equation}{section} %Numerazione equazione per sezione
\renewcommand{\thefootnote}{\fnsymbol{footnote}} 
\hypersetup{linktoc=all}
\newcommand{\sectionbreak}{\clearpage}

%Indice--------------------------------------------------------
\usepackage{tocloft}
\renewcommand{\cftpartfont}{\large\bfseries}
\renewcommand{\cftsecfont}{\bfseries}



\begin{document}%-------------------------------------------------------------

\frontmatter
\pagestyle{empty}

%Copertina----------------------------------------------------------------------
\begin{titlepage}
	\begin{center}

	\begin{center}
		\centering
		Anno accademico 2018-2019\par
		\vspace{.05\textheight}
	\end{center}	

		\vspace{.1\textheight}
		{\LARGE\scshape \textsc{Marco Accerenzi}\par}
		\vspace{.1\textheight}
		\textsc{}\par
		\vspace{.05\textheight}
		{\Huge \textsc{Istituzioni di Fisica Matematica}\par}
		\vspace{.05\textheight}
		{\large Riassunto degli appunti\par}
		\vfill
		{\large
Questo fascicoletto è un riassunto dei miei appunti personali del corso di \textit{Istituzioni di fisica matematica} dell'anno accademico 2018-2019 tenuto dal prof. Fassò. \\Gli esempi riportati sono quelli visti a lezione ed alcuni esercizi che ho risolto personalmente.
  }
  \vspace{2cm} 
\end{center}
\end{titlepage}

%Inizio------------------------------------------------------------------------------
\tableofcontents

%Parte 1-----------------------------------------------------------------
\cleardoublepage
\mainmatter\setcounter{page}{1}
\pagestyle{fancy}
\addcontentsline{toc}{part}{Teoria sui sistemi dinamici}
\thispagestyle{plain}
	\begin{center}
		\vspace{.1\textheight}
		{\LARGE\scshape \textsc{Prima parte}\par}
		\vspace{.1\textheight}
		\textsc{}\par
		\vspace{.05\textheight}
		{\Huge \textsc{Teoria sui sistemi Dinamici}\par}
		\vfill

\end{center}
\fancyfoot[C]{\small Teoria su Sistemi Dinamici}
%------------------------------------------------------------------

\section{Equazioni differenziali}
Lo studio della dinamica si basa sulle equazioni differenziali, per lo più useremo equazioni autonome
\[
\dot{z}=X(z)
\]
ma talvolta anche non autonome, lo specificheremo dove necessario. Dove non specificato altrimenti consideriamo soltanto equazioni autonome.
 
Chiameremo $\Omega \subseteq \mathbb{R}^n $ un aperto di $\mathbb{R}^n$ e $ z\in\Omega $ un suo punto; useremo spesso curve differenziabili su $\Omega$: funzioni $\gamma:J\rightarrow\Omega$ dove $J\subseteq\mathbb{R}$ e la $\gamma$ è differenziabile.

\begin{enumerate}[i.]
\item Sia $\gamma:J\rightarrow\Omega$ una curva differenziabile definita su un intorno $J\subseteq \mathbb{R}$ di $0$ con $\gamma(0)=z\in\Omega$.
\item Il vettore $\dot{\gamma}(0)= \dfrac{\partial\gamma}{\partial t}(0)$ sarà un vettore tangente a $\Omega$ in $z$ e vivrà nello spazio tangente
\[
T_z\Omega= \left\{ \dot{\gamma}(0) \text{ per } \gamma:I_\gamma\rightarrow\Omega , \> \gamma(0)=z  \right\}
\]
$T_z\Omega\subseteq\mathbb{R}^n$ e per semplicità assumeremo spesso $T_z\Omega \approx \mathbb{R}^n$.

\item Lo spazio tangente affine è lo spazio tangente ad $\Omega$ ``incollato'' al punto di $\Omega$ su cui è calcolato, in pratica considereremo l'origine di $T_z\Omega$ sovrapposta a $z$: lo possiamo pensare come $z+T_z\Omega$.
\end{enumerate}

\begin{definizione}
Chiamiamo fibrato tangente $T\Omega$ di $\Omega$ l'insieme delle coppi dei punti di $\Omega$ e dei vettori a loro tangenti, cioè
\[
T\Omega=\left\{ (z,v) : z\in\Omega , \> v\in T_z\Omega   \right\}
\]
\end{definizione}

A volte, sempre per semplicità, identificheremo $T\Omega \approx \Omega \times \mathbb{R}^n$.

Un campo vettoriale su $\Omega$ è una mappa che associa a  $z\in\Omega$ un vettore contenuto in $T_z\Omega$, possiamo cioè definire $X:\Omega\rightarrow T\Omega$ con $X(z)=(z,\tilde{X}(z))$ e $\tilde{X}:\Omega\rightarrow T_z\Omega$: $X(z)$ conserva l'informazione sul punto in cui è stato applicato e $\tilde{X}(z)$ è il campo vettoriale.

\begin{definizione}
Una curva integrale di un campo vettoriale $X$ è una curva $\gamma:J\rightarrow\Omega$ che soddisfi
\[
\dot{z}=\gamma '(t)= X(z) = X(\gamma(t)) \> \> \> \forall t\in J
\]
\end{definizione}

Le curve integrali di $X$ corrisponderanno (biunivocamente) alle soluzioni dell'equazione differenziale $\dot{z}=X(z)$.

Da qui in avanti scriveremo $z(t)$ per indicare una soluzione dell'equazione differenziale al variare del tempo, sottointendendo il concetto di curva integrale.

Chiamiamo \textbf{soluzioni massimali} le soluzioni $\gamma:J\rightarrow\Omega$ di un' equazione differenziale tali per cui non è possibile estendere l'intervallo di definizione, cioè una soluzione $\gamma:J\rightarrow\Omega$ è massimale se $\forall I : J\subset I$ una $\gamma:I\rightarrow\Omega$ non è più una soluzione.

\begin{definizione}
Un problema di Cauchy per un'equazione differenziale del primo ordine consiste nella ricerca della soluzione (massimale) dell'equazione differenziale passante per un punto $z_0$ dato ad un istante $t_0$.
\end{definizione}

Sappiamo che sotto diverse ipotesi sul campo vettoriale dell'equazione il problema di Cauchy avrà esistenza ed unicità per le soluzioni, globali o locali a seconda delle ipotesi. Questo significa che due soluzioni massimali diverse del problema di Cauchy, rispettate le ipotesi, non potranno intersecarsi e che per un determinato punto di $\Omega$ passerà sempre una ed una sola soluzione.

\begin{definizione}
Chiamiamo orbita di una soluzione $\gamma:J\rightarrow\Omega$ di un problema di Cauchy la sua immagine orientata in $\Omega$. 
\end{definizione}

L'insieme di tutte le orbite si chiama \textbf{ritratto in fase}, ovviamente ci si limita a rappresentare le orbite più significative quando si disegna un ritratto in fase.

\begin{esempio}[Equazione logistica]\flushleft
\begin{minipage}[H]{0.6\linewidth}
L'equazione $\dot{z}=kz-hz^2$ con $z\in\mathbb{R}$ e i parametri $h,k>0$ è chiamata equazione logistica. Vogliamo disegnarne il ritratto in fase e studiarla senza calcolarne l'integrale generale.

Vediamo che la velocità sarà nulla in $z=0$ e in $z_1=\sfrac{k}{h}>0$, per $0<z<z_1$ la velocità è positiva mentre per $z<0$ e $z>z_1$ la velocità è negativa. Disegniamo e orientiamo nel senso dei tempi positivi il grafico di $(z,\dot{z})$ e da questo ricaviamo il ritratto in fase. Vedremo più avanti che $z_1$ è un equilibrio attrattivo cioè un punto a cui tende la $z(t)$ per $t\rightarrow+\infty$ mentre l'origine è repulsiva.

L'introduzione di un prelievo a velocità costante $\dot{z}=kz-hz^2-p$ abbassa di $p>0$ il grafico $(z,\dot{z})$. 
\end{minipage}
\begin{minipage}[H]{0.35\linewidth}
\centering
\includegraphics[width=\linewidth]{img/ritrattoLogistica.png}
\end{minipage}


Per prelievi sufficientemente piccoli $p,\overline{p}$ viene ristretta la zona con $z(t)$ crescente ed abbassato il valore $z_1$ attrattivo ma per prelievi troppo forti $p>\overline{p}$ tutta la curva va sotto lo zero, l'equilibrio $z_1$ sparisce e tutte le orbite con valori iniziali $z_0$ positivi tenderanno nei tempi positivi a $z\rightarrow 0$.
\end{esempio}

Dall'unicità delle soluzioni segue che soluzioni diverse non si intersecano mai.

\begin{definizione}\label{defFlusso1}
La mappa $\mathbb{R}\times\Omega \rightarrow \Omega$ definita $\Phi(t,z_0)=z(t;z_0,t_0)$ si chiama flusso dell'equazione differenziale e fornisce, al variare di $z_0$, tutte le soluzioni dell'equazione differenziale.
\end{definizione}

\begin{definizione}
Un sottoinsieme $M$ dello spazio delle fasi $\Omega$ si dice invariante per l'equazione differenziale $\dot{z}=X(z)$, o equivalentemente invariante per il flusso di $X$, se ogni soluzione dell'equazione differenziale con dato iniziale appartenente ad $M$ appartiene ad $M$ per tutti i tempi
\[
\forall z(t;z_0) \text{ con } z_0\in M \>\>\> z(t;z_0)\in M \> \forall t
\]
\end{definizione}
Un invariante è un sottospazio dello spazio delle fasi.


\subsection{Equazioni differenziali autonome}
Indicheremo con $z(t;z_0,t_0)$ la soluzione passante per $z_0\in\Omega$ al tempo $t_0\in J$ del problema di Cauchy.

Se consideriamo equazioni differenziali autonome, come abbiamo implicitamente fatto fino ad ora, allora sappiamo che due soluzioni che passano ad istanti diversi per lo stesso punto di $\Omega$ differiscono solo per una traslazione temporale, cioè
\[ z(t;z_0,t_0)=z(t-t_0;z_0,0) \]
per questo motivo quando tratteremo equazioni differenziali autonome ometteremo di indicare $t_0$ nella definizione della soluzione e considereremo semplicemente $t_0=0$.

\begin{proposizione}
Le orbite di un'equazione differenziale autonoma non si intersecano mai.
\end{proposizione}

\begin{esempio}[Esempi di equazioni autonome scalari]
\begin{itemize}
\item Per $z\in\mathbb{R}$ e $\dot{z}=X(z)=kz$ lo spazio delle fasi è $\Omega=\mathbb{R}$, le soluzioni sono $z(t;z_0,0)=z(t;z_0)=e^{kt}z_0$ e fissato un punto $z_0$ la traslata temporale della soluzione è $z(t;z_0,t_0)=e^{k(t-t_0)}z_0$. Le soluzioni sono definite su tutto $\mathbb{R}$. 
\item Per $z\in\mathbb{R}$ e $\dot{z}=kz^2$ lo spazio delle fasi è $\mathbb{R}$, le soluzioni sono $z(t;z_0)=\dfrac{z_0}{1-ktz_0}$. Le soluzioni massimali hanno un intervallo di esistenza che dipende da $z_0$ poichè non sono definite in $t=\sfrac{1}{kz_0}$ 
\end{itemize}
\end{esempio}


\subsection{Equazioni differenziali del secondo ordine}
Nelle equazioni differenziali del secondo ordine cambiamo leggermente la notazioni: useremo\\ $x\in C \subseteq\mathbb{R}^n , \> \dot{x}\in T_x C \approx\mathbb{R}^n$, $Y:TC\rightarrow TC$.

Un'equazione differenziale del secondo ordine $\ddot{x}=Y(x,\dot{x})$ si potrà sempre ricondurre a una del primo ordine definendo il sistema di equazioni lineari

\begin{equation}\label{equazione secondo ordine come sistema}
\begin{cases}
\dot{x}=v \\
\dot{v}=Y(x,v)
\end{cases}
\end{equation}

Riscrivendolo come equazione del primo ordine possiamo passare alla notazione usata precedentemente definendo un singolo campo vettoriale dovremo definire $X(x,v):=(v,Y(x,v))$ e quindi dovremo ritrovare i risultati della sezione precedente per il campo $X(x,v)=(v_1,...,v_n,Y_1(x,v),...,Y_n(x,v))$.\\
Notiamo che i punti su cui agisce il campo vettoriale sono i $(x,v)$ contenuti in $T C \approx C \times \mathbb{R}^n$. 

\begin{proposizione} Data  $\ddot{x}=Y \left( x,\dot{x} \right)$ e il sistema \eqref{equazione secondo ordine come sistema} valgono:
\begin{itemize}
\item Per $t\mapsto x(t)$ soluzione dell'equazione del secondo ordine la corrispondente \[t\mapsto \left(x(t), \dfrac{\partial x}{\partial t}(t) \right) \in TC\] è una soluzione del sistema del primo ordine \eqref{equazione secondo ordine come sistema}.
\item Se $t \mapsto (x(t),v(t))$ è una soluzione del sistema del primo ordine con $v(t)=\dfrac{\partial x}{\partial t} (t)$ allora $t\mapsto x(t)$ è una soluzione dell'equazione del secondo ordine.
\end{itemize}
\end{proposizione}

Possiamo quindi chiamare spazio delle fasi $\Omega$ dell'equazione del secondo ordine il fibrato tangente $TC$ e le orbite saranno le immagini delle soluzioni $t\mapsto(x(t),v(t))$ e non delle solo $t\mapsto x(t)$.

\begin{definizione}
Per un'equazione differenziale del secondo ordine $\ddot{x}=Y(x,\dot{x})$ con $x\in C$ chiamiamo $C$ spazio delle configurazioni. Le immagini orientate in $C$ delle soluzioni $t\mapsto x(t)$ dell'equazione si diranno configurazioni. 
\end{definizione}



Le orbite, contenute in uno spazio $TC\approx\mathbb{R}^{2n}$, non si intersecano invece le configurazioni potranno intersecarsi in $C$.


\begin{esempio}[Oscillatore e Repulsore armonico]
L'equazione differenziale $\ddot{x}=kx$ con $k\in\mathbb{R}^2$ è l'equazione per oscillatori e repulsori armonici.


\begin{itemize} 
\item Per $k<0$ l'equazione rappresenta un oscillatore armonico. Possiamo riscriverla in modo più esplicito come $\begin{cases}\dot{x}=v \\ \dot{v}=-\omega^2 x  \end{cases}$.

Questa equazione è facilmente integrabile e l'integrale generale è 
\begin{equation}
x(t)=A\cos{\omega t} +B\sin{\omega t}
\end{equation}
\[
v(t)=-A\sin{\omega_t}\omega+B\omega\cos{\omega t}
\]
Volendole esprimere in funzione dei dati iniziali vediamo con $x(0),v(0)$ che $A=x_0,\, B=\dfrac{v_0}{\omega}$. Quindi l'integrale generale è
\[
\left(\begin{array}{c} x(t) \\ v(t) \end{array}\right) = \left(\begin{array}{cc} \cos{\omega t} & \dfrac{\sin{\omega t}}{\omega} \\ -\omega\sin{\omega t} & \cos{\omega t} \end{array}\right) \left(\begin{array}{c} x_0 \\ v_0  \end{array}\right)
\]
che a meno di fattori $\omega$ è una matrice di rotazione oraria di angolo $\omega t$, quindi le orbite quindi sono ellissi centrate nell'origine percorse in senso orario.
L'origine è uno zero del campo e quindi è un'orbita a sè (equilibrio).

\item Per $k>0$ l'equazione rappresenta un repulsore armonico. Possiamo riscriverla in modo più esplicito come $\begin{cases}\dot{x}=v \\ \dot{v}=\omega^2 x  \end{cases}$.

Anche questa equazione, come quella dell'oscillatore armonico, si può integrare con le tecniche normali e avrà integrale generale
\begin{equation}
x(t)=Ae^{\omega t} + Be^{-\omega t}
\end{equation}
esprimendola in funzione dei dati iniziali diventa
\[
x(t)=\sfrac{1}{2} \left(x_0 +\dfrac{v_0}{\omega} \right)e^{\omega t} +\sfrac{1}{2} \left(x_0 - \dfrac{v_0}{\omega}\right)e^{-\omega t}
\]
\[
v(t)=\dfrac{\omega}{2}\left( x_0+\dfrac{v_0}{\omega}\right) e^{\omega t} -\dfrac{\omega}{2}\left(x_0-\dfrac{v_0}{\omega}\right)e^{-\omega t}
\]

Notiamo che oltre all'equilibrio nell'origine abbiamo altre $4$ orbite speciali, infatti per $v_0=\pm\omega x_0$ le orbite diventano $x(t)=x_0e^{\pm\omega t}$ $v(t)=v_0e^{\mp\omega t}$.
Queste orbite nel piano $(x,v)$ sono date dalle due rette $v=\pm x$, poi divise in orbite dall'origine.

Sulla retta $x=y$ i moti si allontaneranno dall'origine nei tempi positivi mentre sulla retta $y=-x$ i moti per i tempi positivi tenderanno ad andare nell'origine.

A partire da queste orbite speciali possiamo poi disegnare tutte le altre orbite, che per l'unicità saranno delimitate da queste, il verso di percorrenza sarà facilmente individuabile graficamente dalle rette in quanto determinato dal termine dominante dell'equazione dell'orbita che sarà lo stesso della retta più vicina punto a punto.

In realtà a partire dalla legge oraria delle orbite possiamo abbastanza facilmente eliminare il tempo e trovare che le orbite sono soluzioni di 
\[
v^2(t)-\omega^2 x^2(t) = v_0^2 - \omega^2 x_0^2
\]
e quindi sono le iperboli di asintoti $v=\pm x$.

\end{itemize}
\end{esempio}


\section{Equilibri attrattivi e repulsivi}

\begin{definizione}
Chiamiamo equilibrio di una $\dot{z}=X(z)$ un punto del suo spazio delle fasi $\overline{z}$ per cui le orbite passanti per $\overline{z}$ rimangano ``ferme'' in $\overline{z}$ per tutti i tempi:
\[
z(t, \overline{z})=\overline{z} \> \> \> \forall t
\]
\end{definizione}

Chiaramente ogni punto di equilibrio costituisce un orbita a sé.

\begin{proposizione}
Un punto $\overline{z}$ è un equilibrio di $\dot{z}=X(z)$ se e solo se $X(\overline{z})=0$.
\end{proposizione}
\begin{dimo}~\begin{itemize}
\item Se $\overline{z}$ è un equilibrio $z(t,\overline{z})=\overline{z}$ $\forall t$ ma allora  $\dfrac{dz(t;\overline{z})}{dt}(t)=0=X(z(t;\overline{z}))=X(\overline{z})$.
\item Se $X(\overline{z})=0=\dot{z}$ allora la curva costante $\varphi(t)=\overline{z}$ soddisfa $\dot{\varphi}(t)=X(\varphi(t))$ e quindi $z(t;\overline{z})=\varphi(t)=\overline{z}$ è una soluzione e $\overline{z}$ è un equilibrio. \qedhere
\end{itemize}
\end{dimo}

\begin{esempio}[Equilibri di alcuni sistemi elementari]
\begin{itemize}
\item Per $\dot{z}=Az$ con $A\in M_{n\times n}$ gli equilibri sono gli zeri di $X(z)=Az$ per cui è l'iniseme $\ker{A}$, se la matrice $A$ diagonalizzabile l'unico equilibrio è l'origine. 
\item Per \[
\left(\begin{array}{c}
\dot{x} \\ \dot{y} \\ \dot{z} \end{array}\right) = \left(\begin{array}{c}
y(z-1) \\ x (1-z) \\ x+z \end{array}\right)
\]
gli equilibri sono gli zeri del campo quindi: per $y\neq0$ il campo è nullo per $x=-1$ e $z=1$ per ogni valore di $y$ mentre per $y=0$ il campo è nullo in $x=z=0$ e in $z=1=-x$. Quindi l'origine è un punto di equilibrio è tutta la $\{x=-1,z=1\}$ è una retta di equilibrio.
\item Per $z\in\mathbb{R}^2$ e $\dot{z}=\left(\begin{array}{cc} 0 & 1 \\ 0 & -1 \end{array}\right)z$ che equivale a $\dot{z}=\left(\begin{array}{c}
z_2 \\ - z-2 \end{array}\right)$ gli equilibri sono tutti i punti con $z_2=0$. 
\end{itemize}
\end{esempio}

\begin{definizione}Dato un equilibrio $\overline{z}$ lo chiameremo:
\begin{itemize}
\item Attrattivo se $\exists V\subset\Omega$ intorno di $\overline{z}$ per cui $\forall z_0 \in V$ ogni soluzione $z_t$ con dato iniziale $z_0$ soddisfi \[\lim_{t\rightarrow +\infty} z_t=\overline{z} \]
\item Repulsivo se $\exists V\subset\Omega$ intorno di $\overline{z}$ per cui $\forall z_0 \in V$ ogni soluzione $z_t$ con dato iniziale $z_0$ soddisfi \[\lim_{t\rightarrow -\infty} z_t=\overline{z} \]
\end{itemize}
\end{definizione}

\begin{proposizione}
In dimensione $1$ per $\dot{z}=X(z)\in\mathbb{R}$ con un equilibrio $\overline{z}$ allora $\overline{z}$ è attrattivo se e solo $X'(\overline{z})<0$ mentre è repulsivo se e solo se $X'(\overline{z})>0$.
\end{proposizione}
\begin{dimo}
Guardando al ritratto in fase gli equilibri sono zeri del campo e sono attrattivi se $X(z)>0$ per $z<\overline{z}$ e $X(z)<0$ per $z>\overline{z}$ che, se $X'(\overline{z})\neq0$, significa che $X'(z)<0$. Equivalentemente per equilibri repulsivi se $X'(\overline{z})\neq0$ allora  $X'(z)>0$.\qedhere
\end{dimo}

\begin{esempio}
Consideriamo l'equazione differenziale scalare $\dot{z}=z(1-z)-\sfrac{1}{4}$. Gli equilibri saranno le soluzioni di $z^2 -z +\sfrac{1}{4} =0$ cioè il punto $z=\sfrac{1}{2}$. $X'(z)=1-2z$ e $X'(\sfrac{1}{2})=0$ quindi  non possiamo usare il criterio semplice, notiamo però che $X(z)$, e quindi anche la velocità di percorrenza delle orbite, è positivo per $z<\sfrac{1}{2}$ e negativo per $z>\sfrac{1}{2}$ quindi l'equilibrio sarà attrattivo.

Essendo in dimensione $1$ con l'equilibrio è la monotonia della velocità è facile costruire il ritratto in fase:\par
\centering
\includegraphics[width=0.7\linewidth]{img/ritrattoInFaseBase1.png}
\end{esempio}

Per un equazione del secondo ordine $\ddot{x}=Y(x,\dot{x})$ come abbiamo visto si ritroveranno i risultati delle equazioni del primo ordine considerando il campo vettoriale $X(x,v)=(v,Y(x,v))$ e quindi ovviamente perché questo campo possa annullarsi servirà che sia $v=0$.
\begin{proposizione}
Gli equilibri di un'equazione del secondo ordine sono tutti e soli i punti dello spazio delle fasi $(\overline{x},0)\in TC$ per cui $Y(\overline{x},0)=0$.
\end{proposizione} 
Tutto ciò è ovvio pensando a partire dalla definizione di equilibrio: se la soluzione deve essere ferma nel punto $\overline{x}$ chiaramente la sua derivata in quel punto sarà nulla, quindi $v=0$.


\subsection{Linearizzazione}\label{sezLinearizzazione}
Una volta determinati gli equilibri di un'equazione differenziale possiamo sfruttarli per studiarla in modo semplificato in un intorno dell'equilibrio:
sia $\overline{z}$ un equilibrio di $\dot{z}=X(z)$, questo significa che $X(\overline{z})=0$ e potremo sviluppare il campo in serie attorno a $\overline{z}$
\[
X(z)=X'(\overline{z})(z-\overline{z}) + o(z-\overline{z})
\] 
e se ci teniamo sufficientemente vicino a $\overline{z}$ i termini di ordine superiore al primo saranno trascurabili. 
Quindi in un intorno di un equilibrio $\overline{z}$ l'equazione si comporterà come un'equazione lineare
\begin{equation}\label{linearizzazione}
\dot{z}=A(z-\overline{z}) \hspace{0.15\linewidth} J_X(\overline{z})=A\in M_{n\times n}(\mathbb{R})
\end{equation}
Potremo anche definire con una traslazione su $z$ la variabile $u=z-\overline{z}$ per la quale la \eqref{linearizzazione} diventerà l'equazione lineare omogenea $\dot{u}=Au$.


Per un'equazione del secondo ordine $\ddot{x}=Y(x,\dot{x})$ possiamo linearizzare sia il sistema del primo ordine sia l'equazione del secondo ordine:
\begin{enumerate}[i.]
\item Considerando il sistema di primo grado sappiamo che un equilibrio sarà un $(\overline{x};0)\in\Omega$ e espandiamo in serie la seconda equazione del sistema ottenendo
\[
\begin{cases}
\dot{x}=v \\
\dot{v}=Y(x,v) \\
\end{cases} \implies
\begin{cases}
\dot{x}=v \\
\dot{v}=\dfrac{\partial Y}{\partial x}(\overline{x},0) (x-\overline{x}) + \dfrac{\partial Y}{\partial v}(\overline{x},0) v 
\end{cases}
\]
\item Equivalentemente sviluppando con Taylor l'equazione di secondo grado avremo
\[
\ddot{x}=\dfrac{\partial Y}{\partial x}(\overline{x},0) (x-\overline{x}) + \dfrac{\partial Y}{\partial v}(\overline{x},0) v 
\]
\end{enumerate}





\section{Equazioni differenziali lineari}
Dato un'equazione differenziale lineare $\dot{z}=Az$ dove $z\in\Omega\subseteq\mathbb{R}^n$ possiamo scrivere il flusso dell'equazione tramite la matrice esponenziale $\exp{A}$

\begin{proposizione}
La soluzione di $\dot{z}=Az$ passante per $z_0$ è
\[
z(t;z_0)=\exp{(tA)}z_0
\]
\end{proposizione}
\begin{dimo}
La derivata della funzione $t\mapsto\exp{(tA)}z_0$ si calcola derivando il termine generale della serie che definisce la matrice esponenziale
\[
\dfrac{d}{dt}\exp{(tA)}z_0=\sum_{k=0}^{+\infty}\dfrac{d}{dt}\left(\dfrac{t^k A^k}{k!}\right)z_0=\sum_{k=0}^{+\infty}\dfrac{t^{k-1}A^k}{(k-1)!}z_0=A\sum_{k=0}^{+\infty}\dfrac{t^{k-1}A^{k-1}}{(k-1)!}z_0=A\exp{(tA)}z_0
\]
Quindi $\varphi(t)=\exp{(tA)}z_0$ soddisfa $\dot{\varphi}(t)=A\varphi(t)$ e $ \exp{(tA)}z_0$ è la soluzione passante per $z_0$.\footnotemark{}  \qedhere
\end{dimo}
\footnotetext{Si veda il materiale di \textit{Analisi 3} per una trattazione più generale di questo argomento. Ai fini di questo corso è sufficiente sapere che nelle equazioni lineari autonome la matrice esponenziale definisce l'unica soluzione passante per un dato punto.} 


Ricordiamo che un sottospazio di $V\subseteq\mathbb{R}^n$ si dice invariante per una matrice $A\in M_n(\mathbb{R})$ se $\forall v\in V$ rimane $Av\in V$. Questo significa che un sottoinsieme dello spazio delle fasi è invariante per l'equazione differenziale se e solo se è invariante per la matrice $\exp{tA}$ ad ogni tempo $t\in \mathbb{R}$.

Si dimostra che un sottospazio dello spazio delle fasi è invariante per l'equazione differenziale $\dot{z}=Az$ se e solo se è invariante per la matrice $A$, useremo quindi ``invariante per l'equazione differenziale'', ``invariante per $\exp{tA}$'' e ``invariante per $A$'' come sinonimi.

\subsection{Matrici diagonalizzabili}\label{sezMatriciDiagonalizzabili}

Per una matrice diagonale $D=\text{diag}(D_{1,1},D_{2,2},...,D_{n,n})$ la matrice esponenziale è 
\[\exp{D}=\exp{\left(\text{diag}(D_{1,1},D_{2,2},...,D_{n,n})\right)}=\text{diag}(e^{D_{1,1}},e^{D_{2,2}},...,e^{D_{n,n}})\]
Da questo si trova che se la matrice $A$ è diagonalizzabile
$\exists P\in M_{n,n}$ tale che 
\[P^{-1}AP=\text{diag}\left( \lambda_1,...,\lambda_n \right)\] 
dove i $\lambda_j$ sono gli autovalori di $A$, in generale complessi.
Si dimostra che
\[
\exp{tA}=P \exp{\left( tP^{-1}AP \right)} P^{-1}
\]


Sappiamo che per la matrice esponenziale (e in realtà per una matrice qualsiasi) vale
\[
\exp{(tA)}(v+w)=\exp{(tA)}v + \exp{(tA)}w \>\>\> \forall t\in\mathbb{R}
\]

Se la matrice $A\in M_n(\mathbb{R})$ è diagonalizzabile su $\mathbb{C}$ vi saranno, in generale, $n$ autovettori complessi linearmente indipendenti su $\mathbb{C}$.

Come sappiamo se un $v\in \mathbb{C}$ è un autovettore complesso con parte immaginaria non nulla di una matrice $A$ reale  anche $\overline{v}\in\mathbb{C}$ sarà un autovettore, relativo al autovalore complesso coniugato di quello del primo.
\begin{dimo}
Sia $Av=\lambda v$ allora coniugando $\overline{Av}=\overline{\lambda v}$ e però dato che $A$ è reale vale $\overline{Av}=A\overline{v}=\overline{\lambda}\overline{v}$.
\end{dimo}

Allora possiamo dire in tutta generalità che $A$ avrà $k$ coppie di autovalori complessi coniugati e $n-2k$ autovalori reali
\[
\alpha_1 \pm i \beta_1 ,..., \alpha_k \pm i \beta_k \hspace{0.15\linewidth} \alpha_{2k+1},...,\alpha_n
\]


Se due vettori complessi coniugati sono linearmente indipendenti su $\mathbb{C}$ allora le loro parti reali e immaginarie sono linearmente indipendenti su $\mathbb{R}$
\begin{dimo}
Siano $z=v+iw$ e $\overline{z}=v-iw$ linearmente indipendenti su $\mathbb{C}$. Siano per assurdo $v,w$ linearmente dipendenti su $\mathbb{R}$, cioè $w=\lambda v$ per un $\lambda\in\mathbb{R}$. Allora $w=\sfrac{z-\overline{z}}{2i} \, =\lambda \sfrac{z+\overline{z}}{2}\,=\lambda v$ ma quindi $z(i\lambda-1)=\overline{z}(1-i\lambda)$ che è assurdo.
\end{dimo}
Potremo scegliere autovettori reali per gli autovalori reali e quindi avremo una base di $\mathbb{R}^n$
\begin{equation}\label{base decomposizione invarianti}
\left\{
v_1,w_1,v_2,w_2,...,v_k,w_k,v_{2k+1},...,v_n
\right\}
\end{equation}
e una decomposizione
\begin{equation}\label{decomposizione invarianti}
\mathbb{R}^n= \left\langle v_1,w_1 \right\rangle \oplus \left\langle v_2,w_2 \right\rangle \oplus ... \left\langle v_k,w_k \right\rangle \oplus \left\langle v_{2k+1} \right\rangle \oplus ... \oplus \left\langle v_{n} \right\rangle 
\end{equation}

Ciascuno di questi spazi sarà per costruzione invariante per la matrice $A$: quelli unidimensionali sono semplicemente autospazi mentre per quelli bidimensionali
\[
A(v+iw)=(\alpha+i\beta)(v+iw)\implies \begin{cases}Av=\alpha v-\beta w \\ Aw=\beta v+\alpha w \end{cases}
\]


Se $B$ è la matrice di una trasformazione lineare $\mathcal{B}:\mathbb{R}^n\rightarrow\mathbb{R}^n$ e $U,W\subset\mathbb{R}^n$ sono due sottospazi invarianti per $B$ che decompongono $\mathbb{R}^n=U\oplus W$ allora la matrice $B$ si potrà scrivere nella base data dall'unione delle basi dei due sottospazi in forma diagonale a blocchi
\begin{equation}\label{matrice blocchi sottospazi invarianti}
\left(
\begin{array}{cc}
B_U & 0 \\
0 & B_W \\
\end{array}
\right)
\end{equation}
e quindi per la trasformazione $\mathcal{A}$, di matrice $A$ nella base canonica, nella base della decomposizione in sottospazi invarianti per $A$ di $\mathbb{R}^n$ la matrice $\hat{A}$ di $\mathcal{A}$ sarà una matrice a blocchi e la sua esponenziale sarà
\[
\exp{t\hat{A}}=\text{diag} \left( \exp{tA}_{\langle w_1,w_1 \rangle},...,\exp{tA}_{\langle v_k,w_k \rangle },e^{t\alpha_{2k+1}},...,e^{t\alpha_{n}}  \right) 
\]
Dove le $\exp{tA}_{\langle w_i,w_i \rangle}$ sono matrici $2\times 2$ e i $\exp{tA_{\langle v_{j} \rangle}}=\exp{t\alpha_{j}}$ sono entrate $1\times 1$ con $\alpha_j$ l'autovalore relativo all'autospazio $\langle v_j \rangle$.

\begin{proposizione}\label{proposizioneRestrizioneMatriceSottospaziInvarianti}Se $A$ è diagonalizzabile su $\mathbb{C}$, con $k$ coppie di autovalori $\alpha_j\pm i\beta_j$ complessi coniugati con autovettori $v_j \pm iw_j$ e $n-2k$ autovalori reali $\alpha_l$ con autovettori $v_l$, per ogni $t\in\mathbb{R}$:
\begin{itemize}
\item La restrizione di $\exp{tA}$ a ciascun sottospazio unidimensionale $\langle v_l \rangle$ è una dilatazione di fattore $e^{t\alpha_l}$
\item La restrizione di $\exp{tA}$ a ciascun sottospazio bidimensionale $\langle v_j,w_j \rangle$ è una composizione tra una dilatazione di fattore $e^{t\alpha_j}$ e  una rotazione di angolo $\beta_j t$
\end{itemize}
\end{proposizione}
\begin{dimo}$ $
\begin{itemize} 
\item La restrizione di $A$ agli spazi unidimensionali ovviamente dà $Au=\alpha_j u$ quindi direttamente dalla scrittura \eqref{matrice blocchi sottospazi invarianti} la matrice $\exp{tA}$ agendo su un vettore di $\langle v_j \rangle$ dà, nella base \eqref{base decomposizione invarianti}, $\exp{(tA)} u= e^{t\alpha_j}u$
\item Consideriamo un sottospazio bidimensionale invariante $\langle v,w \rangle$, chiamiamo $\hat{\mathcal{A}}=\mathcal{A}_{|\langle v,w \rangle}$ con $\mathcal{A}(u)=Au$.
Da $Av=\alpha v -\beta w$ e $Aw=\beta v + \alpha w$ segue che la matrice relativa ad $\hat{\mathcal{A}}$ relativa alla base $\{v,w\}$ è 
\[
\hat{A}=\left( \begin{array}{cc}
\alpha & \beta \\
-\beta & \alpha \\
\end{array}\right) = \left( \begin{array}{cc}
\alpha & 0 \\
0 & \alpha \\
\end{array}\right)+\left( \begin{array}{cc}
0 & \beta \\
-\beta & 0 \\
\end{array}\right)
\]
La matrice della restrizione di $z\mapsto \exp{(tA})z$ a $\langle v, w \rangle$ è $\exp{t\hat{A}}$  
\[
\exp{t\hat{A}}=\exp{  \left( \begin{array}{cc}
\alpha t & 0 \\
0 & \alpha t \\
\end{array}\right) } + \exp{ \left( \begin{array}{cc}
0 & \beta t \\
-\beta t & 0 \\
\end{array}\right)  } =e^{t\alpha}\mathbb{I}+\exp{ \left( \begin{array}{cc}
0 & \beta t \\
-\beta t & 0 \\
\end{array}\right)  }= e^{t\alpha}R(\beta t) \qedhere
\] 
\end{itemize}
\end{dimo}

Quindi la dinamica all'interno di ciascuno dei sottospazi unidimensionali $\langle v_l \rangle$ vista dalla base $\{v_l\}$ è un'espansione esponenziale se $\alpha_l>0$ o una contrazione esponenziale se $\alpha_l<0$.
Nei sottospazi bidimensionali invece, per le basi $\{v_j,w_j\}$ le orbite sono spirali logaritmiche percorse verso l'infinito se $\alpha_j>0$ e verso l'origine se $\alpha_j<0$ o cerchi se $\alpha_j=0$. Ritornando alla base canonica i cerchi diventeranno ellissi e le spirali si deformeranno di conseguenza.



\subsection{Sistemi lineari in $\mathbb{R}^2$}
Sia ora $\dot{z}=Az$ con $z\in\mathbb{R}^2$ e $A\in M_2(\mathbb{R})$ diagonalizzabile,  avremo sostanzialmente $4$ possibilità: 
\begin{itemize}
\item due autovalori complessi coniugati
\item due autovalori reali distinti non nulli
\item due autovalori reali distinti di cui uno nullo 
\item un autovalore reale doppio.
\end{itemize}

Chiaramente nei primi due casi avremo un unico equilibrio che coinciderà con l'origine e potremo facilmente studiare il ritratto in fase passando alla base di autovettori e usando quanto visto nella sezione precedente per poi infine ritornare alla base canonica.

\begin{definizione}
Data in $\mathbb{R}^2$ l'equazione $\dot{z}=Az$, $z\in \mathbb{R}^2$ se $A\in M_2(\mathbb{R})$ è diagonalizzabile in $\mathbb{C}$ con autovalori complessi coniugati $\alpha\pm i \beta$ con $\beta \neq 0$ e autovettori $v\pm iw$ i ritratti in fase si dicono centri se $\alpha=0$, fuochi stabili se $\alpha_j<0$ e fuochi instabili se $\alpha_j>0$.
\end{definizione}

Nella base canonica di $\mathbb{R}^2$ questi $3$ ritratti in fase sono di forma:
\begin{table}[H]
\begin{minipage}{0.2\linewidth}
$ $
\end{minipage}
\begin{minipage}[H]{0.20\linewidth}
\centering
	\begin{figure}[H]
  	\includegraphics[width=\linewidth]{img/fuocostabile.png}
	\end{figure}

\end{minipage}
\begin{minipage}{0.18\linewidth}
\centering
	\begin{figure}[H]
  	\includegraphics[width=\linewidth]{img/centro.png}
	\end{figure}
\end{minipage}
\begin{minipage}[H]{0.20\linewidth}
\centering
	\begin{figure}[H]
  	\includegraphics[width=\linewidth]{img/fuocoinstabile.png}
	\end{figure}
\end{minipage}
\begin{minipage}{0.15\linewidth}
$ $
\end{minipage}
\\
\begin{minipage}{0.2\linewidth}
$ $
\end{minipage}
\begin{minipage}[H]{0.20\linewidth}
\centering
Fuoco stabile
\end{minipage}
\begin{minipage}{0.18\linewidth}
\centering
Centro
\end{minipage}
\begin{minipage}[H]{0.20\linewidth}
\centering
Fuoco instabile
\end{minipage}
\begin{minipage}{0.15\linewidth}
$ $
\end{minipage}

\hspace{0.01\textheight}

\centering\small Ritratti in fase nella base $\{u_1,u_2\}$

\end{table}


\subsubsection{Autovalori reali distinti}
Trattiamo il caso degli autovalori reali distinti un po' più esplicitamente: siano $u_1$ e $u_2$ gli autovettori relativi agli autovalori $\alpha_1$ e $\alpha_2$, la matrice di $A$ sarà diagonalizzabile ed invertibile quindi se definiamo
\[
z_t \text{ soluzione di } \dot{z}=Az \hspace{0.14\linewidth} P=(u_1,u_2) \hspace{0.15\linewidth} y=P^{-1}z
\]
vediamo che per le proprietà della matrice esponenziale
\[
P^{-1}\exp{(tA)}P=PP^{-1}\exp{(P^{-1}tAP)}PP^{-1}=\exp{tP^{-1}AP}=\text{diag}(\alpha_1,\alpha_2)
\]
e quindi per le soluzioni varrà
\[
y_t=P^{-1}z_t=P^{-1}\exp{(tA)}z_0=P^{-1}\exp{(tA)}Py_0=\left( \begin{array}{cc}
\alpha_1 & 0\\ \alpha_2 &0\\
\end{array}\right) y_0
\]
E quindi otteniamo soluzioni
\[
y(t)=\left(\begin{array}{c}
y_1(t)\\
y_2(t)\\
\end{array}\right)=\left(\begin{array}{c}
e^{t\alpha_1}y_{1,0}\\
e^{t\alpha_2}y_{2,0}\\
\end{array}\right)
\]
Questo ci fa subito capire che gli assi dati dai due autovettori $u_1, u_2$ daranno (4) orbite per dati iniziali con una coordinata nulla e quindi che le altre orbite dovranno necessariamente essere confinate all'interno dei $4$ quadranti nella base $\{u_1,u_2\}$.

Per dato iniziale con entrambe le coordinate non nulle possiamo eliminare il tempo e trovare esplicitamente le equazioni delle orbite 
\[
y_2(t)=k\, |y_1|^{\dfrac{\alpha_2}{\alpha_1}}
\]
Quindi nelle coordinate date dalla base $\{u_1,u_2\}$ avremo orbite racchiuse nei quadranti che si dovranno orientare in base al segno degli autovalori.
I ritratti in fase saranno quindi:
\begin{itemize}
\item se entrambi gli autovalori sono negativi le orbite sono curve simili a parabole tangenti nell'origine all'autovalore di valore assoluto minore e dirette da infinito a 0: chiameremo questo ritratto in fase un \textbf{nodo stabile}
\item Se $\alpha_1$ è positivo e $\alpha_2$ è negativo le orbite sono curve simili a iperboli con proiezione lungo $u_1$ andrà ad infinito per $t\rightarrow \infty$ e con proiezione lungo $u_2$ che andrà invece a 0 per $t\rightarrow \infty$: chiameremo questo ritratto in fase una \textbf{sella}
\item se entrambi gli autovalori sono positivi le orbite sono curve simili a parabole tangenti nell'origine all'autovalore di valore assoluto minore e dirette da 0 ad infinito: chiameremo questo ritratto in fase un \textbf{nodo instabile}

\end{itemize}

Nella base canonica di $\mathbb{R}^2$ questi $3$ ritratti in fase sono di forma:
\begin{table}[H]
\begin{minipage}{0.2\linewidth}
$ $
\end{minipage}
\begin{minipage}[H]{0.20\linewidth}
\centering
	\begin{figure}[H]
  	\includegraphics[width=\linewidth]{img/nodostabile.png}
	\end{figure}
Nodo stabile
\end{minipage}
\begin{minipage}{0.20\linewidth}
\centering
	\begin{figure}[H]
  	\includegraphics[width=\linewidth]{img/sella.png}
	\end{figure}
Sella
\end{minipage}
\begin{minipage}[H]{0.20\linewidth}
\centering
	\begin{figure}[H]
  	\includegraphics[width=\linewidth]{img/nodoinstabile.png}
	\end{figure}
Nodo instabile
\end{minipage}
\begin{minipage}{0.15\linewidth}
$ $
\end{minipage}

\hspace{0.014\textheight}

\centering\small Ritratti in fase nella base $\{u_1,u_2\}$
\end{table}

\begin{esempio}[Repulsore armonico]
L'equazione $\ddot{x}=kx$ con $k>0$ è l'equazione del repulsore armonico. In modo più esplicito la possiamo scrivere come $\begin{cases} \dot{x}=v \\ \dot{v}=\omega^2 x \end{cases}$.

Il sistema del primo ordine è lineare con matrice $A=\left(\begin{array}{cc} 0 & 1 \\ \omega^2 & 0\end{array}\right)$ di autovalori $\pm \omega$ e con autovettori $u_\pm =(1,\pm \omega)$.

Da questo sappiamo che l'origine è una sella.
\end{esempio}

\subsubsection{Autovalore doppio}
Se gli autovalori sono $\alpha_1=\alpha_2=\alpha$ la matrice digonalizzata sarà un multiplo dell'identità. Tutti i vettori possibili saranno autovettori e il ritratto sarà quello di un \textbf{nodo a stella stabile} se $\alpha<0$ o di un \textbf{nodo a stella stabile} se $\alpha>0$.

Se invece $\alpha=0$ la matrice $A$ sarà la matrice nulla e tutti i punti saranno equilibri

 
\begin{table}[H]
\begin{minipage}{0.30\linewidth}
$ $
\end{minipage}
\begin{minipage}[H]{0.40\linewidth}
\centering
	\begin{figure}[H]
  	\includegraphics[width=\linewidth]{img/nodiAstella.png}
	\end{figure}
	\begin{minipage}[H]{0.45\linewidth}
		\centering
		Nodo a stella stabile
	\end{minipage}
		\begin{minipage}[H]{0.45\linewidth}
		\centering
		Nodo a stella instabile
	\end{minipage}
\end{minipage}
\begin{minipage}{0.25\linewidth}
$ $
\end{minipage}

\hspace{0.014\textheight}

\centering\small Ritratti in fase nella base $\{u_1,u_2\}$
\end{table}

\subsubsection{Autovalore nullo}

Se uno dei due autovalori è nulla $\alpha_1=0$ le orbite saranno tutti i punti di $\langle u_1 \rangle =\ker{A}$ e tutte le semirette che partono da essi parallele a $\langle u_2\rangle$ con verso di percorrenza determinato dal segno di $\alpha_2$. Chiameremo questi ritratti in fase \textbf{nodi a pettine}.

\begin{table}[H]
\begin{minipage}{0.40\linewidth}
$ $
\end{minipage}
\begin{minipage}[H]{0.20\linewidth}
\centering
	\begin{figure}[H]
  	\includegraphics[width=\linewidth]{img/nodoPettine.png}
	\end{figure}
	Nodo a pettine con $\alpha_2<0$
\end{minipage}
\begin{minipage}{0.35\linewidth}
$ $
\end{minipage}
	
\hspace{0.014\textheight}

\centering\small Ritratto in fase nella base $\{u_1,u_2\}$.

\end{table}

\subsubsection{Matrice nilpotente}

Se la matrice $A\in M_{2\times2}(\mathbb{R})$ non è diagonalizzabile nè in $\mathbb{R}$ nè in $\mathbb{C}$ che ritratti in fase possiamo ottenere?

Questa situazione si può incontrare soltanto quando la matrice ha un autovalore doppio reale $\lambda$ con però un autospazio relativo di dimensione $1$.

Sappiamo che ogni matrice è simile ad una matrice triangolare quindi
\[
A\sim b=\left(\begin{array}{cc}
\lambda & b \\ 0 & \lambda
\end{array}\right)
\]
dove $b\in\mathbb{R}$, $b\neq0$. L'autospazio unidimensionale relativo a $\lambda$ è $\langle (1,0)\rangle$.

Il sistema $\dot{z}=Bz$ è
\[
\dot{x}=\lambda x+by \qquad \quad \dot{y}=\lambda y
\]

e possiamo integrarlo con tecniche elementari ottenendo
\[
x(t)=e^{\lambda t}(x_0+by_0 t) \quad \qquad y(t)=e^{\lambda t}y_0
\]

I ritratti in fase quindi dipenderanno dal segno di $\lambda$ e $b$ e se $\lambda\neq 0$ si chiameranno \textbf{nodi degeneri} mentre per $\lambda=0$ il sistema sarà $\dot{x}=by$ $\dot{y}=0$.

\begin{table}[H]
\begin{minipage}{0.30\linewidth}
$ $
\end{minipage}
\begin{minipage}[H]{0.40\linewidth}
\centering
	\begin{figure}[H]
  	\includegraphics[width=\linewidth]{img/nodiDegeneri.png}
	\end{figure}
	\begin{minipage}[H]{0.45\linewidth}
		\centering
		$b>0, \lambda<0$ \\ Nodo degenere stabile
	\end{minipage}
		\begin{minipage}[H]{0.45\linewidth}
		\centering
		$b>0, \lambda>0$ \\ Nodo degenere instabile
	\end{minipage}
\end{minipage}
\begin{minipage}{0.25\linewidth}
$ $
\end{minipage}
\end{table}

\subsection{Sottospazi stabile, instabile e centrale}
Come abbiamo visto nella sezione \ref{sezMatriciDiagonalizzabili} data una $\dot{z}=Az$ in $\mathbb{R}^n$ con matrice $A$ diagonalizzabile potremo trattare semplicemente l'equazione differenziale usando la base $\{v_1,w_1,...,w_k,v_{2k+1},...,v_n$ di suoi autovettori e decomporre $\mathbb{R}^n$ in sottospazi invarianti
\[
\langle v_1,w_1 \rangle \oplus ... \oplus \langle v_k,w_k \rangle \oplus \langle v_{2k+1} \rangle \oplus ... \oplus \langle v_n \rangle
\]
su cui le proiezione delle orbite sono fuochi, centri o dilatazioni esponenziali.

Possiamo passare dalla base canonica alla base di autovettori e viceversa tramite la matrice $P=(v_1,w_1,...,w_k,...,v_n)$ e quindi, dato che nella base degli autovettori la matrice associata alla trasformazione $z=Az$ è $\hat{A}$, con $\hat{A}=P^{-1}AP$ dove $A$ è la matrice corrispondente alla stessa trasformazione ma nella base canonica.
Allora possiamo scrivere esplicitamente $\hat{A}$:
\[ 
\hat{A}=
\left(
\begin{array}{ccccccc}
B_1 &  &  &  &  &  &   \\
 & B_2 &  &  &  &  &   \\
 &  & \ddots &  &  &  &   \\
 &  &  & B_k &  &  &   \\
 &  &  &  &\alpha_{2k+1} &  &   \\
 &  &  &  &  & \ddots &  \\
 &  &  &  &  &  & \alpha_n   \\
\end{array}
\right)
\hspace{0.15\linewidth}
B_j=\left( \begin{array}{cc}
\alpha_j & \beta_j \\
-\beta_j & \alpha_j\\
\end{array}\right)
\]
e quindi la matrice esponenziale nella base degli autovettori è
\[
\exp{t\hat{A}}=
\left(
\begin{array}{ccccccc}
e^{\alpha_1 t}R(\beta_1 t) &  &  &  &  &  &   \\
 & e^{\alpha_2 t}R(\beta_2 t)  &  &  &  &  &   \\
 &  & \ddots &  &  &  &   \\
 &  &  & e^{\alpha_k t}R(\beta_k t)  &  &  &   \\
 &  &  &  & e^{\alpha_{2k+1}t} &  &   \\
 &  &  &  &  & \ddots &  \\
 &  &  &  &  &  & e^{\alpha_n t}   \\
\end{array}
\right)
\]
dove $R(\beta_jt)$ è la matrice $2\times 2$ di rotazione di angolo $\beta_j t$
\[
R(\beta_j t)=\left( \begin{array}{cc}
\cos{\beta_j t} & \sin{\beta_j t} \\
-\sin{\beta_j t} & \cos{\beta_j t}\\
\end{array}\right)
\]
E con questo abbiamo ripercorso quanto visto con la proposizione \ref{proposizioneRestrizioneMatriceSottospaziInvarianti}.
	
Consideriamo ora una soluzione nelle coordinate della base di autovettori
\[
(y_1,\tilde{y_1},...,y_k,\tilde{y}_k,y_{2k+1},...,y_n)(t) \hspace{0.15\linewidth} t\mapsto z_t \text{ soluzione di } \dot{z}=Az \text{ e } y(t)=P^{-1}z(t)
\]
questa avrà componenti
\[
\left(\begin{array}{c}
y_j(t) \\
\tilde{y_j}(t) \\
\end{array}\right)=
e^{\alpha_j t} R(\beta_j t)
\left(\begin{array}{c}
y_{j,0} \\
\tilde{y_{j,0}} \\
\end{array}\right)
\text{ con } j\in[1,k] \hspace{0.15\linewidth}
y_l(t)=e^{\alpha_l t}y_{l,0} \text{ con } l\in[2k+1,n]
\]

Questo mostra bene come le le componenti delle soluzioni sui singoli sottospazi invarianti siano semplici. Potremo da qui in poi cercare di ricostruire il ritratto in fase complessivo.

Possiamo infine definire a partire dalla decomposizione
\begin{definizione}\textit{Sottospazio stabile} \\
Chiamiamo sottospazio stabile la somma degli autospazi unidimensionali relativi a autovalori reali negativi e degli autospazi bidimensionali relativi ad autovalori complessi con parte reale negativa. Le proiezioni delle soluzioni in questo sottospazio tendono all'origine per $t\rightarrow +\infty$ e all'infinito per $t\rightarrow -\infty$.\\
\end{definizione}

\begin{definizione}\textit{Sottospazio centrale}\\
Chiamiamo sottospazio stabile la somma dell' autospazio relativo all' autovalore nullo, se presente, e dagli autospazi bidimensionali relativi ad autovalori immaginari. Le proiezioni delle soluzioni in questo sottospazio sono limitate per tutti i tempi e nessuna di esse tende all'origine per $t\rightarrow \pm\infty$.\\
\end{definizione}

\begin{definizione}[Sottospazio instabile]
Chiamiamo sottospazio stabile la somma degli autospazi unidimensionali relativi a autovalori reali positivi e degli autospazi bidimensionali relativi ad autovalori complessi con parte reale positiva. Le proiezioni delle soluzioni in questo sottospazio tendono all'origine per $t\rightarrow -\infty$ e all'infinito per $t\rightarrow +\infty$.\\
\end{definizione}

A seconda della dimensione e del sistema uno o più di questi sottospazi possono non essere presenti, per esempio in $\mathbb{R}^2$ se il ritratto in fase è una sella non è presente il sottospazio centrale mentre se il ritratto è un nodo o un fuoco stabile è presente solo il sottospazio stabile mentre se è un centro è presente solo il sottospazio centrale.

\begin{esempio}
La matrice \[A=\left(\begin{array}{ccc} -1 & 0 &0 \\ 0 & 0 & -1 \\ 0 & 1 & 0 \end{array}\right)\] è diagonalizzabile e ha autovalori $-1$ e $\pm i$con autovettori relativi rispettivamente $u=(1,0,0)$ e $(0,0,1)\pm i (0,1,0)$.
Il sottospazio unidimensionale $\langle(1,0,0)\rangle$, cioè le rette parallele all'asse $x$, sarà il sottospazio stabile mentre il sottospazio $\langle \{(0,1,0),(0,0,1)\}\rangle$, cioè i piani paralleli al piano $(y,z)$, è il sottospazio centale.

La base degli autovettori coincide con la base canonica e quindi il flusso dell'equazione $\dot{z}=Az$ è
\[
z(t)=\left(\begin{array}{ccc}
e^{-t} & 0 & 0 \\
0 & \cos{t} & -\sin{t} \\
0 & -\sin{t} & \cos{t} 
\end{array}\right) z_0
\]
da cui vediamo che, come ci aspettavamo dalla decomposizione in autovettori, le orbite sono spirali con proiezione circolare nel piano $(y,z)$.
\end{esempio}



\subsection{Dalla linearizzazione al sistema non lineare}
Come abbiamo visto nella sezione \ref{sezLinearizzazione} data una qualsiasi equazione lineare $\dot{z}=X(z)$ con equilibrio $\overline{z}$ potremo linearizzarla attorno all'equilibrio $\dot{z}=A(z-\overline{z})$ e trasformarla poi in un sistema di equazioni lineari omogenee $\dot{u}=Au$ con una traslazione $z\rightarrow u=z-\overline{z}$.

Abbiamo visto come sono i ritratti in fase di sistemi puramente lineari, ora vogliamo capire se almeno in piccoli intorni dell'equilibrio il ritratto in fase attorno all'origine di $\dot{u}=Au$ assomigli a quello di $z=X(z)$.

Il \textit{Teorema di Grobman-Hartman} ci dice che se tutti gli autovalori della Jacobiana calcolata nell'equilibrio hanno parte reale non nulla i ritratti in fase sono localmente simili. Se gli autovalori hanno parte reale nulla non si sa a priori se si assomiglino.\\È importante sottolineare che questo è un risultato locale e che il ritratto in fase complessivo potrebbe essere completamente diverso.



\begin{definizione}
Un equilibrio di un'equazione differenziale si dice iperbolico se tutti gli autovalori di $\dfrac{\partial X}{\partial z}(\overline{z})$ hanno parte reale non nulla e si dice ellittico se tutti gli autovalori sono puramente immaginaria.
\end{definizione}

Inoltre il \textit{Teorema della varietà stabile} dice che per un equilibrio iperbolico esistono globalmente due sottovarietà invarianti per l'equazione differenziale dette \textbf{varietà stabile} e \textbf{varietà instabile} che hanno la stessa dimensione e proprietà simili ai sottospazi stabile e instabile della linearizzazione attorno all'equilibrio. 
Tutti i moti della varietà stabile tendono all'equilibrio (iperbolico) per $t\rightarrow +\infty$ e la varietà è tangente nell'equilibrio al sottospazio, allo stesso modo i moti della varietà instabile i moti tendono all'equilibrio per $t\rightarrow -\infty$ e questa è tangente al sottospazio nell'equilibrio.

Rami di una varietà stabile e di una varietà instabile dello stesso equilibrio o di equilibri diversi possono coincidere.


\section{Dipendenza dai dati iniziali}

\begin{proposizione}\label{propMappaFlussoDifferenziabile}
Se $X$ è un campo differenziabile allora, per ogni $t\in\mathbb{R}$ la  mappa \[z_0\mapsto z(t,z_0)\] è differenziabile
\end{proposizione}

La differenziabilità implica la continuità quindi ci aspetteremmo che ``piccoli'' cambiamenti sui dati iniziali diano soluzioni ``vicine'' però, come già sappiamo intuitivamente, soluzioni che partono vicine possono andare in direzioni completamente diverse quando il tempo cresce verso infinito poiché questa continuità non è uniforme per $t$.

La continuità è
\[
\forall \epsilon>0 \> \exists \delta_{t,\epsilon}>0 : \left|\left| z_0-z'_0\right|\right| <\delta_{t,\epsilon} \implies \left|\left| z_t(z_0)-z_t(z'_0) \right|\right|<\epsilon   
\]

che possiamo interpretare come
\[
\forall\epsilon,T>0 \exists \Delta_{\epsilon,T}>0 : \left|\left| z_0-z'_0\right|\right| <\Delta_{t,\epsilon} \implies \left|\left| z_t(z_0)-z_t(z'_0) \right|\right|<\epsilon \> \forall t\in[-T,T]
\]
prendendo 
\[\Delta_{\epsilon,T}=\underset{t\in[-T,T]}{\text{min}}{\delta_{\epsilon,t}}\]
e come ci aspettavamo questo perde di senso se $T\rightarrow\infty$ e $\Delta_{\epsilon,t}\rightarrow 0$
Vogliamo però una stima più precisa di quanto due soluzioni diverse si distanziano nel tempo

\begin{proposizione}
Sia $X$ un campo vettoriale definito e differenziabile in $\mathbb{R}^n$ lipschitziano con costante di Lipschitz $K$, allora per ogni coppia di punti $z_0,z'_0\in\mathbb{R}^n$ si ha
\[
\left|\left| z_t(z_0)-z_t(z'_0) \right|\right| \leq e^{K\left|t\right|} \left|\left| z_0-z'_0 \right|\right| \>\>\> \forall t\in\mathbb{R}
\]
\end{proposizione}
\begin{dimo}
Ogni campo vettoriale lipschitziano è completo e quindi ne esistono soluzioni per tutti i tempi.

Siano $z_t=(z,t,z_0)$ e $z'_t=z(t,z'_0)$ con $z_0\neq z'_0$ soluzioni con dato iniziale diverso
\[
\dot{z_t}=X(z_t) \hspace{0.15\linewidth} \dot{z}'_t=X(z'_t)
\]
Definiamo $\delta_t\mathrel{\mathop:}=\left|\left| z_t-z'_t\right|\right|$ e sappiamo che sarà positiva e non nulla ad ogni $t$ dall'unicità delle soluzioni. Useremo la più comoda $\delta_t^2 = (z_t-z'_t)\cdot(z_t-z'_t)$ e derivandola
\[
\dfrac{\partial}{\partial t} \delta_t^2 = 2(\dot{z_t}-\dot{z}'_t)\cdot(z_t-z'_t)=2\left(X(z_t)-X(z'_t)\right)\cdot(z_t-z'_t) 
\]
Quindi
\begin{align*}
\left| \dfrac{\partial}{\partial t} \delta_t^2  \right| &= 2\left|\left| X(z_t)-X(z'_t) \right|\right| \left|\left|  z_t-z'_t \right|\right| \leq 2 K \left|\left| z_t-z'_t \right|\right| \left|\left| z_t-z'_t \right|\right| \\&= 2K \left|\left| z_t-z'_t \right|\right|^2=2K\delta_t^2
\end{align*}
Ma $\left| \dfrac{\partial}{\partial t} \delta_t^2  \right|=2\delta_t\left|\dfrac{\partial}{\partial t}\delta_t\right|$ e, dato che $\delta_t$ non può mai annullarsi, troviamo
$\left|\dfrac{\partial}{\partial t} \delta_t \right|\leq K\delta_t$.\\
Quindi la velocità con cui varia la distanza fra le due soluzioni è compresa tra $-K\delta_t$ e $K\delta_t$.\footnotemark

Allora chiamiamo $\dfrac{\partial}{\partial t}g_\pm = K g_\pm $ e il grafico della $\delta_t$ sarà sempre compreso tra i loro due grafici, le due funzioni si possono integrare facilmente e valgono $g_\pm(t)=\delta_0 e^{\pm Kt}$. La $\delta_t$ non può uscire dalla regione compresa tra i due grafici perché se lo facesse nell'istante $\overline{t}$ di uscita la sua derivata avrebbe valore assoluto maggiore di $K\delta_{\overline{t}}$

Dunque 
\[\left|\left|  z(t,z_0)-z(t,z'_0) \right|\right|\leq e^{K \left| t \right|} \left|\left|  z_0-z'_0 \right|\right|
\qedhere\]
\end{dimo}


\vfill
\begin{footnotesize}

Gli ultimi passaggi della dimostrazione riportata personalmente non mi sembrano troppo convincenti quindi provo a darne una versione alternativa.

Per $t\geq0$ la velocità massima di allontanamento è $K\delta_t$ quindi la distanza massima sarà quella ottenuta allontanandosi a velocità massima ad ogni istante che rispetterà
\[
\dfrac{\partial}{\partial t}\delta_{max}(t)= K \delta_{max}(t) \implies \delta_{max}(t)=e^{KT}\delta_0
\]
Ripetendo per $t\leq0$ si trova $\delta_{max}(t)=e^{-Kt}\delta_0=e^{K \left| t \right|}\delta_0$ e quindi la distanza a tutti i tempi sarà minore di
\[
\delta_t = \left|\left|  z(t,z_0)-z(t,z'_0) \right|\right|\leq e^{K\left|t\right|}\delta_0
\qedhere 
\]
\end{footnotesize}

\begin{esempio}
Per l'equazione differenziale $\dot{z}=kz$ con $z\in\mathbb{R}$ e $k\neq 0$ il campo $X(z)=kz$ è lipschitziano con costante $k$ quindi ci aspettiamo che la separazione delle soluzioni sia minore o uguale a $e^{kt} |z_0-z_0'|$.

L'integrale dell'equazione è $z(t;z_0)=e^{kt}z_0$ quindi la distanza fra due soluzioni con dati iniziali diversi è $|z(t;z_0)-z(t;z_0')|=e^{kt} |z_0-z_0'|$.
\end{esempio}

\section{Flusso e proprietà}

La definizione \ref{defFlusso1} introduce la mappa del flusso di un campo vettoriale, la ridefiniamo qui in modo un po' più preciso:

\begin{definizione}\textit{Flusso di un campo vettoriale}\\
Dato un campo vettoriale $X$ che sia definito e differenziabile su $\Omega\subseteq\mathbb{R}^n$, assumendo la completezza di $X$ possiamo definire
\[
\Phi^X:\mathbb{R}\times\Omega\rightarrow\Omega \hspace{0.15\linewidth} (t,z_0)\mapsto\Phi^X(t,z_0)\mathrel{\mathop:}=z(t,z_0)
\]
flusso del campo vettoriale $X$ per l'equazione differenziale autonoma $\dot{z}=X(z)$.
Fissando un qualsiasi $t\in\mathbb{R}$ definiamo mappa del flusso al tempo $t$ 
\[
\Phi^X_t:\Omega\rightarrow\Omega \hspace{0.15\linewidth} z_0\mapsto\Phi^X_t(z_0)\mathrel{\mathop:}=\Phi^X(t,z_0)=z(t,z_0)
\]
\end{definizione}

\begin{esempio}
L'oscillatore armonico $\dot{x}=v $ $\dot{v}=-\omega^2 z$ ha flusso
\[
(t,z_0)\mapsto\Phi(t,z_0)=\left(\begin{array}{cc}
\cos{\omega t} & \dfrac{1}{\omega} \sin{\omega t} \\
-\omega \sin{\omega t} & \cos{\omega t}
\end{array}\right) \left(\begin{array}{c}
x_0 \\ v_0
\end{array}\right)
\]
Quindi il flusso dell'oscillatore armonico consiste di rotazioni ellittiche.

\end{esempio}

Vediamo un po' delle proprietà del flusso

\begin{proposizione}Se $X$ è un campo differenziabile su $\Omega$ completo \begin{itemize}
\item La mappa $\Phi^X_0$ è la mappa identità su $\Omega$
\item $\forall t,s\in\mathbb{R}$ si ha $\Phi_t^X \circ \Phi^X_s=\Phi^X_{t+s}$
\item $\forall t\in\mathbb{R}$ la mappa del flusso al tempo $t$ $\Phi^X_t$ è invertibile e la sua inversa è la mappa del flusso al tempo $-t$: $(\Phi^X_t)^{-1}=\Phi^X_{-t}$
\item $\forall t\in\mathbb{R}$ la mappa del flusso al tempo $t$ $\Phi^X_t$ è un diffeomorfismo
\end{itemize}
\end{proposizione}
\begin{dimo}$ $ \begin{itemize}
\item Per definizione $\Phi_0^X(z_0)=z(0,z_0)=z_0$ per tutti i dati iniziali quindi è la mappa identità
\item La proprietà significa che per tutti i dati iniziali
\[
\Phi^X_t(\Phi^X_s(z_0))=\Phi^X_t(z(s,z_0))=z(t,z(s,z_0))
\]
ma essendo l'equazione differenziale autonoma $z(t,z(s,z_0))=z(t+s,z_0)$ da cui si prova 
$\Phi_t^X \circ \Phi^X_s=\Phi^X_{t+s}$
\item Unendo le proprietà precedenti è ovvio che $\Phi_t^X \circ \Phi^X_{-t}=\Phi^X_{0}\equiv\mathbb{I}_\Omega$ che equivale a dire che $\Phi_t^X$ è invertibile con inversa $\Phi_{-t}^X$
\item Dalla proposizione \ref{propMappaFlussoDifferenziabile} segue che sia $\Phi^X_t$ sia la sua inversa $\Phi^X_{-t}$ sono differenziabili e quindi $\Phi^X_t$ è un diffeomorfismo
\end{itemize}
\end{dimo}

\section{Equazioni differenziali non autonome}

Un'equazione differenziale dipendente dal tempo, o non autonoma, è un'equazione differenziale nella forma
\begin{equation}
\dot{z}=X(t,z)
\end{equation}
Passando da equazioni autonome ad equazioni non autonome vengono mantenuti i teoremi di esistenza ed unicità delle soluzioni ma si perde la traslabilità nel tempo delle soluzioni.
In generale cioè $z(t,z_0,t_0)$ e $z(t,z_0,t_1)$ sono soluzioni diverse!
Questo significa che per specificare una soluzione non basterà più indicarne un punto nello spazio $\Omega$  ma servirà anche indicare a quale istante vi passa.

\begin{definizione}\textit{Flusso non autonomo}\\
Per un'equazione differenziale non autonoma definiamo flusso non autonomo la mappa $\mathbb{R}\times\mathbb{R}\times\Omega\rightarrow\Omega$ definita da
\[
(t,t_0,z_0)\mapsto\Phi^X_{t,t_0} := z(t;z_0,t_0)
\]
\end{definizione}

Avendo perso la traslabilità nel tempo le orbite di soluzioni diverse potrebbero incontrarsi in uno o più punti di $\Omega$, l'unicità delle soluzioni garantirà solo che passino per lo stesso punto ad istanti diversi.

Possiamo definire uno spazio delle fasi esteso $\mathbb{R}\times\Omega$ che contenga anche la coordinata temporale e su di esso potremo associare ad un'equazione non autonoma il sistema autonomo
\begin{equation}\label{NonAutonomaAdAutonoma}
\dot{z}=X(t,z) \overset{t\rightarrow \tau}{\longrightarrow}
\left(\begin{array}{c}
\dot{\tau}\\
\dot{z} 
\end{array}\right)=\left(\begin{array}{c}
X(\tau,z) \\
1
\end{array}\right)
\end{equation}
e, essendo questo sistema a tutti gli effetti autonomo, le sue orbite in $\mathbb{R}\times\Omega$ non potranno mai intersecarsi.
Sfortunatamente non potremo usare questo trucchetto per semplificare ogni equazione non autonoma in un'equazione autonoma, per esempio infatti il sistema autonomo di \eqref{NonAutonomaAdAutonoma} non potrà avere equilibri.

\section{Cambi di coordinate}

Vogliamo capire come effettuare efficacemente dei cambiamenti di coordinate nelle equazioni differenziali.

In pratica un cambio di coordinate si può pensare come una mappa invertibile e differenziabile dello spazio delle fasi in sé stesso, che però ne cambia la base. Studieremo quindi come cambia l'equazione differenziale ed il suo flusso sotto un diffeomorfismo che manda lo spazio delle fasi $\Omega$ in uno spazio $\tilde{\Omega}$ avente la stessa dimensione.

\begin{definizione}
Dato un campo vettoriale $X$ su $\Omega$ e un campo $\tilde{X}$ su $\tilde{\Omega}$ diremo che i loro flussi e le loro equazioni differenziali sono coniugate dal diffeomorfismo $C:\Omega\rightarrow\tilde{\Omega}$ se le curve integrali di $\tilde{X}$ sono tutte e le sole immagini sotto $C$ delle curve integrali di $X$.
\end{definizione}


Riportiamo con uno schemino i collegamenti tra i campi vettoriale, i flussi e gli spazi.\\
$ $\\
\begin{minipage}[H]{0.45\linewidth}
\[
\begin{array}{ccccc}
&\Omega & \overset{C}{\longrightarrow} & \tilde{\Omega}&\\
\Phi_t^X & \downarrow & & \downarrow & \Phi_t^{\tilde{X}}\\
& \Omega & \overset{C}{\longrightarrow} & \tilde{\Omega} & \\
\end{array}
\]
\end{minipage}
\begin{minipage}[H]{0.45\linewidth}
\[
\begin{array}{ccccc}
&\Omega & \overset{C}{\longrightarrow} & \tilde{\Omega} & \\
X & \downarrow & & \downarrow &  \tilde{X}\\
& T_z\Omega &  & T_{\tilde{z}}\tilde{\Omega} & \\
\end{array}
\]
\end{minipage}

Ricordiamo che il campo vettoriale ha valori nello spazio (fibrato) tangente e non semplicemente in $\mathbb{R}^n$ come sarebbe comodo pensare, quindi non potremo dedurre immediatamente la scrittura del campo $\tilde{X}$ a partire da $X$ come faremmo per una funzione. 

\begin{proposizione}\label{propCampoConiugato}
Dati il campo vettoriale $X$ su $\Omega$ ed il diffeomorfismo $C:\Omega\rightarrow\tilde{\Omega}$ esiste un unico campo vettoriale $\tilde{X}$ su $\tilde{\Omega}$ coniugato ad $X$ da $C$ ed è
\[
\tilde{X}:\tilde{\Omega}\rightarrow T_{\tilde{z}}\tilde{\Omega} \hspace{0.15\linewidth} \tilde{X}=(C'X)\circ C^{-1}
\]
\end{proposizione}
\begin{dimo}
Sia $t\mapsto z_t$ una curva integrale di $X$, allora $t\mapsto \tilde{z}_t=C(z_t)$ dovrà essere una curva integrale di $\tilde{X}$. Derivandola otteniamo
\[
\tilde{X}(\tilde{z}) = \dfrac{\partial}{\partial t} \tilde{z}_t = C' \dot{z}_t = C'X(z_t)
\]
Vogliamo che il campo a destra sia un $\tilde{\Omega}\rightarrow\mathbb{R}^n$ e non un $\Omega\rightarrow\mathbb{R}^n$ quindi, da $z_t=C^{-1}(\tilde{z}_t)$, componiamo con $C^{-1}$ ottenendo
\[
\tilde{X}(\tilde{z}) = C'X (C^{-1}(\tilde{z}_t))\hspace{0.15\linewidth}\tilde{X}=(C'X)\circ C^{-1} \qedhere
\]
\end{dimo}

Questa formula per il campo $\tilde{X}$ è comoda nella teoria ma nella pratica sarà più semplice ripetere il procedimento seguito nella dimostrazione, derivando $\tilde{z}$ lungo una curva integrale con la regola della catena, evitandoci di calcolare esplicitamente i prodotti tra matrici e le matrici inverse.

\subsection{Sollevamento tangente}\label{sezSollevamentoTangente}
Consideriamo una mappa differenziabile tra varietà $\Psi:M\rightarrow P$, in generale di dimensioni diverse.
$\Psi$ induce una legge di trasformazione dei vettori tangenti alle due varietà, infatti per definizione un vettore tangente ad un punto di $z$ di $M$ si può esprimere come
\[
\gamma(0)=z \hspace{0.15\linewidth} v=\dot{\gamma}(t) 
\]
ed un vettore tangente ad un punto $\tilde{z}$ di $P$ come
\[
\phi(0)=\tilde{z}=\Psi(\gamma(0)) \hspace{0.15\linewidth} \tilde{v}=\dot{\phi}(t)
\]
ed allora $\dfrac{d}{dt} \left(\Psi\circ\gamma\right)(t)|_{t=0}$ è tangente a $P$ in $\Psi(z)$. Esplicitando
\[
\tilde{v}=\Psi'(\gamma(0))\dot{\gamma}(0)=\Psi'(z) v
\]
Ad ogni vettore $v\in T_zM$ viene associato un vettore $\tilde{v}=\Psi'(z)v\in T_{\Psi(z)}P$, cioè la mappa $\Psi:M\rightarrow P$ induce una famiglia di mappe lineari fra gli spazi tangenti
\[
T_zM\rightarrow T_{\Psi(z)}P \hspace{0.15\linewidth} \quad \forall z\in M
\]
Possiamo considerare questa diverse mappe come un' unica $T\Psi$ tra i due fibrati tangenti definita da
\begin{equation}
T\Psi:TM\rightarrow TP \hspace{0.15\linewidth} T\Psi(z,v)=(\Psi(z),\Psi'(z)v)
\end{equation}
che chiameremo \textbf{sollevamento tangente} della mappa $\Psi$.

Notiamo che se $\Psi$ è un diffeomorfismo anche $T\Psi$ è un diffeomorfismo.
\par
Tornando nello specifico ai cambi di coordinate, $C:\Omega\rightarrow\tilde{\Omega}$ indurrà una $TC:T\Omega\rightarrow T\tilde{\Omega}$ e i vettori negli spazi tangenti saranno associati secondo $v\mapsto C'(z)v$.
\[
\begin{array}{ccccc}
 & \Omega & \overset{C}{\longrightarrow} & \tilde{\Omega} &  \\
X &\downarrow &  & \downarrow & \tilde{X} \\
 & T_z\Omega & \xrightarrow[]{v\mapsto vC'(z)} & T_{\tilde{z}}\tilde{\Omega}& \\
 \\
 & T\Omega & \overset{TC}{\longrightarrow} & T\tilde{\Omega} &
\end{array}
\]


\subsection{Push-forward e pull-back}
Dato un diffeomorfismo $C:\Omega\rightarrow\tilde{\Omega}$ per una normale funzione $f:\Omega\rightarrow\mathbb{R}$ chiamiamo push-forward di $f$ la funzione $\tilde{\Omega}\rightarrow\mathbb{R}$
\[
C_{*}f=\tilde{f}=f\circ C^{-1}
\]

Analogamente con i campi vettoriali $X$ e $\tilde{X}$ coniugati da $C$ chiameremo il campo $\tilde{X}$ il \textbf{push-forward} di $X$ e $X$ il \textbf{pull-back} del campo $\tilde{X}$\footnotemark[2].

In generale il push-forward di un campo vettoriale $X$ tramite un diffeomorfismo $C$ è
\begin{equation}\label{push-forward}
C_{*}X=(C'X)\circ C^{-1}
\end{equation}
e il suo pull-back
\begin{equation}\label{pull-back}
C^{*}X=((C^{-1})'\tilde{X})\circ C
\end{equation}

Schematicamente le relazioni sono

\[
\begin{array}{ccccc}
 & \Omega & \overset{C}{\longrightarrow} & \tilde{\Omega} &  \\
X\equiv C^{*}\tilde{X} &\downarrow &  & \downarrow & C_{*}X\equiv\tilde{X} \\
 & T_z\Omega & \xrightarrow[]{v\mapsto vC'(z)} & T_{\tilde{z}}\tilde{\Omega}& \\
\end{array}
\]

è poi abbastanza ovvio quindi che $C^{*}C_{*}X=X$.


\footnotetext[2]{
Il diffeomorfismo che porta $\tilde{\Omega}$ in $\Omega$ è $C^{-1}$ e non $C$ quindi dalla \eqref{pull-back} troviamo
\[
C^{*}\tilde{X}=(C'X)\circ C^{-1}\equiv X
\]
quindi effettivamente $X$ è il pull-back di $\tilde{X}$
}

\subsection{Cambi di coordinate dipendenti dal tempo}\label{sezCambiCoordinateDipendentiTempo}
Per ora abbiamo trattato cambi di coordinate rappresentati da diffeomorfismi $\Omega\rightarrow\tilde{\Omega}$, cioè essenzialmente non dipendenti dal tempo. A volte però sarà necessario usare dei cambi di coordinate descritti da mappe
\[
C:\mathbb{R}\times\Omega\rightarrow\tilde{\Omega} \hspace{0.15\linewidth} (t,z)\mapsto C(t,z)
\]
differenziabili in $t$ e dove per ogni $t$ fissato $C_t:\Omega\rightarrow\tilde{\Omega}$ è un diffeomorfismo.

Usando cambi di coordinate dipendenti dal tempo è più naturale considerare equazioni non autonome, infatti anche per equazioni autonome $\dot{z}=X(z)$ il trasformato del campo $\tilde{X}$ ovviamente dipenderà dal tempo.

Come per i cambi non dipendenti dal tempo diremo che $C$ coniuga $X$ e $\tilde{X}$ se le curve integrali di $\tilde{X}$ sono tutte e sole trasformate mediante $C$ delle curve integrali di $X$.
Quindi dati $X$ e $\tilde{X}$ coniugati $ z\mapsto z(t)$ è una curva integrale di $X$ se e solo se $\tilde{z}(t)=C(t,z(t))$ è una curva integrale di $\tilde{X}$.

Da questo, derivando con la regola della catena
\[
\tilde{X}(t,\tilde{z}_t)=\dfrac{\partial}{\partial t} \tilde{z}(t) = \dfrac{\partial C}{\partial t}(t,z_t)+\dfrac{\partial C_t}{\partial z_t} \dot{z}(t)
\]
e quindi
\begin{equation}
\tilde{X}:\mathbb{R}\times\tilde{\Omega}\rightarrow\mathbb{R}^n \hspace{0.15\linewidth}
\tilde{X}=\left[\dfrac{\partial C_t}{\partial t} + \dfrac{\partial C_t}{\partial z}X
\right] \circ C_t^{-1}
\end{equation}

\section{Integrali primi}

Come abbiamo già visto un insieme $M\subset\Omega$ si dice invariante per il flusso di un'equazione differenziale $\dot{z}=X(z)$ se ogni soluzione dell'equazione con dato iniziale in $M$ rimane in $M$
\[
z\in M \implies \Phi_t^X(z)\in M \quad \forall t
\]
e nel caso l'invarianza valga non per tutti i $t$ ma solo per $t>0$ la chiameremo invarianza nel futuro mentre se vale solo per $t<0$ la chiameremo invarianza nel passato.

Chiaramente quindi, dato che stiamo considerando equazioni autonome, un insieme invariante contiene completamente ogni orbita che lo interseca.

\begin{definizione}\textit{Integrale primo}\\
Una funzione $f:\Omega\rightarrow\mathbb{R}$ è un integrale primo del campo vettoriale $X$ su $\Omega$ e della sua equazione differenziale $\dot{z}=X(z)$ se i suoi insiemi di livello sono invarianti per il flusso di $X$.
\end{definizione}

Dire che $f$ è un integrale primo per $X$ equivale a richiedere
\[
f\circ\Phi^X_t=f \hspace{0.15\linewidth} f(z(t,z_0))=f(z_0) \quad \forall t , \forall z_0
\]
cioè che $f$ dia sempre lo stesso valore lungo tutta un'orbita: l'insieme di livello è l'insieme in cui $f$ assume un solo valore e se questo insieme è un invariante tutte le orbite che vi passano vi rimangono.

\begin{esempio}
Per l'oscillatore armonico $m\ddot{x}=-k x$ con $k>0$ l'energia $E(x\dot{x})=\sfrac{1}{2} m\dot{x}^2 +\sfrac{1}{2}kx^2$ è un integrale primo.

Proviamo a verificarlo derivando rispetto al tempo lungo una soluzione $t\mapsto x_t$
\[\dfrac{d}{dt}E(x,\dot{x})=m\dot{x}_t\ddot{x}_t+k x_t \dot{x}_t = (m\ddot{x}_t +k x_t )\dot{x}_t =0
\]
Quindi lungo le orbite $E$ è cotante e quindi è un integrale primo. Questo significa che le orbite dell'equazioni sono tutte contenute in insiemi di livello dell'energia. Ma gli insiemi di livello dell'energia sono ellissi nel piano $(x,\dot{x})$ e, esclusa l'origine, dato che queste ellissi non contengono equilibri ciascuna di esse è un'orbita.
\end{esempio}


Gli integrali primi permettono di suddividere lo spazio delle fasi e quindi permettono un'analisi qualitativa del ritratto in fase molto immediata.

\begin{proposizione}[Integrali primi e attrattività]\label{propIntegraliPrimiAttattività}
Un'equazione differenziale che ha un equilibrio attrattivo non possiede nessun integrale primo che non sia costante in un intorno dell'equilibrio
\end{proposizione}
\begin{dimo}
Sia $f$ un integrale primo e $\overline{z}$ un equilibrio. Allora $\exists V$ intorno di $\overline{z}$ per cui $\lim_{t\rightarrow +\infty}z_t=\overline{z}$ per ogni soluzione $z_t$ passante per $V$. Un integrale primo è continuo in $\Omega$ quindi $\lim_{t\rightarrow +\infty}f(z_t)=f(\overline{z})$ ed essendo $f$ costante sulle soluzione $f(z)=f(\overline{z})$ $\forall z\in V$. \qedhere
\end{dimo}

\subsection{Derivata di Lie}

Un integrale primo è utile quando non si conoscono direttamente le soluzioni quindi vogliamo un metodo per verificare se una data funzione $f$ è un integrale primo che non usi il calcolo diretto dalla definizione.

Consideriamo una curva integrale $t\mapsto z_t$ del campo $X$ e dato il nostro candidato integrale primo $F$ allora lungo l'orbita avremo la funzione composta $t\mapsto f(z_t)$ e se la deriviamo con la regola della catena troviamo
\[
\dfrac{\partial}{\partial t}f(z_t)=f'(z_t)\dot{z}_t=f'(z_t)X(z_t)=X(z_t)\cdot\nabla f(z_t)\footnotemark
\]
\footnotetext{Quando scriviamo $f'X$ stiamo considerando il prodotto righe per colonne della matrice Jacobiana di $f$, avente $1$ riga e $n$ colonne, per la matrice di $X$ data da $n$ righe e $1$ colonna. Questo prodotto è uguale al prodotto scalare dei due vettori colonna con $n$ componenti $X$ e $\nabla f$.}
da questo definiamo
\begin{definizione}\textit{Derivata di Lie}\\
Chiamiamo derivata di Lie di una funzione $f:\Omega\rightarrow\mathbb{R}$ associata ad un campo vettoriale $X$ su $\Omega$ la funzione
\[
L_Xf:\Omega\rightarrow\mathbb{R} \hspace{0.15\linewidth}
L_Xf=f'X=X\cdot\nabla f
\]
\end{definizione}


\begin{proposizione}
Siano $X$ un campo vettoriale su $\Omega$ e $f$ una funzione su $\Omega$ allora
\begin{enumerate}[i.]
\item La derivata di $f$ lungo una qualsiasi curva integrale $t\mapsto z_t$ sarà
\[
\dfrac{\partial}{\partial t}f(z_t)=L_Xf\,(z_t) \quad \forall t
\]
ovvero
\[
\dfrac{\partial}{\partial t} ( f\circ\Phi^X_t )= L_Xf \Phi^X_t \quad \forall t
\]
\item $f$ è un integrale primo di $X$ se e solo se $L_Xf\equiv 0$
\end{enumerate}
\end{proposizione}
\begin{dimo}
Derivando lungo una curva integrale con la regola della catena troviamo  
\[
\dfrac{\partial}{\partial t}f(z_t)=f'(z_t)\dot{z}_t=f'(z_t)X(z_t)=L_Xf
\]
\begin{itemize}
\item Sia $L_Xf\equiv 0$. 
Allora per quanto detto nella prima parte della proposizione significa che $f$ è costante lungo ogni orbita e quindi che $f$ è un integrale primo.
\item  Sia $f$ un integrale primo. 
Allora dovrà essere costante lungo una qualsiasi orbita $t\mapsto z_t$ ad ogni tempo 
\[
L_Xf\left(z(t;z_0)\right)=0 \quad \forall t, \forall z_0
\]
quindi prendendo anche per $t=0$ che dà $L_Xf\left(z(0;z_0)\right)=L_Xf(z_0)=0$ $\forall z_0$ cioè $L_Xf\equiv 0$.\qedhere
\end{itemize}
\end{dimo}

\begin{esempio}
Per l'oscillatore armonico $m\ddot{x}=-kx$ possiamo verificare se l'energia è un integrale primo con la derivata di Lie.

Il campo per il sistema del primo ordine dell'oscillatore armonico è $X(x,v)=\left(\begin{array}{c} v \\ -\dfrac{k}{m} x \end{array}\right)$ quindi la derivata di Lie di $E(x\dot{x})=\sfrac{1}{2} m\dot{x}^2 +\sfrac{1}{2}kx^2$ è
\[
L_x E= X\cdot\nabla E(x,\dot{x}) =\sfrac{1}{2}\left(\begin{array}{c}
V \\ \dfrac{-kx}{m} \end{array}\right) \cdot \left(\begin{array}{c}
kx \\ mv \end{array}\right) = \sfrac{1}{2} (vkx - \dfrac{k}{m}xmv)=0
\]
\end{esempio}

\subsection{Foliazione e abbassamento dell'ordine}\label{sezFoliazione}

Come sappiamo se l'integrale primo $f$ di un campo vettoriale $X$ su $\Omega$ è tale che il suo gradiente non si annulli in nessun punto di $\Omega$
\[
\nabla f (z)\neq 0\quad \forall z\in\Omega
\]
i suoi insiemi di livello sono varietà differenziali di dimensione $n-1$.

Per ciascun punto dello spazio delle fasi passerà un unico insieme di livello di $f$, diremo che la presenza dell'integrale primo produce una $\textbf{foliazione}$ dello spazio delle fasi in sottovarietà invarianti di dimensione $n-1$.

Se però il campo ha $k>1$ integrali primi ciascuno di loro avrà le stesse proprietà e quindi ogni orbita dovrà appartenere completamente ad un'intersezione di loro insiemi di livello
\[
\left\{ z\in\mathbb{R}^n:f_1(z)=\alpha_1,...,f_k(z)=\alpha_k \right\}
\]
Se gli integrali primi sono \textbf{funzionalmente indipendenti} in $z$ potremo definire una  
\[
F(z):\Omega\rightarrow\mathbb{R}^k \qquad F(z)=\left(f_1(z),...,f_k(z)\right)
\] 
che sarà una sommersione in $z$ $\text{rk}(F'(z))=k$. Se quindi gli integrali primi sono ovunque indipendenti gli insiemi dati dalle intersezioni degli insiemi di livello
\begin{equation}
\mathcal{Z}_\alpha=\{ z\in\Omega:F(z)=\alpha\in\mathbb{R}^k \}=\left\{ z\in\Omega:f_1(z)=\alpha_1,...,f_k(z)=\alpha_k \right\}
\end{equation}
saranno sottovarietà invarianti di dimensione $n-k$ dello spazio delle fasi.

Questo è particolarmente utile all'aumentare di $k$:
\begin{itemize}
\item Per $k=n$ le intersezioni sono solo singoli punti e quindi tutte le orbite sono singoli equilibri: questo può verificarsi solo per $X=0$.
\item per $k=n-1$ le intersezioni saranno curve e quindi potremo direttamente individuare orbite ed equilibri.
\end{itemize}

In particolare lavorando in $\Omega\subseteq\mathbb{R}^2$ basterà individuare un integrale primo per poter rappresentare direttamente le orbite a partire dai suo insiemi di livello.

\begin{proposizione}\textit{Abbassamento dell'ordine}\\
Per l'equazione differenziale $\dot{z}=X(z)$ su $\Omega\subseteq\mathbb{R}^n$ dati $k$ integrali primi $f_1,...,f_k$ funzionalmente indipendenti in un punto $z\in\Omega$ esiste un sistema di coordinate
\[
(u,y)=(u_1,...,u_k,y_1,...,y_{n-k})
\]
definito in un intorno di $z$ col quale l'equazione differenziale prende la forma
\[
\begin{cases}
\dot{u}=0 \\
\dot{y}=Y(u,y)
\end{cases}
\]
per qualche funzione $Y=(Y_1,...,Y_{n-k})$.
\end{proposizione}

\subsection{Equazione di Newton unidimensionale}\label{subsezNewton1}

L'equazione di Newton è un equazione del secondo ordine con campo conservativo
\begin{equation}\label{eqNewtonUnidim}
m\ddot{x}=-\dfrac{dV}{dx}(x) \qquad \quad x\in\mathbb{R}
\end{equation}

Questa equazione descrive il moto di una particella di massa $m$ che si muove senza attrito lungo una guida liscia e rettilinea, sotto l'effetto di forza conservative con potenziale $V(x)$.

In realtà un'ampia classe di problemi meccanici integrabili si riduce allo studio di questa equazione.

L'equazione \eqref{eqNewtonUnidim} ha l'integrale primo dell'energia $E(x,v)=\dfrac{1}{2}mv^2+V(x)$ e, essendo il problema unidimensionale, sarà sufficiente disegnarne gli insiemi di livello
\[
C_e=\{(x,v)\in\mathbb{R}^2:E(x,v)=e\}
\]

Ricordiamo che $C_e\subseteq \mathbb{R}^{2}$ e che le orbite saranno date dalle sue cerve di livello e dagli equilibri.


Possiamo disegnare facilmente i $C_e$ a partire dal grafico della $V$ tramite un semplice procedimento basato su fatti noti. 
\begin{enumerate}[i.]
\item L'energia cinetica $\sfrac{1}{2}\,mv^2$ è non negativa quindi necessariamente $E(x,v)\geq V(x)$ $\forall x,v$. 
\item Fissando $E(x,v)=e$, con $e>\text{min}V(x)$, l'insieme $C_e$ si proietterà sull'asse $x$ nel sottoinsieme
\[
J_e=\{ x\in\mathbb{R}:V(x)\leq e \}
\]
\item dato che $\dfrac{\partial E}{\partial x}=\dfrac{\partial V}{\partial x}=V$ e $\dfrac{\partial E}{\partial v}=mv$ $C_e$ è una curva regolare ad eccezione che negli equilibri $(\overline{x},0)$ di energia $e$, che inoltre soddisferanno $V(\overline{x})=e$.
\item Nei sottoinsiemi $V(x)<e$ di $C_e$ questo consiste di due rami, simmetrici rispetto all'asse $x$, di equazioni cartesiane $v=\pm \tilde{v}(x,e)$ con
\begin{equation}
\tilde{v}(x,e):=\sqrt{\dfrac{2}{m}(e-V(x)}
\end{equation}
Questi due rami si incontreranno nei punti $V(x)=e$ dell'asse $x$ con quindi $\tilde{v}=0$. Questi punti possono essere equilibri, in cui la $C_e$ non è una curva regolare, o punti non di equilibrio in cui quindi la $C_e$ è regolare e la sua tangente quindi è verticale.
\item Il verso di percorrenza delle orbite è verso detra nel semipiano superiore e verso sinistra nel semipiano inferiore.
\end{enumerate}

\begin{esempio}[Ricostruiamo alcune curve di livello particolari di $C_e$]
Considerando una componente connessa di $J_e$ alla volta:\begin{itemize}[label=$\cdot$]
\item Un minimo stretto del potenziale è un equilibrio. In particolare $\overline{x}\in J_e$ tale che $e=E(\overline{x},0)=V(\overline{x})$.
\item Fissando $e$ con $C_e$ regolare:
\begin{enumerate}[i.]
	\item Se $J_e$ è un intervallo chiuso e limitato con estremi in cui $V(x)=e$. La curva di livello di $C_e$ corrispondente sarà una curva regolare chiusa, unione dei due rami $v=\pm\tilde{v}(x,e)$ che si raccorderanno nei due estremi di $J_e$ con tangente verticale. Non contiene punti di equilibrio è un'unica orbita, percorsa periodicamente. Chiameremo i due estremi di $J_e$ \textbf{punti di inversione} perchè in essi la velocità si annulla e cambia segno. Questa situazione si incontra per energie ``poco'' superiori di un minimo dell'energia (che invece corrispone ad un equilibrio). Un minimo di energia è sempre circondata orbite chiuse periodiche.
	\item Se $J_e$ è una semiretta con $V(x)=e$ solo nel suo estremo la situazione è simile alla precedente: i due rami della $C_e$ si incontreranno nell'estremo e ad infinito con tangente verticale, anche in questo caso $C_e$ è una sola orbita ma questa non sarà periodica. 
	\item Se $J_e$ è l'intera retta e $V(c)\leq e$ in ogni punto allora la $C_e$ consiste di due rami disgiunti. Questa situazione spesso si trova ``sopra'' (cioè a velocità maggiori) le due precedenti.
\end{enumerate}
\item UN massimo stretto del potenziale è un equilibrio. Sia $V(\overline{x})=e$ allora localmente $V(x)<e$ sia a destra sia a sinistra e l'equilibrio divide in quattro le orbite che lo circondano (è una sella).
\end{itemize}
Mettendo insieme questi casi particolari di $J_e$ si può costruire rapidamente un ritratto in fase qualsiasi.
\end{esempio}

\subsubsection{Equazione ristretta per l'equazione di Newton unidimensionale}

Possiamo sfruttare l'integrale primo per abbassare la dimensione ottenendo un'equazione ristretta integrabile direttamente.
Ogni insieme $C_e$ non critico, esclusi eventuali intersezioni con l'asse $x$, è dato da due rami simmetrici rispetto all'asse $x$ ottenuti da 
\[
x\mapsto\tilde{v}(x,e)=\sqrt{\dfrac{2}{m}\left(e-V(x)\right)}
\]

In ciascun semipiano $v>0$, $v<0$ pensiamo quindi di usare $(E,x)$ come coordinate per la foliazione invariante data dall'energia. In questo modo la $x$ fornisce una coordinata sul ramo $C_e$ del sottospazio.

Consideriamo per esempio il semipiano superiore: il campio di coordinate e la sua inversa saranno
\[
(x,v)\mapsto\left(E(x,v),x\right) \qquad (E,x)\mapsto(x,\tilde{v}(x,E))
\]
Derivando lungo le soluzioni l'inversa troviamo che il sistema $\dot{x}=v$ $\dot{v}=-V'(x)$ diventa
\begin{equation}
\dot{E}=0 \qquad \dot{x}=\tilde{v}(x,E)
\end{equation}

Quindi per ogni $E=e$ fisatto la seconda equazione sarà ristretta nel ramo di $C_e$ del semipiano. 

Nel semipiano inferiore le cose saranno analoghe, con un segno meno nella seconda equazione.

Consideriamo una soluzione avente dato iniziale $(x_0,v_0)$ ad un tempo $t_0$, assumendo per semplicità $v_0>0$.
Allora per $E(x_0,v_0)=e$ la $t\mapsto x_t$ di una soluzione dovrà soddisfare l'equazione ristretta $\dot{x}_t=\sqrt{\dfrac{2}{m}\left(e-V(x_t)\right)}$ fino a quando la velocità non si annulla. Vi è allora un intervallo di $t_0$ in cui la $t\mapsto x_t$ è monotona ed invertibile, la sua inversa soddisferà
\[
\dfrac{dt_x}{dx}=\sqrt{\dfrac{\sfrac{m}{2}}{e-V(x)}}
\]
ed integrandola otteniamo
\begin{equation}
t_x=t_0+\sqrt{\dfrac{m}{2}}\int_{x_0}^{x}\dfrac{dx}{\sqrt{e-V(x)}}
\end{equation}
Questa ovviamente sarà valida solo nell'intervallo in cui la velocità non si annulla. 

La presenza dell'integrale primo permette di integrare l'equazione, a meno dell'inversione di una funzione. Diremo che questo problema è risolto alle \textbf{quadrature}.

\begin{esempio}[Periodo dei moti limitati]
Consideriamo per esempio i moti periodici dell'esempio precedente. Le loro orbite toccheranno l'asse $x$ nei punti di inversione $x_\pm(e)$ ed il loro periodo sarà il doppio del tempo necessario per arrivare ad un punto di inversione partendo dall'altro, cioè
\[
T(e)=\sqrt{2m}\int_{x_-(e)}^{x_+(e)} \dfrac{dx}{\sqrt{e-V(x)}}
\]
Notiamo che il periodo è una funzione della sola energia e quindi tutti i moti limitati aventi la stessa energia saranno \textbf{isocroni} indipendentemente dai dati iniziali. Diremo che il sistema è \textbf{asinocrono}.
\end{esempio}

\section{Stabilità degli equilibri}

Fin'ora abbiamo parlato di stabilità solo in relazione all'attrattività e corrispondemente di instabilità solo in relazione alla repulsività. Esiste però un altro concetto di stabilità che riguarda punti fissi che, pure senza proprietà attrattive, hanno la proprietà che le orbite ad essi vicine rimangano vicine.

\begin{definizione}[Stabilità secondo Lyapunov]
Un punto di equilibrio $\overline{z}$ di una $\dot{z}=X(z)$ può essere
\begin{itemize}
\item Stabile se per ogni intorno $U$ di $\overline{z}$ esiste un secondo intorno $U_0$ di $\overline{z}$ per cui
\[
\Phi_t^X(U_0)\subseteq U \quad \forall t>0
\]
\item Stabile per tutti i tempi se per ogni intorno $U$ di $\overline{z}$ esiste un secondo intorno $U_0$ di $\overline{z}$ per cui
\[
\Phi_t^X(U_0)\subseteq U \quad \forall t\in\mathbb{R}
\]
\item Instabile se non è stabile secondo Lyapunov
\item Asintoticamente stabile se è sia stabile che attrattivo
\end{itemize}
\end{definizione}

I due intorni $U$ e $U_0$ nella definizione svolgono un ruolo simile ai due intorni nella definizione di continuità di una funzione: ``per ogni intorno grande esiste un intorno piccolo per cui gli evoluti dei punti dell'intorno piccolo rimangono sempre nell'intorno grande''.

Stabilità, instabilità e stabilità asintotica sono proprietà topologiche e quindi sono invarianti per cambiamenti di coordinate.


La determinazione della stabilità o instabilità di un dato equilibrio di un'equazione in generale è un problema molto complicato, essenzialmente studieremo due metodi per risolverlo: uno che costruisce insiemi invarianti mediante apposite funzioni e l'altro che utilizza le proprietà spettrali della linearizzazione del campo vettoriale.

\subsubsection{Stabilità per equazioni del secondo ordine}

Per equazioni del secondo ordine gli equilibri sono della forma $(x,v)=(\overline{x},0)$ e diremo quindi che una configurazione di equilibrio $\overline{x}$ è stabile (o instabile o asintoticamente stabile) se il corrispettivo equilibrio $(\overline{x},0)$ è stabile (o instabile o asintoticamente stabile).
Quindi per esempio $\overline{x}\in\Omega$ se per ogni intorno $U\subset T\Omega$ di $(\overline{x},0)$ esiste un intorno $U_0\subset T\Omega $ di $(\overline{x},0)$ per cui
\[
(x,v)\in U_0 \implies (x_t,v_t)\in U \> \forall t>0 
\]



\subsection{Metodo delle funzioni Lyapunov}

Consideriamo una funzione $\mathcal{F}$ che sia continua nello spazio delle fasi. Allora i suoi insiemi di sopralivello e sottolivello $\{z:\mathcal{F}(z)<c\}$ saranno aperti. Se questa funzione è tale da avere un minimo (massimo) in un equilibrio $\overline{z}$ allora i suoi insiemi di sottolivello (sopralivello) si stringeranno attorno all'equilibrio e formeranno una sorta di base per gli intorni dell'equilibrio. Allora
\begin{enumerate}
\item Se $L_X\mathcal{F}=0$ allora $\mathcal{F}$ è un integrale primo e le soluzioni saranno costrette a rimanere all'interno del loro insieme di sottolivello e quindi se partono in un intorno dell'equilibrio, che corrisponderà ad un insieme di sottolivello, non potranno allontanarsi.
\item Se $L_X\mathcal{F}\leq 0$ allora $\mathcal{F}$ sarà non crescente lungo le soluzioni nel senso dei tempi positivi, cioè al crescere del tempo le soluzioni dovranno rimanere nel loro intorno o al massimo andare a regioni con minore $\mathcal{F}$. Quindi se $\mathcal{F}$ ha un minimo stretto in un equilibrio le soluzioni che vi passano vicine saranno costrette a ``non allontanarsi''.
\item Se $L_X\mathcal{F} < 0$ allora $\mathcal{F}$ sarà decrescente lungo le soluzioni nel senso dei tempi positivi, cioè al crescere del tempo le soluzioni dovranno andare verso regioni a minor valore di $\mathcal{F}$. Quindi se $\mathcal{F}$ ha un minimo relativo in un equilibrio allora le soluzioni vicine saranno costrette ad andare nell'equilibrio.
\end{enumerate}

A partire da questa idea potremo enunciare uno dei teoremi di Lyapunov.

\begin{proposizione}[Secondo teorema di Lyapunov]\label{secondoTeoLyapunov}
Sia $\overline{z}$ un equilibrio di $\dot{z}=X(z)$, $V$ un intorno di $\overline{z}$ e $\mathcal{W}:V\rightarrow\mathbb{R}$ una funzione differenziabile che ha un minimo relativo stretto in $\overline{z}$. Allora: 
\begin{enumerate}[i.]
\item Se $L_X\mathcal{W}(z)=0$ $\forall z\in V$ allora $\overline{z}$ è stabile per tutti i tempi.
\item Se $L_X\mathcal{W}(z)\leq 0$ $\forall z\in V$ allora $\overline{z}$ è stabile.
\item Se $L_X\mathcal{W}(z)<0$ $\forall z\in V\setminus\{\overline{z}\}$ allora $\overline{z}$ è asintoticamente stabile. 
\end{enumerate}
Una funzione che soddisfa le ipotesi del teorema e  rientra in uno dei $3$ casi si chiama funzione di Lyapunov per l'equilibrio $\overline{z}$.
\end{proposizione}

Per poter dimostrarlo useremo il lemma seguente:

\begin{lemma}\label{lemmaSecondoTeoLyapunov}
Sia $\mathcal{W}:V\rightarrow\mathbb{R}$ differenziabile con minimo stretto in $\overline{z}$ allora per ogni intorno $U$ di $\overline{z}$ esiste un $c\in\mathcal{W}(V)$ con $c>\mathcal{W}(\overline{z})$ tale che $W_{c,\overline{z}}\subset U$, dove $W_{c,\overline{z}}$ è la componente connessa contenente $\overline{z}$ dell'insieme di sottolivello $\mathcal{W}_c=\{z:\mathcal{W}(z)<c\}$.
\end{lemma}
\begin{dimo}[Dimostrazione lemma]
Dato $U$ prendiamo una palla aperta $B$ centrata in $\overline{z}$ sufficientemente piccola perchè sia\begin{itemize}
\item $B\subset U\cap V$
\item $\overline{z}$ sia un punto di minimo assoluto di $\mathcal{W}$ in $\overline{B}$.
\end{itemize}
Sia $c$ il valore minimo di $\mathcal{W}$ su $\partial B$ (che esiste  $\partial B$ è compatta) allora l'insieme $\mathcal{W}_c$ di sottolivello di $c$ conterrà $\overline{z}$ ma non conterrà alcun punto di $\partial B$.
La $\mathcal{W}_{c,\overline{z}}$ parte connessa di $\mathcal{W}_c$ contenente $\overline{z}$ deve essere contenuta in $B$ perché se avesse punti esterni a $B$ dovrebbe anche intersecarne la frontiera, ma quindi $W_{c,\overline{z}}\subset U$. \qedhere
\end{dimo}

Il secondo teorema di Lyapunov ha un corollario che individua equilibri asintoticamente stabili con ipotesi più deboli rispetto a quelle della proposizione. Lo dimostreremo all'interno della dimostrazione del teorema.

\begin{corollario}[Principio di La Salle-Krasovskii]\label{corollario2TeoremaLyapunov}
Sia $V$ un intorno di un equilibrio $\overline{z}$ di $\dot{z}=X(z)$ e $\mathcal{W}:V\rightarrow\mathbb{R}$ una funzione liscia avente minimo stretto in $\overline{z}$. Se $L_X\mathcal{W}\leq 0$ in tutto $V$ e l'insieme $\{ z:L_X\mathcal{W}(z)=0 \}$ non contiene completamente nessun orbita eccetto $\{ \overline{z} \}$, allora $\overline{z}$ è asintoticamente stabile. 
\end{corollario}

\begin{dimo}[Dimostrazione proposizione \ref{secondoTeoLyapunov} e corollario]
Dato un intorno $U$ di $\overline{z}$ $\exists c>\mathcal{W}(\overline{z})$ tale che $\mathcal{W}_{c,\overline{z}}\subset U$. 
Poichè $\mathcal{W}$ è continua $\mathcal{W}_c$ è aperto e quindi $W_{c,\overline{z}}$ è un intorno (le componenti connesse di un aperto sono aperte).

\begin{enumerate}[i.]
\item $L_x\mathcal{W}=0$ quindi $\mathcal{W}_{c,\overline{z}}$ è invariante e dunque ogni soluzione che parte da esso vi rimane per tutti i tempi. Allora $\forall U \> \exists U_0=W_{c,\overline{z}}$ per cui $\phi_t^X(U_0)\subseteq U_0\subset U\> \forall t\in\mathbb{R}$ e quindi $\overline{z}$ è stabile per tutti i tempi.
\item\label{2TeoLyapDim2} $L_x\mathcal{W}\leq0$ quindi $\mathcal{W}_{c,\overline{z}}$ è invariante solo per i tempi positivi ($\mathcal{W}$ è non crescente lungo le soluzioni quindi  le soluzioni che partono da un insieme di sottolivello rimarranno al suo interno per i tempi positivi) e quindi $\forall U \> \exists U_0=W_{c,\overline{z}}$ per cui $\phi_t^X(U_0)\subseteq U_0\subset U\> \forall t>0$ e quindi $\overline{z}$ è stabile.
\item Da \ref{2TeoLyapDim2} sappiamo che $\overline{z}$ è stabile. 
Dimostriamo ora che l'equilibrio è attrattivo per $L_X\mathcal{W}\equiv 0$ ma volutamente usando solo ipotesi più deboli in modo da dimostrare anche il corollario \ref{corollario2TeoremaLyapunov}.

Sia $L_X\mathcal{W}\leq 0$ in tutto $V$ e supponiamo che se $y\in V$ è tale che $\mathcal{W}(\Phi_t^X(y))=\mathcal{W}(y)$ per tutti i $t\geq0$ allora $y=\overline{z}$ (cioè che l'unica orbita completamente in $\{z:L_X\mathcal{W}(z)=0\}$ è $\{\overline{z}\}$).
Chiaramente se $L_X\mathcal{W}(y)<0$ $\forall y\neq\overline{z}$ la seconda ipotesi è soddisfatta.

Scegliamo $U$ come uno degli intorni $\mathcal{W}_{c,\overline{z}}$ con $c>\mathcal{W}(\overline{z})$ in modo che $\overline{U}$ sia limitata e contenuta in $V$, che sappiamo essere possibile dal lemma.
Allora per \ref{2TeoLyapDim2} $U$ è invariante per i tempi positivi.

Sia $z_0\in U$ e $z_0\neq\overline{z}$, allora la funzione $t\mapsto\mathcal{W}\left(\Phi_t^X(z_0)\right)$ è una funzione decrescente da $\mathbb{R}_+$ in $\mathbb{R}$ con minimo $\mathcal{W}(\overline{z})$ e quindi esiste
\[
w:=\lim_{t\rightarrow+\infty}\mathcal{W}(\Phi_t^X(z_0))\geq \mathcal{W}(\overline{z})
\]
Dato che la soluzione con dato iniziale $z_0$ apparterrà ad $U$ per tutti i tempi positivi e che $\overline{U}$ è compatto esisterà $y\in\overline{U}$ e una successione $t_n\rightarrow +\infty$ tali che 
\[
\lim_{n\rightarrow\infty} \Phi_{t_n}^X(z_0)=y
\]
e quindi per la continuità di $\mathcal{W}$ e del flusso $\mathcal{W}(y)=\mathcal{W}\Phi_{\lim_{n\rightarrow\infty}t_n}^X(y)=w$.

Per la continuità del flusso si avrà, per ogni $t\geq0$
\[
\Phi_t^X(y)=\Phi_t^X\left( \lim_{n\rightarrow\infty}\Phi_{t_n}^X(z_0) \right)=\lim_{n\rightarrow\infty} \Phi_t^X\circ\Phi_{t_n}^X (z_0)=\lim_{n\rightarrow\infty}\Phi_{t+t_n}^X (z_0)
\]
e quindi
\[
\mathcal{W}\left(\Phi_t^X(y)\right)=\mathcal{W} \left( \lim_{n\rightarrow\infty}\Phi_{t+t_n}^X (z_0) \right)=\lim_{n\rightarrow\infty}\mathcal{W} \left(\Phi_{t+t_n}^X (z_0) \right)=w
\]
quindi $\mathcal{W}(y)=\mathcal{W}(\Phi_t^X(y))$ e cioè $y=\overline{z}$, ma $y$ è l'evoluto ad infinito di un qualsiasi $z_0$ preso in $W_{c,\overline{z}}$ e quindi $\overline{z}$ è attrattivo. \qedhere
\end{enumerate}
\end{dimo}

Potremo usare le funzioni di Lyapunov anche per mostrare l'instabilità di un equilibrio: se una $\mathcal{W}$ ha un minimo relativo stretto in $\overline{z}$ e $L_X\mathcal{W}>0$ in $\Omega\setminus\{\overline{z}\}$ allora $\overline{z}$ è instabile.

\begin{esempio}
Studiamo la stabilità dell'origine per l'equazione $\dot{x}=-x-y+x\ln{(1+y)}$ $\dot{y}=x-y-xy$. 

Usiamo come funzione di Lyapunov $\mathcal{W}=\sfrac{1}{2}(x^2+y^2)$, differenziabile e con minimo stretto nell'origine.

La derivata di Lie di $\mathcal{W}$ è
\[
L_X \mathcal{W}=-x^2-y^2+xy^2+x^2\ln{(1+y)}
\]
Vediamo che in un intorno dell'origine $ln{(1+y)}=y+o(t)$ e quindi $L_X \mathcal{W}=-x^2-y^2+xy(y+x)+O(x^2y)=-x^2-y^2+O_3$ quindi la derivata di Lie sarà negativa in ogni intorno bucato dell'origine. Questo ci dice che l'origine è asintoticamente stabile.
\end{esempio}

\subsection{Metodo spettrale}

Potremo ottenere informazioni quasi complementari a quelle che ci danno le funzioni di Lyapunov tramite l'analisi dello spettro della linearizzazione attorno all'equilibrio.
Questo metodo fornisce delle condizioni sufficienti alla stabilità asintotica e condizioni necessarie per la stabilità non asintotica.

\begin{proposizione}[Primo teorema di Lyapunov]
Sia $\overline{z}$ un equilibrio dell'equazione differenziale $\dot{z}=X(z)$ con $z\in\mathbb{R}^n$. Allora: \begin{itemize}[i.]
\item Se la jacobiana calcolata in $\overline{z}$ del campo ha tutti gli autovalori con parte reale negativa allora $\overline{z}$ è asintoticamente stabile.
\[
\alpha_j<0 \> \forall \alpha_j\in[\alpha_1,...,\alpha_k,\alpha_{2k+1},...,\alpha_n] \implies \overline{z} \text{ asintoticamente stabile}
\]
\item Se la jacobiana calcolata in $\overline{z}$ del campo ha almeno un autovalore con parte reale positiva allora $\overline{z}$ è instabile.
\[
\exists\alpha_j\in[\alpha_1,...,\alpha_k,\alpha_{2k+1},...,\alpha_n] \> : \alpha_j>0 \implies \overline{z} \text{ instabile}
\]
\end{itemize}
\end{proposizione}
\begin{dimo}[Dimostrazione del primo punto]
Dimostriamo solo il primo punto della proposizione, lo facciamo in due modi diversi per essere chiari.
\begin{itemize}
\item Il \textit{Teorema di Grobman-Hartman} ci assicura che vicino ad un equilibrio iperbolico l'equazione si comporti come la sua linearizzazione e quindi, essendo $\overline{z}$ iperbolico ed asintoticamente stabile per la linearizzazione, quindi $\overline{z}$ sarà asintoticamente stabile.
\item Sia la matrice Jacobiana $\dfrac{\partial X}{\partial z}(\overline{z})$ diagonalizzabile. Allora possiamo usare quanto visto nella sezione \ref{sezMatriciDiagonalizzabili} e quindi $X'(\overline{z})$ avrà $k$ coppie di autovalori complessi coniugati e $n-2k$ autovalori reali e quindi esiste una matrice invertibile $P$ (cambio di base nelle coordinate degli autovettori) per cui la matrice $A=P^{-1}X(\overline{z})P$ è in forma diagonale a blocchi
\[
A=\left(\begin{array}{cccccc}
B_1 & & & & & \\
 & \ddots  & & & & \\
 &  & B_k & & & \\
 &  & & \alpha_{2k+1} & & \\
 &  & & & \ddots & \\
 &  & & & & \alpha_{n} \\
\end{array}\right) \hspace{0.15\linewidth}
B_j=\left(\begin{array}{cc}
\alpha_j & \beta_j \\
-\beta_j & \alpha_j
\end{array}\right)
\]
Possiamo usare le coordinate $y=P^{-1}(z-\overline{z})=:C(z)$ in cui l'equazione diventa
\[
\dot{z}=X(z) \rightarrow \dot{y}=Y(y)=C'X\circ C^{-1}(y) = P^{-1}X\left(\overline{z}+P(y)\right)
\]
perchè $Y=C_*X$. Nelle coordinate $y$ l'equilibrio $\overline{z}$ diventa $C(\overline{z})=0$ e quindi la linearizzazione in $0$ è data da 
\[
\dot{y}=Ay
\]
infatti \[
Y'(y)=\left( P^{-1} X' (\overline{z}+Py) \right) Py = P^{-1} X'_{|z} P y \overset{z=\overline{z}}{\longrightarrow} A y
\]

Mostrare che $0$ è asintoticamente stabile per $\dot{y}=Y(y)$ equivarrà a mostrare che $\overline{z}$ lo sia per $\dot{z}=X(z)$.

Sia $\mathcal{W}=\dfrac{1}{2}\left|\left|y\right|\right|^2$ vediamo che $\mathcal{W}$ avrà un minimo assoluto nell'origine e che 
\[
L_Y\mathcal{W}(y)=Y(y)\cdot\nabla\mathcal{W}(y)=Y(y)\cdot y
\]
Dato che $Y'(0)=0$ e $Y'(y)\overset{y\rightarrow 0}{\rightarrow} Ay$ possiamo usare la formula di Taylor al primo ordine attorno all'origine ottenendo
\[
Y(y)=Ay+r(y)
\]
dove il resto $r(y)$ è un campo vettoriale su $\Omega$ che soddisferà $r(0)=0$ e $r'(0)=0$.
Quindi
\[
L_y\mathcal{W}=y\cdot \left( Ay+r(y) \right)
\]
Possiamo vedere $A$ come 
\[
A=\text{diag}(\alpha_1,...,\alpha_k,\alpha_{2k+1},...,\alpha_n) + Q = D+Q
\] 
con $Q$ antisimmetrica
\[
Q=\left(\begin{array}{cccccc|ccccc}
0 & \beta_1  & 0 & \dots &   0 &0 &0 & &\dots & &  0 \\ 
-\beta_1 & 0 & \beta_2  & 0 & \dots &0  &0 &  &\dots & &  0 \\  
0 & -\beta_2 &  \ddots  & & & & 0 & & \dots & & 0 \\
\vdots &  & & \ddots &  & & \vdots & & & & \vdots \\ 
 &  & & & 0 &\beta_k & 0 & & \dots &  & 0 \\ 
 &  & & & -\beta_k & 0 & 0 & &\dots & &  0 \\
 \hline
 &  & & &  &  &  & &  & & \\
 &  & \textbf{0} & &  &  &  & & \textbf{0} & &  \\
 &  & & &  & & & & & & \\
\end{array}\right)
\]
quindi la derivata di Lie sarà
\[
L_y\mathcal{W}=y\cdot \left(Dy + Qy +r(y) \right) =
y\cdot \left( Dy +r(y) \right) = y Dy + r(y)^T y
\]
Usando il \textit{secondo teorema di Lyapunov} se la funzione $yDy+r(y)^Ty$ è negativa in un intorno bucato di $0$, ovvero se ha massimo stretto nell'origine, allora avremo provato la stabilità asintotica.

Gradiente ed Hessiana saranno
\begin{gather*}
\nabla L_Y\mathcal{W}(y)=Dy+yD+r(y)^T+r'(y)^Ty=2Dy+r(y)^T+r'(y)^Ty \\
\left( L_Y\mathcal{W} \right)''(y) = 2D + 2r'(y)^T + r''(y)^Ty
\end{gather*}
I termini di $Dy$ e $r''(y)^Ty$ si annullano in $y=0$, $r(y)=o(y)$ per definizione e quindi
\[
\nabla L_Y\mathcal{W}(0)= 0 \hspace{0.15\linewidth} \left(L_Y\mathcal{W} \right)''(0)=2D
\]
ed essendo $D$ definita negativa l'origine è un massimo stretto. \qedhere
\end{itemize}
\end{dimo}

Il metodo spettrale ci dice che per alcuni equilibri le proprietà di stabilità dell'equazione differenziale attorno all'equilibrio sono le stesse della linearizzazione, cioè che i termini non lineari dell'equazione sono ininfluenti.
Se però le ipotesi del teorema non sono verificate questo non è detto: per esempio se gli autovalori hanno tutti parte reale negativa eccetto uno o più che sono completamente immaginari l'equilibrio sarà stabile per la linearizzazione ma non potremo sapere se lo sarà anche per l'equazione non lineare.

Nell'importante caso dei sistemi meccanici conservativi, in cui la stabilità asintotica è impossibile per la proposizione \ref{propIntegraliPrimiAttattività}, il metodo spettrale fornisce ``solo'' una condizione sufficiente per l'instabilità e quindi una condizione necessaria per la stabilità (stabile $\implies$ non instabile).

\section{Equazioni differenziali su sottovarietà}
Le equazioni differenziali che abbiamo visto fin'ora erano su aperti di $\mathbb{R}^n$, sarà però necessario estendere la teoria ad equazioni differenziali su sottovarietà di $\mathbb{R}^n$. 

In realtà abbiamo già visto che se l'equazione ha un integrale primo la dinamica si svolge sui suoi insiemi di livello che, sotto alcune ipotesi, sono sottovarietà dello spazio delle fasi. La restrizione ad essi della dinamica sarà descritta dall'equazione ristretta che è proprio un'equazione su una sottovarietà.

\subsection{Sottovarietà invarianti}\label{sezSottovarietàInvarianti}
Mostrare che un generico sottoinsieme dello spazio delle fasi sia un invariante può essere complesso mentre è più semplice per una sottovarietà.

Sia $M$ una sottovarità di $\mathbb{R}^n$ di dimensione $m$, i vettori tangenti ad $M$ in un suo punto $z$ saranno, come già visto, le derivate di curve differenziabili a valori in $M$ passanti per $z$, il loro insieme è $T_zM$ è un sottospazio vettoriale dello spazio $T_z\mathbb{R}^n$, a sua volta formato da tutti i vettori tangenti a $\mathbb{R}^n$, ed avrà dimensione $m$.

\begin{definizione}
Sia $M$ una sottovarietà di $\mathbb{R}^n$, un campo vettoriale $X$ su un aperto $\Omega\subseteq\mathbb{R}^n$ si dirà tangente ad $M$ se
\[
X(z)\in T_zM \quad\forall z\in M
\]
\end{definizione}

Ovviamente non ha senso controllare la tangenza nei punti al di fuori di $M$ e quindi l'essere tangente si può vedere come una proprietà della restrizione $X|_{M}$ di $X$.

\begin{proposizione}
Sia $M$ una sottovarietà di $\mathbb{R}^n$ e sia $X$ un campo vettoriale definito in un aperto $\Omega\subseteq\mathbb{R}^n$ contente $M$.
Se $M$ è invariante per $\dot{z}=X(z)$ allora $X$ è tangente a $M$. 
\end{proposizione}
\begin{dimo}
Se $M$ è invariante allora per ogni $z_0\in M$ la curva integrale $t\mapsto z_t$ con dato iniziale $z_0$ è una curva differenziabile a valori in $M$ e di conseguenza
\[
\dfrac{d}{dt}z_t|_{t=0}=\dot{z_0}=X(z_0)\in T_{z_0}M \qedhere
\]
\end{dimo}

Il viceversa, $X$ tangente a $M$ ``$\implies$'' $M$ invariante per $X$, è più difficile da formulare.

\begin{proposizione}
Sia $M$ un sottovarietà di $\mathbb{R}^n$ e sia $X$ un campo vettoriale completo definito in un aperto $\Omega\subseteq\mathbb{R}^n$ contenente $M$. Allora 
\begin{enumerate}[i.]
\item Se $X$ è tangente ad $M$ una curva integrale di $X$ che passi per un punto di $M$ allora se esce da $M$ lo fa passando per un punto di $\overline{M}\setminus M$.
\item Se $X$ è tangente ad $M$ e $M$ è topologicamente chiuso allora $M$ è invariante sotto il flusso di $X$.
\end{enumerate}
\end{proposizione}

\begin{dimo}
Consideriamo un sistema di coordinate locali di $\mathbb{R}^n$ definite in un intorno $V$ di $z_0$ ed adattate ad $M$, cioè un diffeomorfismo
\[
C:V\rightarrow Y\times U \subseteq \mathbb{R}^m\times\mathbb{R}^{n-m} \hspace{0.15\linewidth} C(z)=(y,u)
\]
con $Y\subset\mathbb{R}^m$ e $U\subseteq\mathbb{R}^{n-m}$ aperti, per cui $M\cap V$ è un $C(M\cap V)=\{(y,u):u =\text{costante}\}=Y\times \text{cost}$.

L'esistenza di queste coordinate è garantita dalla definizione di sottovarietà.\footnotemark
In queste coordinate $\dot{z}=X(z)$ diventerà
\[
\dot{y}=\tilde{X}_y(y,u) \hspace{0.15\linewidth} \dot{u}=\tilde{X}_u(y,u)
\]
dove la $\tilde{X}$ è $C_*X$. 
I punti di $M\cap V$ sono caratterizzati da un valore di $u$, per esempio $u=0$, ed essendo $X$ tangente ad $M$ nelle nuove coordinate $\tilde{X}$ sarà tangente a $C(M\cap V)=\{u=0\}$ e quindi
\[
\tilde{X}_u(y,0)=0 \quad \forall y\in Y
\]
e questo implica che ogni soluzione di $\left(\dot{y},\dot{u}\right)=\tilde{X}(y,u)$ con punto iniziale avente $u_0=0$ avrà $u_t=0$ per ogni tempo nel suo intervallo di esistenza, la corrispondente curva integrale di $X$ potrebbe uscire dall'intorno $V$ che abbiamo usato per definire le coordinate locali e quindi, anche se il campo è completo, non avremo direttamente invarianza per tutti i tempi.\par 

Per ogni $z_0\in Y$ la $t\mapsto y(t,y_0)$ è soluzione di $\dot{y}=\tilde{X}_y(y,0)$. Dal teorema di esistenza ed unicità la soluzione sarà definita in un intorno $J$ di $0$. 
La curva
\[
t\mapsto (y_t,0) \in Y\times U \hspace{0.15\linewidth} t\in J
\]
è la soluzione di $\left(\dot{y},\dot{u}\right)=\tilde{X}(y,u)$ con dato iniziale $(z_0,0)$.
La sua preimmagine $t\mapsto C^{-1}(y_t,0)$ apparterrà quindi ad $M$ per tutti i $t\in J$.

Quindi per $t\in J$, $\forall y_t \in Y$ si ha $\Phi_t^X\left(C^{-1}(y_t,0)\right)=\left(C^{-1}(y_t,0)\right)\in M$
\begin{enumerate}[i.]
\item Supponiamo che la $t\mapsto\Phi_{t}^X(z_0)=z_t$ esca da $M$ per la prima volta ad un tempo $\overline{t}>0$. Questo significa 
\begin{itemize}
\item $z_t\in M \quad \forall t\in [0,\overline{t}\, )$
\item In ogni intorno destro di $\overline{t}$ vi è un $\tilde{t}$ tale che $z_{\tilde{t}} \not\in M$
\end{itemize}
e ovviamente $\overline{z}=z_{\overline{t}}\in\overline{M}$.

Sia per assurdo $\overline{z}\in M$. Allora $\exists b>0$ tale che $t\mapsto\Phi_{t}^X(\overline{z})\in M \quad \forall t\in(-b,b)$ per quanto visto sopra.
Però 
\[\Phi_t^X(\overline{z})=\Phi_t^X \left( \Phi_{\overline{t}}^X(z_0) \right)=\Phi_{t+\overline{t}}^X(z_0)=z_{t+\overline{t}} \]
che implicherebbe che $z_t\in M \quad \forall t\in (\overline{t},\overline{t}+b)$ che è assurdo dato che la $t\mapsto z_t$ dovrebbe uscire da $M$ a $\overline{t}$.

Quindi i punti di uscita delle soluzioni sono in $\overline{M}\setminus M$.
\item Se $M$ è chiusa $\overline{M}=M$ e quindi le curve per uscire dovrebbero passare da $M\setminus M=\emptyset$ \qedhere
\end{enumerate}
\end{dimo}
\footnotetext{Si veda la definizione della descrizione cartesiana di una varietà.}



%Parte 2-----------------------------------------------------------------
\cleardoublepage
\addcontentsline{toc}{part}{Meccanica}
\thispagestyle{plain}
	\begin{center}
		\vspace{.1\textheight}
		{\LARGE\scshape \textsc{Seconda parte}\par}
		\vspace{.1\textheight}
		\textsc{}\par
		\vspace{.05\textheight}
		{\Huge \textsc{Meccanica}\par}
		\vfill

\end{center}

%------------------------------------------------------------------


\fancyfoot[C]{\small Meccanica}
\section{Modello}


\newcommand{\proSc}{\odot}
\newcommand{\proVe}{\otimes}

Per poter descrivere la meccanica avremo bisogno di definire alcuni concetti ``primitivi'': \begin{itemize}
\item Il \textbf{tempo}, che assumiamo assoluto nel senso che l'intervallo di tempo fra due eventi sarà uguale per ogni osservatore, unidimensionale e con struttura affine.
\item Lo \textbf{spazio} in cui avvengono i fenomeni $\mathbb{E}^3$, che assumiamo tridimensionale e con struttura affine euclidea.
\item I \textbf{punti materiali} o particelle, oggetti che si muovono in $\mathbb{E}^3$ occupandone istantaneamente un solo punto allo scorrere del tempo, dotati di \textbf{massa}.
\item Le \textbf{forze} agenti sui punti materiali.
\end{itemize}

Lo spazio $\mathbb{E}^3$ è uno spazio affine reale, cioè ad ogni coppia di punti $P_1,P_2\in\mathbb{E}^3$ sarà associato univocamente un vettore di uno spazio vettoriale tridimensionale, che chiameremo $\mathbb{V}^3$, che li ``congiunge''.
Denoteremo $P_1P_2$ o $P_2-P_1$ il vettore in $\mathbb{V}^3$ tra $P_1$ e $P_2$.

Se fissiamo un punto $O\in\mathbb{E}^3$ possiamo definire una biiezione tra lo spazio affine e lo spazio vettoriale
\[
\mathbb{E}^3\rightarrow\mathbb{V}^3 \hspace{0.15\linewidth} P\mapsto OP\in\mathbb{V}^3
\]
Essendo euclideo $\mathbb{V}^3$ è dotato di un prodotto scalare (che chiameremo $\proSc$) con cui possiamo definire una norma $\left|\left| \> \right|\right|_{\mathbb{V}^3}$ in $\mathbb{V}^3$, un prodotto vettore $\proVe$ e distanze ed angoli in $\mathbb{E}^3$.


Le forze potranno essere \textbf{interne}, se sono dovute ad interazioni tra punti del sistema, o \textbf{esterne}. Supporremo che valga il \textbf{principio di sovrapposizione} cioè che forze dovute a cause diverse siano indipendenti fra loro.

Un \textbf{moto} è una curva differenziabile in $\mathbb{E}^3$ $t\mapsto P(t)$.


\subsection{Sistemi di riferimento}

Un \textbf{riferimento spaziale} $\Sigma$ in $\mathbb{E}^3$ è dato dalla scelta di un punto $O\in\mathbb{E}^3$ e di una base $\{E_1,E_2,E_3\}$ ortonormale di $\mathbb{V}^3$
\[
\Sigma = \left\{ 0 ; E_1,E_2,E_3 \right\}
\]
chiameremo $O$ l'origine di $\Sigma$ e i vettori della base di $\mathbb{V}^3$ gli assi.

D'ora in avanti useremo sempre implicitamente basi ortonormali \textbf{levogire} cioè tali per cui
\begin{equation}
E_1\proVe E_2 =E_3  \quad E_2\proVe E_3 = E_1 \quad E_3\proVe E_1=E_2
\end{equation}

La scelta di una base di $\mathbb{V}^3$ definisce un isomorfismo tra $\mathbb{V}^3$ e $\mathbb{R}^3$ infatti
\[
U\in\mathbb{V}^3 \implies U=\sum_i^3 u_i E_1 \quad u_i\in\mathbb{R}^3
\]
e quindi possiamo associare ad ogni vettore il suo \textbf{rappresentativo} $U\mapsto u=(u_1,u_2,u_3)\in\mathbb{R}^3$.

Per un riferimento $\Sigma$ chiameremo \textbf{coordinate cartesiane} $x$ di un punto $P\in\mathbb{E}^3$ o anche \textbf{vettore posizione} di P in $\Sigma$ il rappresentativo del vettore $OP$
\[
OP=\sum_i^3 x_iE_i \hspace{0.15\linewidth} x=(x_1,x_2,x_3)
\]

Considereremo anche sistemi di riferimento dipendenti dal tempo $\Sigma_t$, cioè funzioni differenzibili $\mathbb{R}\rightarrow\mathbb{E}^3\times\mathbb{V}^3\times\mathbb{V}^3\times\mathbb{V}^3$ 
\[t\mapsto \Sigma_t=\{ O(t); E_1(t),E_2(t),E_3(t) \}
\]
la scelta di un $\Sigma_t$ fornisce unisce un'identificazione di $\mathbb{E}^3$ con $\mathbb{R}^3$ dipendente dal tempo, cioè ad ogni punto $P\in\mathbb{E}^3$ si associa come vettore posizione una terna di funzioni differenziabili
\[
O(t)P=\sum_i^3 x_i(t)E_i(t)
\]

Un moto $P(t)$ in coordinate cartesiane sarà quindi in generale una curva $t\mapsto x(t)\in\mathbb{R}^3$ definita da
\begin{equation}
O(t)P(t)=\sum_i^3 x_i(t) E_i(t)
\end{equation}

Scegliendo un sistema di riferimento potremo valutare le variazioni temporali di funzioni vettoriali del tempo
\begin{definizione}[Derivata temporale relativa a $\Sigma$]
Per $U:\mathbb{R}\rightarrow\mathbb{V}^3$ differenziabile ne definiamo la derivata temporale in $\Sigma$, per $\Sigma$ riferimento di $\mathbb{E}^3$, la funzione vettoriale
\[
\left. \dfrac{dU}{dt} \right|_{\Sigma}(t) := \sum_i^3 \dot{u}_i E_i
\]
dove le $\dot{u}_i(t)$ sono le derivate delle $t\mapsto u(t)=U(t)\proSc E_i(t)$.
\end{definizione}

\begin{definizione}
Per il moto di un punto materiale $P(t)$ definiremo velocità ed accelerazione 
\[
V_P^\Sigma(t)=\sum_u^3 \dot{x}_i(t)E_i(t) \hspace{0.15\linewidth} A_P^\Sigma(t)=\sum_i^3 \ddot{x}_i(t)E_i(t)
\]
\end{definizione}

\subsection{Forze}

Supponiamo che la forza che agisce su un punto materiale $P_\alpha$, $\alpha\in[1,N]$, in un sistema di riferimento $\Sigma$ sia una funzione differenziabile
\[
F_\alpha^\Sigma : \left(\mathbb{V}^{3}\right)^N \times \left(\mathbb{V}^{3}\right)^N \times \mathbb{R} \rightarrow \mathbb{V}^3
\]
con dominio lo spazio degli atti di moto (eventualmente $\times\mathbb{R}$ se dipende anche dal tempo) e valori nei vettori tridimensionali.

Le forze che agiscono sul sistema meccanico si divideranno in \textbf{forze interne}, causate da interazioni tra le particelle, e \textbf{forze esterne}.
Dal principio di sovrapposizione
\[
F_\alpha^\Sigma=F_\alpha^{\Sigma,int}+F_\alpha^{\Sigma,ext}
\]

Tipicamente le forze interne sono interazioni posizionali tra le particelle $F_\alpha^{\Sigma,int}=\sum^N F_{\beta\alpha}(OP_\alpha,OP_\beta)$ mentre le forze esterne sono funzioni di posizione e velocità della sola particella $F_\alpha^{\Sigma,ext}(OP_\alpha,V_\alpha^\Sigma,t)$.

\clearpage
Spesso le forze esterne sono dovute ad approssimazioni nella definizione del sistema, cioè al non considerare alcuni dei punti materiali presenti come parte del sistema, nell'ambito della meccanica classica però anche la forza magnetica per esempio dovrà considerarsi una forza esterna.
\begin{esempio}[Forza peso e problema ristretto dei 3 corpi]
Come già sappiamo dai corsi di fisica precedenti la forza peso di un corpo vicino alla superficie della terra è in realtà una forza di attrazione gravitazionale reciproca tra il corpo e la terra, dato che però l'effetto del corpo sulla terra è generalmente trascurabile la si considera una forza esterna costante ``vicino'' alla superficie terreste.

Consideriamo come sistema meccanico due corpi celesti $M_1,M_2$ ed una sonda $m$: le forze di interazione tra i corpi saranno solo interne e corrisponderanno alle attrazioni gravitazionali. 
Essendo la forza gravitazionale proporzionale alla massa l'effetto della sonda sui due corpi celesti sarà però trascurabile e potremo studiare il moto della sonda in modo molto più semplice approssimando e considerando il sistema della sola sonda su cui agiscono due forze di attrazione esterne.
\end{esempio}

\subsection{Sistemi meccanici}\label{sezSistemiMeccanici}
Un \textbf{sistema meccanico} è costituito da $N\geq1$ punti materiali $P_1,...,P_N$ da un sistema di riferimento e dalle forze che agiscono, per quel sistema di riferimento, sui punti materiali.


Per $\alpha=1,...,N$ le equazioni del moto di una particella $P_\alpha$ nelle coordinate affini di $\Sigma$ sono
\begin{equation}\label{eqMotoPartSistemaCarte}
m_\alpha \ddot{x}_\alpha (t) = f(x_1(t),...,x_N(t),\dot{x}_1(t),...,\dot{x}_N(t),t)
\end{equation}


Chiameremo \textbf{configurazione} del sistema il vettore
\[
X=(x_1,...,x_N)\in\mathbb{R}^{3N}
\]
\textbf{velocità del sistema} il vettore
\[
\dot{X}=(\dot{x}_1,...,\dot{x}_N)\in\mathbb{R}^{3N}
\]
e \textbf{atto di moto} del sistema il vettore
\[
(X,\dot{X})\in\mathbb{R}^{6N}
\]

Per scrivere in modo più compatto le equazioni del moto useremo anche il vettore \textbf{forza sul sistema}
\[
F=(f_1,...,f_N)\in\mathbb{R}^{3N}
\]
e la \textbf{matrice delle masse} 
\[
M=\left(\begin{array}{ccccccc}
m_1 \\
& m_1 \\
& & m_1 \\
& & & \ddots \\
& & & & m_N \\
& & & & & m_N \\
& & & & & & m_N
\end{array}\right)=\text{diag}(m_1,m_1,m_1,...,m_N,m_N,m_N)
\]
e quindi il sistema delle \eqref{eqMotoPartSistemaCarte} si riscrive come
\[
M\ddot{X}=F(X,\dot{X},t)
\]
ed assumendo che le masse delle particelle del sistema siano positive la matrice delle masse è invertibile e quindi
\begin{equation}
\ddot{X}=M^{-1}F(X,\dot{X},t)
\end{equation}
o equivalentemente col sistema di equazioni differenziali del primo ordine
\begin{equation}
\begin{cases}
\dot{X}=V \\
\dot{V}=M^{-1}F(X,\dot{X},t)
\end{cases}
\end{equation}


\begin{definizione}
Per un sistema meccanico libero l'energia cinetica è una funzione $T:\mathbb{R}^{3N}\times\mathbb{R}^{3N}\rightarrow\mathbb{R}$ che associa ad ogni atto di moto $(X,\dot{X})$ la somma delle energie cinetiche dei punti materiali del sistema in quell'atto di moto
\[
T(X,\dot{X})=\sum_{\alpha=1}^N \dfrac{1}{2} m_{\alpha} \left|\left| \ddot{x}_\alpha \right|\right|^2 = \dfrac{1}{2}  \dot{X} \cdot M\dot{X}
\]
\end{definizione}

\section{Relazioni tra sistemi di riferimento diversi}

Dato un sistema meccanico nel sistema $\Sigma_t=\{O(t);E_1(t),E_2(t),E_3(t)\}$ vogliamo tradurre le forze e i moti in un secondo sistema $\Sigma^*_t=\{O^*(t);E_1^*(t),E_2^*(t),E_3^*(t)\}$ che si muove rispetto a $\Sigma_t$ con un moto noto.

Gli assi formano basi di uno stesso spazio vettoriale quindi sappiamo che il cambiamento di base che porta vettori da $\{E_i^*\}_i$ a vettori in $\{E_i\}$ è
\begin{equation}
Rij(t)=E_i(t)\proSc E_j^*(t)
\end{equation}
La matrice $R$ è ortogonale per costruzione (le sue colonne sono i rappresentativi dei $E_i$ nella base $\{E_i^*\}_i$) ed è una rotazione.
I rappresentativi di uno stesso vettore nelle due basi sono legati da
\begin{equation}
u=Ru^* \hspace{0.15\linewidth} u^*=R^{-1}u=R^Tu
\end{equation}

Essendo la matrice di rotazione ortogonale possiamo ricavare
\begin{align*}
RR^T=\mathbb{I}_3 \implies \dot{R} R^T +R\dot{R}^T = \dot{R}R^T+\left(\dot{R}R^T\right)^T=0 
\end{align*}
e quindi la matrice reale $\dot{R}R^T$ è antisimmetrica. Una matrice $3\times 3$ antisimmetrica ha $3$ parametri indipendenti e quindi esisterà un vettore $\omega(t)\in\mathbb{R}^3$ con componenti i tre parametri della matrice.

\begin{definizione}[Velocità angolare]
Definiamo velocità angolare ale tempo $t$ di $\Sigma^*$ rispetto a $\Sigma$ il vettore $\Omega(t)\in\mathbb{V}^3$ con rappresentativo avente per componenti i tre parametri della matrice $\hat{\omega}(t)$ 
\[
\hat{\omega}(t)=\left(\begin{array}{ccc}
0 & -\omega_{3} & \omega_{2} \\
\omega_{3} & 0 & -\omega_{1} \\
-\omega_{2} & \omega_{1} & 0
\end{array}\right) 
\hspace{0.15\linewidth}
\omega(t)=\left(\begin{array}{c}
\omega_1 \\ \omega_2 \\ \omega_3 
\end{array}\right) \hspace{0.15\linewidth}
\Omega(t)=\sum_i^3 \omega_i E_i
\]
e velocità angolare di $\Sigma^*$ rispetto a $\Sigma$ la funzione $t\mapsto \Omega(t)$.

La velocità angolare è tale che
\[
\hat{\omega}u=\omega \times u \quad \forall u\in\mathbb{R}^3 \text{ che sarà il rappresentativo di } \Omega\proVe U \in\mathbb{V}^3 \quad \forall U\in\mathbb{V}^3
\]
\end{definizione}

La velocità angolare si potrà definire in questo modo come vettore soltanto in spazi tridimensionali: solo in $M_{3\times 3}$ le matrici antisimmetriche hanno proprio $3$ parametri indipendenti.

\begin{proposizione}[Formule di Poisson]
Dati due sistemi di riferimento $\Sigma_t$ e $\Sigma^*_t$per ogni $t$ esiste uno ed un solo vettore $\Omega(t)\in\mathbb{V}^3$ tale che
\[
\left. \dfrac{dE_i^*}{dt} \right|_{\Sigma} (t)=\Omega\proVe E_i^*(t)
\]
\end{proposizione}
\begin{dimo}
Per quanto detto sulla velocità angolare i rappresentativi degli assi dei due riferimenti sono legati da\footnotemark[2]
\[
e'_i(t)=R(t)e^*_i \hspace{0.15\linewidth} e^*_i=R(t)^Te'_i(t)
\]
dove chiamiamo $e'_i(t)$ le componenti del rappresentativo dei vettori $E^*_i$ nella base $\{E_i\}_i$.
Derivando rispetto al tempo avremo $\dot{e}'_i(t)=\dot{R}(t)e^*_i=\dot{R}(t)R^T(t)e'_i(t)$
Per quanto visto sulla velocità angolare
\[
\dot{e}'_i(t)=\hat{\omega}(t) e'_i(t)=\omega(t)\times e'_i(t)
\]
e quindi
\[
\sum_i^3 \dot{e}'_i(t) E_i = \Omega(t)\times\sum_i^3  e'_i(t) E_i \implies \left. \dfrac{dE_i^*}{dt} \right|_{\Sigma} (t)=\Omega\proVe E_i^*(t)   \qedhere
\]

\end{dimo}
\footnotetext[2]{Usiamo che chiaramente i rappresentativi dei vettori che compongono una base rispetto a quella stessa base sono costanti $(1,0,0)..$.}
Questo ci permette di dedurre le relazioni fra le derivate temporali della stessa funzione nei due sistemi di riferimento diversi

\begin{proposizione}
Per una qualunque funzione vettoriale $t\mapsto U(t)\in\mathbb{V}^3$ in ogni istante $t$ vale
\begin{equation}
\left. \dfrac{d U}{dt} \right|_\Sigma (t) = \left. \dfrac{dU}{dt} \right|_{\Sigma^*}(t) + \Omega(t) \proVe U(t)
\end{equation}
dove $\Omega(t)$ è la velocità angolare di $\Sigma^*$ rispetto a $\Sigma$.
\end{proposizione}
\begin{dimo}
Derivando in $\Sigma$ la $U(t)=\sum u_i^*(t)E_i^*(t)$ troviamo\footnotemark{}
\begin{align*}
\left. \dfrac{dU}{dt}\right|_{\Sigma} &= \left. \dfrac{d}{dt} u^*_i(t) E_i^*(t) \right|_{\Sigma} =\\
& = \sum_i^3 \dot{u}^*_i(t) E_{i}^*(t) + \sum_i^3 u^*_i(t) \Omega\proVe E_i^*(t) = \\
& = \left. \dfrac{dU}{dt} \right|_{\Sigma^*}(t) + \Omega(t) \proVe U(t) \qedhere
\end{align*}
\end{dimo}


Quindi in particolare avremo $\left.\dfrac{d\Omega}{dt}\right|_\Sigma (t) =\left. \dfrac{d\Omega}{dt}\right|_{\Sigma^*} (t)+ \Omega (t) \proVe \Omega(t) \implies \left.\dfrac{d\Omega}{dt}\right|_\Sigma (t)=\left. \dfrac{d\Omega}{dt}\right|_{\Sigma^*}(t) \quad \forall t$


\begin{proposizione}[Moti nei diversi sistemi di riferimento]
Ad ogni istante $t$ le velocità ed accelerazione di un punto $P$ relative a due sistemi di riferimento $\Sigma_t$, $\Sigma^*_t$ sono legate da
\begin{gather}
V_P^\Sigma(t) = V_P^{\Sigma^*} (t) + V^{Tr}(t) \\
A_P^{\Sigma}(t)= A_P^{\Sigma^*}(t) + A^{Tr}(t) + A^{Cor}(t)
\end{gather}
con \begin{itemize} 
\item $V^{Tr}(t)=V_{O^*}^{\Sigma}(t) + \Omega(t)\proVe O^*P(t)$ \textbf{velocità di trascinamento}.
\item $A^{Tr}(t)=A_{O^*}^{\Sigma}(t)+ \Omega(t)\proVe \left[ \Omega(t) \proVe O^*P(t) \right] + \dot{\Omega}(t)\proVe O^*P(t)$ 
\item $A^{Cor}(t)= 2 \Omega(t) \proVe V_P^{\Sigma^*}(t)$ \textbf{accelerazione di Coriolis}.
\end{itemize}
Dove $\Omega(t)$ è la velocità angolare di $\Sigma^*$ rispetto a $\Sigma$.
\end{proposizione}
\begin{dimo}
Per avere una notazione più snella eviteremo di indicare le dipendenze temporali e scriveremo $V$ per $V_P^\Sigma$, $V^*$ per $V_P^{\Sigma^*}$, $A$ per $A_P^\Sigma$ e $A^*$ per $A_P^{\Sigma^*}$.\par 
Useremo la scomposizione $OP=O^*P+OO*$ derivandola in $\Sigma$: 
\begin{align*}
V & =\left. \dfrac{d O^*P}{dt}\right|_{\Sigma} + V_{O*}^\Sigma = \\ &= \left. \dfrac{d O^*P}{dt}\right|_{\Sigma^*} + \Omega\proVe O^*P + V_{O*}^\Sigma =\\ &= V^* + V_{O^*}^\Sigma + \Omega\proVe O^*P
\end{align*}
Derivando ancora in $\Sigma$ entrambi i membri otteniamo
\begin{align*}
A &=\left. \dfrac{d V^*}{dt}\right|_\Sigma + A_{O^*}^\Sigma + \left. \dfrac{d \Omega \proVe O^*P}{dt}\right|_{\Sigma}=\\
&= A^* + \Omega \proVe V^* + A_{O^*}^\Sigma + \left. \dfrac{d \Omega \proVe O^*P}{dt}\right|_{\Sigma^*} + \Omega\proVe\left[  \Omega \proVe O^*P\right]
\end{align*}
in cui vediamo che
\begin{align*}
\left. \dfrac{d \Omega \proVe O^*P}{dt}\right|_{\Sigma^*} &= \left. \dfrac{d \Omega }{dt}\right|_{\Sigma^*}\proVe O^*P + \Omega \proVe \left. \dfrac{d  O^*P}{dt}\right|_{\Sigma^*} =\dot{\Omega}\proVe O^*P + \Omega \proVe V^* \qedhere
\end{align*}
\end{dimo}

Chiameremo \begin{itemize}
\item \textbf{Forza di trascinamento}
\[
F_P^{Tr}(O^*P,t):=-m\left[ A_{O^*}^{\Sigma}(t)+ \Omega(t)\proVe \left[ \Omega(t) \proVe O^*P(t) \right] + \dot{\Omega}(t)\proVe O^*P(t) \right]
\]
\item \textbf{Forza di Coriolis} 
\[
F_P^{Cor}(V_P^{\Sigma^*},t):=-2m \Omega(t) \proVe V_P^{\Sigma^*}(t)
\]
\end{itemize}
queste due forza sono esterne.

\begin{proposizione}
Ad ogni istante $t$ le forze agenti su un punto materiale in moto $P$ in due sistemi di riferimento $\Sigma$ e $\Sigma^*$ sono legate da
\begin{align*}
F_P^{\Sigma^*}(O^*P,V_P^{\Sigma^*},t)= F_P^{\Sigma} (OO^*+O^*P, V_P^\Sigma+V_{O^*}^{\Sigma}+\Omega\proVe O^*P, t) + F_P^{Tr}(O^*P,t) + F_P^{Cor}(V_P^{\Sigma^*},t)
\end{align*}
\end{proposizione}

Questa dimostrazione è immediata sapendo che $mA^\Sigma=F^\Sigma$, il suo interesse (ed il motivo per la scrittura un po' pesante del suo enunciato) è nell'avere le corrette dipendenze funzionali.

Se consideriamo una trasformazione $G:\mathbb{V}^3\times\mathbb{V}^3\times\mathbb{R}\rightarrow\mathbb{V}^3\times\mathbb{V}^3\times\mathbb{R}$ che manda la posizione e velocità di un punto in un sistema di riferimento in quelle in un altro in funzione delle quantità misurate nel primo riferimento:
\begin{equation}
G(O^*P,V_P^{\Sigma^*},t)=(O^*O+O^*P, V_P^{\Sigma^*}+V_{O^*}^\Sigma+\Omega\proVe O^*P, t)
\end{equation}
potremo scrivere in modo più compatto
\[
F_P^{\Sigma^*}(O^*P,V_P^{\Sigma^*},t)=F_P^\Sigma \circ G (O^*P,V_P^{\Sigma^*},t) + F_P^{Tr}(O^*P,t) + F_P^{Cor} (V_P^{\Sigma^*},t)
\]
la trasformazione $G$ è chiamata \textbf{trasformazione di Galilei}.




\section{Punto vincolato ad una curva}
Consideriamo un punto materiale che sia vincolato a rimanere su una curva $\Gamma$
\[
x_t=\gamma(s_t) \quad \forall t \hspace{0.15\linewidth} s\mapsto\gamma(s)\in\Gamma\subset\mathbb{R}^3
\]
Perchè questa situazione possa essere compatibile con la legge di Newton servirà che la curva o qualche altro dispositivo eserciti una forza $\phi$ sul punto in modo tale che, per ogni scelta della forza $f$ agente sul punto, le traiettorie rimangano sulla curva
\[
m\ddot{x}=f(x,\dot{x})+\phi(x,\dot{x})
\] 
Chiameremo il campo di forza $\phi$ \textbf{reazione vincolare} e l'ordinaria forza agente sul punto $f$ \textbf{forza attiva}.

Scegliamo la parametrizzazione $s\mapsto\gamma(s)$ tale che $\gamma'(s)=:e_\tau (s)$ abbia norma unitaria in ogni punto, questa parametrizzazione esiste sempre e il suo parametro $s$ si dirà \textbf{parametro d'arco}.
In ciascun punto $\gamma(s)\in\Gamma$ consideriamo la base di $T_{\gamma(s)}\mathbb{R}^3$ costituita da 
\begin{equation}
\left\{ e_\tau(s),e_n(s),e_b(s) \right\} \qquad e_n(s):=\rho(s)e'_\tau(s) \quad e_b(s):=e_\tau(s)\times e_n(s) \quad \rho(s):=\left|\left| e_\tau(s)\right|\right|^{-1}
\end{equation}
questa base è chiamata \textbf{base di Frenet}, $e_n$ è detto versore \textbf{normale principale} , $e_b$ versore \textbf{binormale} e $\rho$ \textbf{raggio di curvatura principale}.

Considerando un moto $t\mapsto x(t)=\gamma(s_t)\in\Gamma$:
\begin{gather*}
\dot{x}_t=\dfrac{d\gamma}{dt}(s_t) \dot{s}_t = \dot{s}_t e_\tau(s_t)\\
\ddot{x}_t=\ddot{s}_t e_\tau(s_t) + \dot{s} e'_\tau (s_t) \dot{s}_t = \ddot{s}_t e_\tau(s_t) + \dfrac{\dot{s}^2_t}{\rho(s_t)}e_n(s_t)
\end{gather*}

Ora dividiamo in componenti nella base di Frenet la forza attiva e valutiamo componente per componente la $m\ddot{x}=f+\phi$
\begin{align}
&e_\tau (s_t): \qquad m\ddot{s}_t = f_\tau(s,\dot{s}) + \phi(s,\dot{s}) \label{eqPuntoCurva1}\\
&e_n(s_t): \qquad m \dfrac{\dot{s}^2_t}{\rho(s_t)} = f_n (s,\dot{s}) + \phi(s,\dot{s}) \label{eqPuntoCurva2}\\
&e_b(s_t): \qquad 0 = f_b(s,\dot{s}) + \phi_b(s,\dot{s}) \label{eqPuntoCurva3}
\end{align}
Supponendo che la $f$ sia fissata:\begin{enumerate}[i.]
\item Per ogni scelta della componente $\phi_\tau$ la \eqref{eqPuntoCurva1} è un'equazione differenziale del secondo ordine in $s\in\mathbb{R}$ che si può mettere in forma normale: per ogni $(s_0,\dot{s}_0)$ esiste una ed una sola soluzione $t\mapsto s_t$.
\item La corrispondente curva $t\mapsto\gamma(s_t)\in\Gamma$ è una soluzione valida del sistema solo se $t\mapsto s_t$ soddisfa anche \eqref{eqPuntoCurva2} e \eqref{eqPuntoCurva3}. 
Assumendo che le altre due componenti della reazione vincolare siano date da
\begin{equation}
\phi_n(s,\dot{s})=m\dfrac{\dot{s}^2}{\rho(s)}-f_n(s,\dot{s}) \hspace{0.15\linewidth} \phi_b(s,\dot{s})=-f_b(s,\dot{s})
\end{equation}
la $t\mapsto s_t$ è una soluzione valida.
\item La componente tangente non è determinata dalla richiesta che le traiettorie siano sulla curva: potrà essere una qualunque funzione $(s,\dot{s})\mapsto \phi(s,\dot{s})$
\end{enumerate}

Quindi assegnate le forze attive vi sarà un'unica reazione vincolare che soddisfa la richiesta di traiettorie sulla curva, per poter risolvere il problema del moto vincolato serviranno delle ipotesi sulla $\phi_\tau$ dato che altrimenti per ogni scelta si troverebbero sempre una componente normale e una binormale che soddisfano le equazioni.
La funzione di $\phi_n$ e $\phi_b$ è soltanto mantenere il punto sulla curva, esse non influenzano in alcun modo il suo moto tangente.

Diremo che la curva è \textbf{liscia} se la componente tangente della reazione vincolare è nulla $\phi_\tau=0$: il moto del punto tangenzialmente alla curva è determinato dalle sole forze attive $f_\tau$.



\section{Vincoli olonomi} 

Consideriamo un sistema di $N$ punti materiali in un riferimento $\Sigma$ nel quale sui punti agiscono forze, con rappresentativi $f_1,...,f_N:\mathbb{R}^{3N}\times\mathbb{R}^{3N}\rightarrow\mathbb{R}^3$, che d'ora in avanti diremo \textbf{attive}.

\begin{definizione}
Un vincolo è una restrizione sugli atti di moto $(x_\alpha,\dot{x}_\alpha)$ consentiti ai singoli punti che si traduce nel fatto che gli atti di moto del sistema siano costretti a rimanere in un sottoinsieme dello spazio degli atti di moto $\mathbb{R}^{3N}\times\mathbb{R}^{3N}$.
\end{definizione}
Noi tratteremo soltanto vincoli olono, cioè vincoli che nascsono come restrizione sulle posizioni dei punti costringendo la configurazione del sistema $X$ ad appartenere ad un $Q\subset\mathbb{R}^{3N}$, per poter definire una dinamica differenziabile per il sistema vincolato serve che $Q$ sia una sottovarietà dello spazio delle configurazioni $\mathbb{R}^{3N}$.

Un vincolo olonomo causerà anche un vincolo sulle velocità dei punti, se i punti sono vincolati a stare su una curva o su una superficie le loro velocità non potranno essere normali e quindi gli \textbf{atti di moto} del sistema apparterranno al fibrato tangente
\[
TQ=\left\{ (X,\dot{X})\in\mathbb{R}^{6N} : X\in Q, \dot{X}\in T_XQ \right\}
\]


Un vincolo si dirà \textbf{fisso} se non dipende dal tempo e \textbf{mobile} se dipende dal tempo.

\begin{definizione}
Un sistema meccanico di $N$ punti è soggetto ad un vincolo olonomo fisso se il vettore configurazione del sistema $X$ è costretto ad appartenere ad una sottovarietà $Q$ di $\mathbb{R}^{3N}$, che supponiamo sia connessa.

La sottovarietà Q si dirà varietà delle configuarazioni o varietà vincolare.

Se $Q$ ha dimensione $n$ diremo che il sistema vincolato ha $n$ gradi di libertà.
\end{definizione}

Un sistema di punti senza vincoli si può riguardare come un sistema olonomo con varietà vincolare $Q=\mathbb{R}^{3N}$ quindi i risultati che troveremo si potranno interpretare anche per il caso libero.

In quanto sottovarietà la varietà vincolare $Q$ si potrà esprimere come :\begin{itemize}
\item insieme di livello di una sommersione

 Per ogni punto $X\in Q$ esistono un intorno $\tilde{V}$ di $X\in\mathbb{R}^{3N}$ ed una mappa differenziabile sommersiva $\Psi:\tilde{V}\rightarrow\mathbb{R}^{3N-n}$ per cui $Q\cap\tilde{V}=\Psi^{-1}(0)$
\item immagine di una parametrizzazione

 Per ogni punto di $X\in Q$ esistono un intorno $\tilde{U}$ di $X$ in $Q$, un aperto $U\subseteq\mathbb{R}^{n}$ ed una mappa differenziabile immersiva $\tilde{X}:U\rightarrow\tilde{U}\subset Q$, $\tilde{X}(q)=(\tilde{X}_1(q),...,\tilde{X}_{3N}(q))\in Q$
\end{itemize}


La descrizione come insieme di livello in meccanica si usa soprattutto per definire la varietà delle configurazioni traducendo i vincoli in $r=3N-n$ equazioni
\[\Psi_1(X)=0 ,..., \Phi_r(X)=0\]
la descrizione in coordinate (parametrica) è quella più usata nello studio dei sistemi vincolati: in particolare per un moto $t\mapsto X_t\in Q$ se ne studia la $t\mapsto q_t\in U\subseteq\mathbb{R}^n$ con \[ X_t=\tilde{X}(q_t)
\]
Le coordinate $q=(q_1,...,q_n)$ sulla varietà vincolare di un sistema olonomo sono chiamate \textbf{coordinate lagrangiane}. Per evidenziare la parametrizzazione delle configurazioni dei singoli punti useremo
\[
\tilde{X}=(\tilde{x}_1,...,\tilde{x}_N) \hspace{0.15\linewidth} \tilde{x}_\alpha:U\rightarrow\mathbb{R}^3
\]
dove le mappe $\tilde{x}_\alpha$ danno le posizioni in coordinate affine dei singoli punti in funzione delle coordinate lagrangiane.


\subsection{Invarianza dello spazio degli atti di moto vincolati} \label{sezInvarianzaSpazioVincolato}
Un atto di moto $(x,\dot{X})$ si dirà vincolato se $X\in Q$ e $\dot{X}\in T_XQ$. Chiameremo reazione vincolare la $\Phi(X,\dot{X})=(\phi_1(X,\dot{X}),...,\phi_N(X,\dot{X}))$ che sarà una mappa da $\mathbb{R}^{6N}$ a $\mathbb{R}^{3N}$ come le forze attive ma che avrà significato solo sugli atti di moto vincolati.
Le equazioni del moto allora si potranno riscrivere come
\begin{equation}
M\ddot{X}=F(X,\dot{X})+\Phi(X,\dot{X})
\end{equation}
oppure come sistema del primo ordine in $\mathbb{R}^{6N}$
\begin{equation}\label{eqAttiMotoVincolato}
\begin{cases}
\dot{X}=V \\
\dot{V}=M^{-1}F(X,V)+M^{-1}\Phi(X,V)
\end{cases}
\end{equation}
e per definizione la $\Phi$ dovrà essere tale che ogni soluzione $t\in(X,\dot{X})$ con dato iniziale in $TQ$ appartenga a $TQ$ per ogni tempo, cioè che $TQ$ sia una sottovarietà invariante per la \eqref{eqAttiMotoVincolato}, per quanto visto nella sezione \ref{sezSottovarietàInvarianti} questo vuol dire che la $\Phi$ dovrà rendere il campo $\left(V,M^{-1}(F+\Phi)\right)$ tangente a $TQ$. Abbiamo visto che per un singolo punto vincolato ad una curva esiste sempre una reazione vincolare che soddisfa queste richieste, ora vediamo il caso generale.

In ogni $(X,V)\in\mathbb{R}^{3N}\times\mathbb{R}^{3N}$ la forza agente e la reazione vincolare saranno vettori di $T_X\mathbb{R}^{3N}$. In ciascun $X\in Q$ della sottovarietà vincolare potremo definire un complemento ortogonale allo spazio tangente alla sottovarietà per ciascun spazio tangente a $\mathbb{R}^{3N}$
\[
\left(T_XQ\right)^\perp \oplus T_XQ = T_X\mathbb{R}^{3N}
\]
Ovviamente questa decomposizione varrà soltanto nei punti $X\in Q$ dato che altrove non è definito $T_X Q$.
Allora data una mappa $\Phi:\mathbb{R}^{3N}\times\mathbb{R}^{3N}\rightarrow\mathbb{R}^{3N}$ potremo definirne le componenti tangenti e ortogonali a $Q$ come le mappe
\[
\Phi^\parallel , \Phi^\perp : Q\times\mathbb{R}^{3N}\rightarrow \mathbb{R}^{3N}
\]
dove $\forall (X,V)\in Q\times\mathbb{R}^{3N}$
\[
\Phi(X,V)=\Phi^\parallel(X,V) + \Phi^\perp (X,V) \hspace{0.15\linewidth} \Phi^\parallel(X,V)\in T_XQ \quad \Phi^\perp(X,V)\in (T_XQ)^\perp
\]

Abbiamo definito le componenti della reazione vincolare in $Q\times\mathbb{R}^{3n}$, cioè per $X\in Q$ sono definite per ogni $V\in T_X\mathbb{R}^{3N}$, adesso però avremo bisogno soltanto delle velocità nel sottospazio tangente a $Q$ cioè delle restrizioni a $TQ$ delle due mappe.

\begin{proposizione}\label{propUnicaReazione}
Dati un sistema meccanico di $N$ punti soggetti a forza $F$ attiva e $Q$ una sottovarietà dello spazio delle configurazioni $\mathbb{R}^{3N}$ per ogni scelta di $\left. \Phi^\parallel \right|_{TQ}$ vi è una ed una sola $\left. \Phi^\perp \right|_{TQ}$ tale per cui il campo $\left( V, M^{-1} (F+\Phi) \right)$ è tangente a $TQ$.
\end{proposizione}
\begin{dimo}
Assumiamo per semplicità che $Q$ sia individuata globalmente da $\Psi^{-1}(0)$ con un' unica $\Psi:\mathbb{R}^{6N}\rightarrow\mathbb{R}^r$.
\begin{itemize}
\item \emph{Il fibrato tangente $TQ$ è un insieme di livello del sollevamento $T\Psi$}:

In ogni punto $X\in Q$ per definizione $T_XQ=\ker{\Psi'}$ e quindi $\Psi'(X)\dot{X}=0$ da cui
\begin{equation} (X,V)\in TQ \iff \begin{cases} \Psi(X)=0 \\ \Psi'(X)V=0\end{cases} \end{equation}
La matrice Jacobiana di $T\Psi$ è 
\[
(T\Psi)'=\left(\dfrac{\partial T\Psi}{\partial X}, \dfrac{\partial T\Psi}{\partial V}\right)=\left(
\begin{array}{cc}
\Psi'(X) & 0 \\
\Psi''(X)V & \Psi'(X) \\
\end{array}\right)
\]
e dato che $\Psi'(X)$ ha rango $r$ (per ipotesi) la $(T\Psi)'$ avrà rango $2r$ in ogni $X\in Q$ e quindi $T\Psi$ è una sommersione.

Chiameremo $B(X,V):=\Psi''(X)V$.

\item Lo spazio tangente a $TQ$ sarà quindi dato dal $\ker$ di $(T\Psi)'$ quindi il campo $(V,M^{-1}(F+\Phi))$ per essere tangente a $TQ$ dovrà dare
\[
(T\Psi)' \left(V,M^{-1}(F+\Phi)\right)=\left(
\begin{array}{cc}
\Psi'{\scriptstyle(X)} & 0 \\
\Psi''{\scriptstyle (X,V)} V\cdot V & \Psi'{\scriptstyle(X)} \\
\end{array}\right) \left( \begin{array}{c}
V \\
M^{-1}(F+\Phi){\scriptstyle(X,V)} \\
\end{array}\right)=\left(\begin{array}{c}
0 \\ 0 
\end{array}\right)
\]
in ogni punto $(X,V)\in TQ$.

La $\Psi'{\scriptstyle (X)}V=0$ è automaticamente soddisfatta perchè $V\in T_XQ=\ker{\Psi'}$ e quindi la condizione perchè $\left(V,M^{-1}(F+\Phi)\right)$ sia tangente a $TQ$ è la sola 
\[
B{\scriptstyle (X,V)}V + \Psi'{\scriptstyle (X)}M^{-1}(F+\Phi){\scriptstyle (X,V)}=0 \quad \forall (X,V)\in TQ
\]
che riscriviamo come
\[
\left[ BV + \Psi' M^{-1} (F+\Phi) \right|_{TQ}=0
\]

\item Decomponendo $\left.\Phi\right|_{TQ}= \left.\Phi^\parallel\right|_{TQ} + \left.\Phi^\perp\right|_{TQ}$ possiamo riscrivere la condizione di tangenza come
\[
\Psi'M^{-1}\left.\Phi^\perp\right|_{TQ} = \left[ -BV- \Psi'M^{-1}\left(F+\Phi^\parallel\right)\right|_{TQ}
\]
\item Il complemento ortogonale a $T_XQ$ in $T_X\mathbb{R}^{3N}$ è generato dalla $\nabla\Psi_1,...,\nabla\Psi_r$ (ricordiamo che $T_XQ=\ker{\Psi'}$ è generato da $3n-r$ vettori ortogonali alle $\nabla\Psi_i$) quindi
\[
\Phi^\perp=\sum_i^{r} \lambda_i \nabla\Psi_i = \Psi'^T \lambda
\]
dove i $\lambda_i$ variano da punto a punto di $TQ$ (dato che i $\nabla\Psi_i$ danno $T_XQ$) e quindi si possono considerare come funzioni $\lambda=(\lambda_1,...,\lambda_r) :TQ\rightarrow\mathbb{R}^r$.
\item Possiamo quindi riscrivere la condizione per cui $\left(V,M^{-1}(F+\Phi)\right)$ è tangente a $TQ$ come
\[
\begin{cases}
\left.\Phi^\perp\right|_{TQ}=\Psi'^T\left.\lambda\right|_{TQ}\\
\Psi'M^{-1}\Psi'^T \left.\lambda\right|_{TQ} = -\left[BV - \Psi'M^{-1}(F+\Phi^\parallel)\right|_{TQ}
\end{cases}
\]
La matrice $M$ è definita positiva e $\Psi'$ ha rango $r$ in $Q$ quindi la matrice $A:=\Psi'(X)M^{-1}\Psi'^T(X)$ (simmetrica $r\times r$) è definita positiva ed è quindi invertibilein ogni $X\in Q$. Quindi possiamo ricavare
\[
\lambda=A^{-1} \left[ BV+\Psi'M^{-1}(F+\Phi^\parallel) \right|_{TQ}
\]
e quindi $(V,M^{-1}(F+\Phi))$ è tangente a $TQ$ solo per
\begin{equation}\label{eqReazioneVincolarePerpLagrange1}
\begin{split}
\left.\Phi^\perp\right|_{TQ} &= \Psi'^T A^{-1} \left[ BV+\Psi'M^{-1}(F+\Phi^\parallel)  \right|_{TQ} =\\
& = \Psi'^T \left[ \Psi'M^{-1}\Psi'^T \right]^{-1} \left[ \Psi''V\cdot V+\Psi'M^{-1}(F+\Phi^\parallel)  \right|_{TQ} 
\end{split} 
\end{equation}
che significa che la componente ortogonale è unica in $TQ$.
\end{itemize}
\end{dimo}

\paragraph{Riassumendo} un vincolo olonomo fisso è dato dall'assegnazione di una varietà vincolare $Q$ e dalla specificazione della componente tangente a $Q$ della reazione vincolare. Il vincolo eserciterà complessivamente una forza $\Phi^\parallel+\Phi^\perp$ sugli atti di moto vincolati dove per ogni $\Phi^\parallel_{TQ}$ esiste un'unica $\Phi^\perp|_{TQ}$, data dalla \eqref{eqReazioneVincolarePerpLagrange1}, che mantiene gli atti di moto in $TQ$.
L'equazione del moto vincolato sarà la restrizione a $TQ$ del sistema \[
\begin{cases} \dot{X}=V \\ \dot{V}=M^{-1}(F+\Phi)\end{cases}
\] per una qualunque $\Phi:\mathbb{R}^{3N}\times\mathbb{R}^{3N}\rightarrow\mathbb{R}^{3N}$ la cui restrizione a $TQ$ sia $\Phi|_{TQ}=\Phi^\parallel|_{TQ} +\Phi^\perp|_{TQ}$.


\subsection{Vincoli olonomi fissi ideali}
\begin{definizione}
Un vincolo olonomo fisso con varietà vincolare $Q$ si dice ideale se la componente tangente delle reazioni vincolari che esso esercita è identicamente nulla sugli atti di moto vincolati.
\[\left.\Phi^\parallel\right|_{TQ}\equiv0\]
\end{definizione}

Un vincolo ideale mantiene le configurazioni sulla varietà senza però influenzare in alcun modo il moto lungo $Q$, questo significa che le reazioni vincolari esercitate da un vincolo ideale su ogni atto di moto vincolato $(X,V)$ sono ortogonali\footnotemark{} a $T_XQ$
\begin{equation}
\Phi_{Id}(X,V)\in \left( T_XQ \right)^\perp \quad \forall (X,V)\in TQ
\end{equation}
cioè $\forall W\in T_XQ$ vale $\Phi_{Id}(X,V)\cdot W =0 \quad \forall (X,V)\in TQ$.

\footnotetext{Attenzione: l'ortogonalità è in $\mathbb{R}^{3N}$ non nello spazio vettoriale $\mathbb{V}^{3N}$! Se il sistema è composto da più di un punto materiale la differenza è significativa.}

\begin{proposizione}
In ogni moto vincolato da un vincolo olonomo fisso ideale la reazione vincolare non compie lavoro sul sistema vincolato.
\end{proposizione}
\begin{dimo}
Il lavoro compiuto da una forza $\varphi_\alpha(X_t,\dot{X}_t)$ su un punto $P_\alpha$ del sistema è
\[
\int_{t_1}^{t_2} \varphi_\alpha(X_t,\dot{X}_t) \cdot \dot{x}_\alpha(t) \, dt
\]
e il lavoro totale sul sistema sarà la somma di tutti i lavori sui singoli punti, che per come abbiamo definito le forze e velocità del sistema corrisponde a:
\[
\int_{t_1}^{t_2} \Phi(X_t,\dot{X}_t)\cdot\dot{X}_t \, dt
\]
chiaramente dato che per ogni $t$ $\dot{X}_t\in T_{X_t}Q$ e $\Phi\in \left(T_{X_t}Q\right)^\perp$ il lavoro della reazione vincolare ideale è nullo.
\end{dimo}

Questo significa che se un sistema è conservativo anche introducendo un vincolo olonomo fisso ideale continuerà ad esserlo.

Quindi da quanto visto nella sezione \ref{sezInvarianzaSpazioVincolato} e dall'equazione \eqref{eqReazioneVincolarePerpLagrange1} si trova che:
\begin{importante}
L'equazione del moto soggetto a vincoli olonomi fissi ideali è
\begin{equation}\label{eqLagrangePrimaSpecieFisso}
M\ddot{X}=F(X,\dot{X})+\Phi_{Id}(X,\dot{X})
\end{equation}
dove la reazione vincolare ideale è data da
\begin{equation}\label{eqReazioneLagrangePrimaSpecieFisso}
\Phi_{Id}:=-\Psi'^T\left[ \Psi' M^{-1} \Psi'^T \right]^{-1} \left[ \Psi'' V\cdot V+\Psi' M^{-1}F \right] 
\end{equation}
\end{importante}

Queste equazioni sono chiamate \textbf{equazioni di Lagrange di prima specie}.

In ogni punto $X\in Q$ la matrice $\Pi(X)$
\[
\Pi:=\Psi'^T\left( \Psi'M^{-1}\Psi'^T \right)^{-1}\Psi'M^{-1}
\]
è un proiettore ortogonale ($\Pi^2=\Pi$) di $\text{Im}\Psi'^T(X)=\left(\ker{\Psi'(X)}^\perp\right)=(T_XQ)^\perp$ e nella \eqref{eqReazioneLagrangePrimaSpecieFisso} il termine $-\Pi F$ di $\Phi$ compensa le componenti ortogonali della forza attiva, il termine quadratico in $V$ fornisce invece la \textit{forza centripeta} necessaria a curvare la traiettoria per mantenerla su $Q$.


Queste equazioni sono scomode, nella sezione \ref{sezLagrangeFisso} ne ricaveremo una controparte ristretta a $TQ$ ed espressa in coordinate locali in cui, come vedremo, la reazione vincolare non compare.

\subsection{Vincoli olonomi mobili}

Perchè un sistema di $N$ punti sia sottoposto ad un vincolo olonomo mobile se ad ogni tempo $t$ il vettore delle configurazioni $X_t$ è ristretto istantaneamente ad appartenere ad una ``sottovorietà mobile'' $Q_t$.

\begin{definizione}[Sottovarità mobile]
Una famiglia di sottovarietà $\{Q_t : t\in\mathbb{R}\}$ di $\mathbb{R}^{3N}$ è una sottovarietà mobile di $\mathbb{R}^{3N}$ se esiste una mappa differenziabile
\[
D:Q_0\times\mathbb{R}\rightarrow\mathbb{R}^{3N} \hspace{0.15\linewidth} D(X,T)=D_t(X)
\]
dove per ogni $t$ fissato la mappa $D_t:Q_0\rightarrow\mathbb{R}^{3N}$ \begin{itemize}
\item $D_t(Q_0)=Q_t$
\item $D_t$ è un diffeomorfismo tra $Q_0$ e $Q_t$.
\end{itemize}
\end{definizione}

A partire dalla sommersione $\Psi_0:\mathbb{R}^{3N}\rightarrow\mathbb{R}^r$ che definisce $Q_0=\Psi_0^{-1}(0)$ potremo definire  $\Psi_t$
\[
\Psi_t:=\Psi_0\circ D_t^{-1} :\mathbb{R}^{3N}\rightarrow\mathbb{R}^r
\]
che darà $Q_t=\Psi_t^{-1}(0)$.\footnote{Per ogni $X\in Q_t$ $D_t^{-1}(X)=X_0\in Q_0$ e $\Psi_0(X_0)=0$ per ipotesi.}
La mappa 
\[
\Psi:\mathbb{R}^{3N}\times\mathbb{R}\rightarrow\mathbb{R}^r \hspace{0.15\linewidth} \Psi(X_t,t)\mapsto\Psi_t(X_t)
\]
è una \textbf{sommersione dipendente dal tempo}

Similmente data la descrizione parametrica $\tilde{X}_0:U\rightarrow Q_t$ di $Q_0$ potremo definire una $\tilde{X}_t$
\[
\tilde{X}_t:=D_t\circ\tilde{X}_0 :U\rightarrow Q_t
\]
che parametrizzerà $Q_t$ con le stesse $u_1,...,u_n$ con cui $\tilde{X}_0$ parametrizza $Q_0$. La mappa
\[
\tilde{X}:U\times\mathbb{R}\rightarrow\mathbb{R}^{3N} \hspace{0.15\linewidth} \tilde{X}(q,t)\mapsto\tilde{X}_t(q)
\]
è una \textbf{parametrizzazione dipendente dal tempo}.


Ad ogni istante $t$ chiameremo la varietà $Q_t$ \textbf{varietà vincolare istantanea}.

La principale differenza tra gli atti di moto vincolati da un vincolo olonomo fisso e quelli vincolati da un vincolo olonomo mobile è che in questi ultimi il vettore velocità non è in generale tangente alla varietà vincolare istantanea. Infatti derivando rispetto al tempo la $\Psi(X_t,t)=0$\footnotemark{} troviamo
\[
\Psi'_t(X_t)\dot{X}_t +\dfrac{\partial \Psi}{\partial t}(X_t,t)=0
\]
e quindi se $t\mapsto X_t\in Q_t$ è una curva differenziabile la velocità $\dot{X}_t$ non sarà in $T_X Q_t$ a meno che il termine $\dfrac{\partial \Psi}{\partial t}$ non sia nullo.

\footnotetext{
Per essere chiari: la $\Psi$ dipende sia dalle configurazioni del sistema sia dal tempo, la $\Psi_t$ invece è calcolato per un $t$ fissato. Quindi la $\Psi_t'$ è data dalle derivate rispetto alle configurazioni e le funzioni contenute nella $\Psi$ qui saranno costanti.

$\Psi'$ è il differenziale di una funzione composta e, come visto in \textit{Analisi 2}, sarà dato da
$\sum_i \dfrac{\Psi_t}{\partial x_i} dx_i + \dfrac{\partial \Psi}{\partial t}dt$ dove indicheremo $dx_i$ con $\dot{x}_i$.}

Un atto di moto $(X,V)$ si dirà vincolato a $\{ Q_t : t\in\mathbb{R}\}$ all'istante $t$ se
\begin{equation}
\Psi'_t(X_t) V + \dfrac{\partial \Psi}{\partial t} (X_t,t)=0
\end{equation}
e diremo anche che $(X,V)$ è compatibile con il vincolo all'istante $t$.


Denoteremo con $\mathcal{V}_XQ_t$ il sottoinsieme di $T_X\mathbb{R}^{3N}$ formato dai vettori velocità compatibili col vincolo all'istante $t$
\[
\mathcal{V}_XQ_t=\left\{  V\in T_X\mathbb{R}^{3N} : \Psi'_t(X)V+\dfrac{\partial \Psi}{\partial t}(X_t,t)=0\right\}
\]
e con $\mathcal{V}Q_t$ il sottoinsieme di $\mathbb{R}^{3N}\times\mathbb{R}^{3N}$   contenente gli atti di moto compatibili col vincolo
\[
\mathcal{V}Q_t=\left\{  (X,V) : X\in Q_T , V\in\mathcal{V}_XQ_t \right\}
\]
\begin{esempio}[Pendolo di lunghezza variabile]
Consideriamo un pendolo la cui lunghezza vari nel tempo secondo una $R(t)$ nota. La varietà vincolare istantanea $Q_t$ sarà descritta da
\[
\Psi(x,y,z,t)=\left(\begin{array}{c} y \\ x^2+y^2+z^2-R^2(t)=0 \end{array}\right)=\left(\begin{array}{c} 0 \\ 0 \end{array} \right)
\] 
Allora
\[
\Psi'_t=\left(\begin{array}{ccc} 0 & 1 & 0 \\ 2x & 2y & 2z \end{array}\right)
\hspace{0.15\linewidth} \dfrac{\partial \Psi_t}{\partial t} =\left(\begin{array}{c} 0 \\ -2R(t)\dot{R}(t) \end{array}\right)
\]
Lo spazio tangente a $Q_t$ è $T_{(x,0,z)}Q_t=\left\{ \left(\begin{array}{c} u \\ v \\ w \end{array} \right) : v=0,\> xu + zw=0 \right\} = \left\{ \left\langle(z,0,-x)\right\rangle \right\}$ mentre $\mathcal{V}_{(x,0,z)}Q_t$ 
\begin{align*}
\mathcal{V}_{(x,0,z)}Q_t & =\left\{ \left(\begin{array}{c} u\\v\\w\end{array}\right) : \left(\begin{array}{ccc} 0 & 1 & 0 \\ 2x & 0 & 2z \end{array}\right) \left(\begin{array}{c} u \\ v \\ w \end{array}\right) = \left(\begin{array}{c} 0 \\ 2R(t)\dot{R}(t) \end{array}\right)\right\} \\
\end{align*}
$\mathcal{V}_{(x,0,z)}Q_t$ è uno spazio affine di $\mathbb{R}^3$ parallelo a $T_{(x,0,z)}Q_t$ ma traslato rispetto a lui 
\[\mathcal{V}_{(x,0,z)}Q_t=T_{(x,0,z)}Q_t + \dfrac{R\dot{R}}{\left|\left| (x,0,z) \right| \right|}\left(\begin{array}{c}x \\ 0 \\ z \end{array}\right)\]
\end{esempio}

La trattazione delle reazioni vincolari per i vincoli olonomi mobili è più complicata e la saltiamo.
Considerando anche forze che dipendono dal tempo l'equazione del moto è
\[
M\ddot{X}=F(X,\dot{x},t)+\Phi(X,\dot{X},t)
\]
Esiste un analogo della proposizione \ref{propUnicaReazione} per le reazioni vincolari per cui il sistema rimane in $\{Q_t:t\in\mathbb{R}\}$ che ci assicura che
\begin{proposizione}
Per ogni scelta di $\Phi^\parallel|_{\mathcal{V}Q_t}$, con $t\in\mathbb{R}$, esiste una ed una sola $\Phi^\perp|_{\mathcal{V}Q_t}$, con $t\in\mathbb{R}$, per cui per ogni dato iniziale $(X_0,\dot{X_0})\in\mathcal{V}Q_t$, e istante iniziale $t_0\in\mathbb{R}$, la soluzione di $M\ddot{X}=F(X,\dot{x},t)+\Phi(X,\dot{X},t)$ con quel dato iniziale appartiene a $\{Q_t:t\in\mathbb{R}\}$, che sarà quasi invariante per il l'equazione differenziale.
\end{proposizione}

Potremo definire i vincoli ideali in modo analogo a quanto fatto per quelli fissi ma facendo attenzione che la reazione vincolare sarà ortogonale a $T_XQ_t$  e non allo spazio delle velocità $\mathcal{V}Q_t$.

\begin{definizione}
Un vincolo olonomo mobile si dice ideale se, ad ogni istante $t\in\mathbb{R}$ e per ogni configurazione $X\in Q_t$, la componente tangente a $Q_t$ della reazione vincolare esercitata negli atti di moto vincolati è nulla, ovvero
\[
\Phi(X,V,t)\cdot W = 0 \quad \forall t\in\mathbb{R}, \forall (X,V)\in\mathcal{V}Q_t, \forall W\in T_XQ_t
\]
\end{definizione}

Con questa definizione è chiaro che le reazioni vincolari per vincoli ideali mobili \textbf{possono compiere lavoro} sul sistema.

\begin{importante}
I moti vincolati da un vincolo olonomo mobile ideale sono le soluzioni con dati iniziali compatibili col vincolo delle \textbf{equazioni di Lagrange di prima specie}

\begin{equation} \label{eqLagrangePrimaSpecieMobile}
M\ddot{X}=F(X,\dot{X},t)+\Phi_{id}(X,\dot{X},t)
\end{equation}

\end{importante}


\section{Equazioni di Lagrange} \label{sezLagrange}

Le equazioni del moto di un sistema, libero o soggetto a vincoli olonomi ideali, scritte in coordinate locali sulla varietà vincolare sono equazioni differenziali del secondo ordine e si chiamano equazioni di Lagrange.

\subsection{Parametrizzazione degli atti di moto}

Fino ad ora abbiamo scritto le equazioni del moto per i rappresentativi in $\mathbb{R}^{3N}$ però sappiamo che le varietà delle configurazioni ammette una descrizione parametrica (almeno locale) con le coordinate lagrangiane $q_1,...,q_n$.


\paragraph{Coordinate e vincolo non dipendenti dal tempo}
Sia $Q\subseteq\mathbb{R}^{3N}$ la varietà delle configurazioni e $\tilde{X}:U\rightarrow \tilde{U}\subseteq Q$ una sua parametrizzazione locale (con $U\subseteq\mathbb{R}^n$).

$\tilde{X}$ è un diffeomorfismo fra $U$ e $\tilde{U}$ e quindi, come visto nella sezione \ref{sezSollevamentoTangente}, $T\tilde{X}$ è un diffeomorfismo fra $TU$ e $T\tilde{U}$
\[
(q,v)\mapsto T\tilde{X}(q,v)=(\tilde{X}(q),\tilde{X}'(q) v)\in T\tilde{U}
\]
quindi $T\tilde{X}$ sarà una parametrizzazione di $TQ$ e la sua inversa sarà una carta locale
\[
(q,v):T\tilde{U}\rightarrow U\times\mathbb{R}^n
\]
Le coordinate $q$ sono chiamate coordinate lagrangiane e le $v$, tradizionalmente indicate con $\dot{q}$ anche se \textbf{non} sono le derivate temporali delle $q$, si chiamano \textbf{velocità generalizzate}.

L'atto di moto di coordinate $(q,\dot{q})$ è
\begin{equation}
T\tilde{X}(q,\dot{q}) = \left( \tilde{X}(q), \tilde{W}(q,\dot{q}) \right) = \left( \tilde{X}(q) , \tilde{X}'(q)\dot{q}\right)
\end{equation}

Chiaramente $\tilde{W}(q,\dot{q})\in T_{\tilde{X}(q)} Q$ e quindi $T\tilde{X}(q,\dot{q})$ è un atto di moto vincolato.

I vettori $\dfrac{\partial \tilde{X} } {\partial q_i} (q)$ sono indipendenti (perchè $\tilde{X}$ è un diffeomorfismo) e tangenti a $Q$ in $\tilde{X}(q)$, quindi sono $n$ vettori indipendenti in $T_{\tilde{X}(q)}Q$ spazio con dimensione $n$ e ne formano una base, detta \textbf{base naturale}.
Scrivendo in moto più esplicito la $\tilde{W}(q,\dot{q})$ 
\[
\tilde{X}'(q)\dot{q}=\sum_i^n \dfrac{\partial \tilde{X} }{\partial q_i}(q)\, \dot{q}_i \in T_{\tilde{X}(q)}Q
\]
Quindi se un atto di moto $\left(\overline{X},\dot{\overline{X}}\right)$ ha coordinate $(\overline{q},\dot{\overline{q}})$ vuol dire che $\overline{X}=\tilde{X}(\overline{q})$ e che le $\dot{\overline{q}}_i$ sono le componenti di $\dot{\overline{X}}$ nella base naturale di $T_{\overline{X}}Q$.

La base naturale (e quindi anche la parametrizzazione delle velocità) dipendono dalle $q$.

Evidenziando i singoli punti materiali $\tilde{X}=(\tilde{x}_1,...,\tilde{x}_N)$ la velocità del punto $\alpha$ nell'atto di moto a coordinate $(q,\dot{q})$ è
\begin{equation}
\tilde{w}_\alpha(q,\dot{q})=\sum_{i=1}^n \dfrac{\partial \tilde{x}_\alpha}{\partial q_i} (q) \, \dot{q}_i
\end{equation}
e questa in generale dipenderà da tutte le $q_1,...,q_n,\dot{q}_1,...,\dot{q}_n$.

\begin{esempio}[Punto libero in coordinate cilindriche]
$Q=\mathbb{R}^3$, possiamo parametrizzare in coordinate cilindriche
\[
\tilde{X}(r,\varphi,z)=(r\cos{\varphi}, r\sin{\varphi} , z)
\]
La base naturale sarà data da
\[
\dfrac{\partial \tilde{X} }{\partial r} = ( \cos{\varphi} , \sin{\varphi} , 0 ) \hspace{0.15\linewidth} 
\dfrac{\partial \tilde{X}}{\partial \varphi} = ( -r\sin{\varphi} , r\cos{\varphi} , 0 ) \hspace{0.15\linewidth}
\dfrac{\partial \tilde{X}}{\partial z} = ( 0 , 0 , 1 ) 
\]
e la velocità dell'atto di moto $(r,\varphi,z)$ è
\[
\tilde{w}(r,\varphi,z)= \left(\dot{r} \cos{\varphi} -\dot{\varphi}r\sin{\varphi},
\dot{r}\sin{\varphi}+\dot{\varphi}r\cos{\varphi}, \dot{z}\right)
\]
Se ora ponniamo $r=R$ costante e $z=0$ avremo un punto vincolato a stare su una circonferenza, che avrà
\[
\tilde{X}=(R\cos{\varphi},R\sin{\varphi},0) \qquad
\dfrac{ \partial \tilde{X} }{\partial \varphi} = ( -R\sin{\varphi} ,R\cos{\varphi},0)
\qquad
\tilde{w}(\varphi)=\dot{\varphi}(-R\sin{\varphi},R\cos{\varphi},0)
\]
\end{esempio}

\paragraph{Vincolo fisso e parametrizzazione dipendente dal tempo}
Una parametrizzazione dipendente dal tempo è una mappa
\[
\tilde{X}:U\times{R}\rightarrow \tilde{U}\subseteq Q \hspace{0.15\linewidth} \tilde{X}(q,t)=\tilde{X}_t(q) 
\]
dove la mappa $X_t$ parametrizza la varietà $Q_t$.

Il sollevamento tangente di $\tilde{X}$ parametrizza il fibrato tangente $TQ$ ma essendo $\tilde{X}$ dipendente dal tempo  comparirà un termine aggiuntivo nel suo push-forward, cioè nelle velocità (si veda la sezione \ref{sezCambiCoordinateDipendentiTempo}).

La velocità del moto $t\mapsto X_t$ con $X_t=\tilde{X}(q,t)$ sarà
\[
\dfrac{\partial}{\partial t} X_t = \dfrac{\partial \tilde{X}}{\partial q}(q,t)\dot{q} + \dfrac{\partial \tilde{X}}{\partial t}(q,t)
\]
e quindi un atto di moto vincolato con coordinate $(q,\dot{q})$\footnotemark sarà
\footnotetext{Specifichiamo che anche nella parametrizzazione dipendente dal tempo le $\dot{q}$ \textbf{non} sono le derivate temporali delle $q$, sono vettori tangenti a $U$ in $q$, quindi indipendenti dalle $q$.}
\begin{equation}
\left(\tilde{X}(q,t), \tilde{W}(q,\dot{q},t)\right)=\left( \tilde{X}(q,t), \dfrac{\partial\tilde{X}}{\partial q}(q,t)\dot{q} + \dfrac{\partial \tilde{X}}{\partial t}(q,t) \right)
\end{equation} 

In questo caso dato che $\tilde{X}(q,t)\in Q$ per tutti i $t$ la $\tilde{W}(q,\dot{q},t)$ è tangente a $Q$ in $\tilde{X}(q,t)$.

\paragraph{Vincoli mobili}

Per un vincolo mobile la varietà dipendente dal tempo è $\{ Q_t:t\in\mathbb{R}\}$ e avrà parametrizzazione locale 
\[
\tilde{X}:U\times{R}\rightarrow \mathbb{R}^{3N} \hspace{0.15\linewidth} \tilde{X}(q,t)=\tilde{X}_t(q)\in Q_t 
\]


Un atto di moto vincolato con coordinate $(q,\dot{q})$ sarà
\begin{equation}
\left(\tilde{X}(q,t), \tilde{W}(q,\dot{q},t)\right)=\left( \tilde{X}(q,t), \dfrac{\partial\tilde{X}}{\partial q}(q,t)\dot{q} + \dfrac{\partial \tilde{X}}{\partial t}(q,t) \right)
\end{equation} 
La velocità non è tangente a $Q_t$! Il primo termine $\dfrac{\partial \tilde{X}}{\partial q}\dot{q}$ sarà in $T_{\tilde{X}(q,t)}Q_t$ ma, dato che la varietà cambia nel tempo, il secondo termine in generale no.

\subsection{Energia cinetica}

\begin{definizione}
Definiamo energia cinetica di un sistema meccanico vincolato la restrizione agli atti di moto vincolati $\tilde{T}|_{TQ}$ dell'energia cinetica $\tilde{T}$
\[
\tilde{T}|_{TQ}:TQ\rightarrow\mathbb{R} \overset{\text{con}}{\hspace{0.15\linewidth}} \tilde{T}:\mathbb{R}^{3N}\times\mathbb{R}^{3N}\rightarrow \mathbb{R} \qquad \tilde{T}(X,V)=\dfrac{1}{2} V \cdot M V
\]
\end{definizione}
\begin{table}[H]
Chiameremo $T$ l'energia cinetica espressa nelle coordinate lagrangiane locali 
\begin{itemize}
\item Vincoli e coordinate fissi $T:U\times\mathbb{R}^n \rightarrow \mathbb{R}$ 
\[
T(q,\dot{q})=\tilde{T}\left(\tilde{X}(q),\tilde{X}'(q)\dot{q} \right)
\]
\item Vincoli mobili e parametrizzazioni dipendenti dal tempo $T:U\times\mathbb{R}^n \times\mathbb{R} \rightarrow \mathbb{R}$ 
\begin{equation}
T(q,\dot{q},t)=\tilde{T} \left( \tilde{X}(q), \dfrac{\partial\tilde{X}}{\partial q} (q,\dot{q},t)\dot{q}+\dfrac{\partial \tilde{X}}{\partial t} (q,\dot{q},t)\right)
\end{equation}
\end{itemize}
\end{table}

\begin{proposizione}
Per un sistema meccanico vincolato da vincoli olonomi l'energia cinetica nelle coordinate lagrangiane è data dà:
\begin{enumerate}[i.]
\item per vincoli e parametrizzazioni fissi 
\[
T(q,\dot{q})=\dfrac{1}{2}\dot{q} \cdot A(q) \dot{q}
\]
con $A(q)\in M_{n\times n}$ simmetrica e definita positiva $\forall q\in U$.
\item per vincoli e parametrizzazioni dipendenti dal tempo
\begin{align*}
T(q,\dot{q},t)&= \dfrac{1}{2} \dot{q} \cdot A(q,t)\dot{q} + b(q,t) \dot{q} + T_0(q,t)=\\
&=T_2(q,\dot{q},t)+T_1(q,\dot{q},t)+T_0(q,t)
\end{align*}
con $A(q,t)\in M_{n\times n}$ simmetrica e definita positiva $\forall q\in U$, $b(q,t)\in\mathbb{R}^n$ e $T_0:\mathbb{R}^n\times\mathbb{R}\rightarrow\mathbb{R}$.
\end{enumerate}
\end{proposizione}
\begin{dimo}
\begin{enumerate}[i.]
\item $\tilde{T} =\dfrac{1}{2}V\cdot MV \longrightarrow T(q,\dot{q})=\dfrac{1}{2} \tilde{X}'\dot{q} \cdot M \tilde{X}'\dot{q} = \dfrac{1}{2} \dot{q} \cdot\left((\tilde{X}')^T  M \tilde{X}' \right)\dot{q} $ e $A(q)=(\tilde{X}')^T \cdot M \tilde{X}'$ è simmetrica e definita positiva.
\item $\tilde{T} =\dfrac{1}{2}V\cdot MV \longrightarrow T(q,\dot{q})=\dfrac{1}{2} \left[ \tilde{X}'\dot{q} \cdot M \tilde{X}'\dot{q} + 2 (\tilde{X}')^T M \dfrac{\partial \tilde{X}}{\partial t}  \dot{q} + \dfrac{\partial \tilde{X}}{\partial t}\cdot M \dfrac{\partial \tilde{X}}{\partial t}
\right] $ 

Che dà la tesi con:
$A(q,t)=\left(\tilde{X}'(q,t)\right)^T M \tilde{X}'(q,t) $ simmetrica e definita positiva\\ $b(q,t)= \left(\tilde{X}'(q,t)\right)^T M \dfrac{\partial \tilde{X}}{\partial t} (q,\dot{q},t) $ \\ $T_0(q,t) = \dfrac{1}{2}\dfrac{\partial \tilde{X}}{\partial t}(q,t)\cdot M \dfrac{\partial \tilde{X}}{\partial t}(q,t) :U\times\mathbb{R}\rightarrow\mathbb{R}$\qedhere
\end{enumerate}
\end{dimo}

La matrice $A(q,t)$ si chiama \textbf{matrice cinetica} del sistema relativa alle coordinate $(q,\dot{q})$.

Per calcolare l'energia cinetica è spesso più utile però l'espressione data dalla sommatoria delle energie dei singoli punti
\begin{equation}
T(q,\dot{q},t)=\dfrac{1}{2} \sum_{\alpha=1}^N m_\alpha \left| \left| \left( \sum_{j=1}^n \dfrac{\partial \tilde{x}_\alpha}{\partial q_j} (q,t) \dot{q}_j \right) + \dfrac{\partial \tilde{x}_\alpha}{\partial t} (q,t) \right| \right|^2
\end{equation}

\begin{esempio}[Due punti liberi in coordinate del centro di massa e relative]
Per il sistema di due punti liberi (lo spazio delle configurazioni è $\mathbb{R}^6$) con vettori posizioni $x_1$ e $x_2$ possiamo usare le coordinate 
\[
R=\dfrac{x_1m_1+x_2m_2}{m_1+m_2} \hspace{0.15\linewidth}
r=x_1-x_2
\]
che, chiamando $M=m_1+m_2$, daranno
\[
\tilde{x}_1 (R,r)=R+\dfrac{m_2}{M} r \hspace{0.15\linewidth}
\tilde{x}_2 (R,r) = R-\dfrac{m_1}{M}r
\]
con corrispondenti velocità
\[
\tilde{w}_1(R,r) = \dot{R} + \dfrac{m_2}{M}\dot{r} \hspace{0.15\linewidth}
\tilde{w}_2(R,r) = \dot{R} - \dfrac{m_1}{M}\dot{r}
\]
Quindi l'atto di moto vincolato di coordinate $(R,r)$ è
\[
(x_1,x_2,\dot{x}_1,\dot{x}_2) = \left( \dot{R} + \dfrac{m_2}{M}\dot{r} , R-\dfrac{m_1}{M}r , \dot{R} + \dfrac{m_2}{M}\dot{r} , \dot{R} - \dfrac{m_1}{M}\dot{r} \right)
\]
e l'energia cinetica 
\begin{align*}
T(R,r,\dot{R},\dot{r})&=\dfrac{1}{2} m_1 \left|\left| \tilde{w} (R,r) \right|\right|^2 + \dfrac{1}{2} m_2 \left|\left| \tilde{w} (R,r) \right|\right|^2 =\\
& = \dfrac{1}{2} m_1 \left( \left|\left|\dot{R}\right|\right|^2 + \dfrac{m_2^2}{M^2}\left|\left|\dot{r}\right|\right|^2 \right)+ \dfrac{1}{2} m_2 \left(  \left|\left|\dot{R}\right|\right|^2 - \dfrac{m_1^2}{M^2}\left|\left|\dot{r}\right|\right|^2 \right)
\end{align*}
si dà per scontato che $\dot{R}\cdot\dot{r}=0$ dato che $\{\dot{R},\dot{r}\}$ è una base di $T_{R,r}\mathbb{R}^6$.

Chiamando $\mu=\dfrac{M-1m_2}{m_1+m_2}$ si trova
\[
T(R,r,\dot{R},\dot{r})=\dfrac{1}{2}M\left|\left| \dot{R} \right|\right|^2+ \dfrac{1}{2}\mu\left|\left|\dot{r}\right|\right|^2
\]
\end{esempio}

\subsection{Derivata lungo le curve}

Ricordiamo che per una parametrizzazione locale della varietà vincolare $\tilde{X}:U\times\mathbb{R}$ con $\tilde{X}(t,q)=X(t)\in Q_t$ gli atti di moto vincolati sono parametrizzati da $(q,\dot{q},t)\mapsto\left(\tilde{X}(q,t),\tilde{W}(q,\dot{q},t)\right)$ dove le $q\in U$ si chiamano coordinate lagrangiane locali e le $\dot{q}\in T_q U$ velocità generalizzate.

Da ora in avanti indicheremo con $(q,\dot{q})$ le coordinate locali in $TU$ e con $(q_t,\dot{q}_t)$ una curva differenziabile in $U$ ed il suo vettore tangente: ripetiamo che $\dot{q}$ non dipende da $q$ ed entrambi sono indipendenti dal tempo mentre $\dot{q}_t$ è la derivata di $q_t$ che è una funzione del tempo.

Dato l'insieme di funzioni differenziabili $C^{\infty}(U,\mathbb{R}^k) $ con $U\subseteq\mathbb{R}^n$ aperto
\begin{enumerate}[i.]
\item \textbf{Per funzioni di posizione}:

Chiamiamo derivata lungo le curve o derivata totale l'operatore
\begin{equation}
\dfrac{d}{dt} :C^{\infty}(U,\mathbb{R}^K) \rightarrow C^{\infty}(U\times\mathbb{R}^n,\mathbb{R}^k) \hspace{0.1\linewidth}
\dfrac{dg}{dt} (q,\dot{q}) = g'(q) \dot{q} =\sum_{i=1}^n \dfrac{\partial g}{\partial q_i}(q) \dot{q}_i
\end{equation}
con $\dot{q}\in T_qU$.\\
Lungo una qualsiasi curva $t\mapsto q_t$ l'operatore coinciderà con la normale derivate in $t$ della funzione composta $g(q_t)$:
\[
\dfrac{dg}{dt}(q,\dot{q}) \equiv \dfrac{\partial g(q_t)}{\partial t}=g'(q_t)\dfrac{\partial q_t}{\partial t}= g'(q_t)\dot{q}_t
\]

Potremo quindi calcolare la derivata lungo la curva ``immaginando'' che le $q$ siano funzioni del tempo e ottenere il risultato voluto, la funzione che si ottiene però sarà definita su $TU$ e non su $U$ dato che in realtà le $\dot{q}$ sono indipendenti dalle $q$.

Possiamo estendere questa operazione a funzioni dipendenti anche dal tempo definendo
\[
\dfrac{d}{dt} : C^{\infty}(U\times\mathbb{R},\mathbb{R}^k)\rightarrow C^{\infty}(U\times\mathbb{R}^n\times\mathbb{R},\mathbb{R}^k) \hspace{0.1\linewidth}
\dfrac{dg}{dt} (q,\dot{q}) =\sum_{i=1}^n \dfrac{\partial g}{\partial q_i}(q) \dot{q}_i + \dfrac{\partial g}{\partial t}(q,t)
\]
e anche questa coinciderà con la $\dfrac{\partial}{\partial t} g(q_t,t)$ lungo la curva $t\mapsto q_t$.

\item \textbf{Per funzioni di posizione e velocità}:
Chiamiamo derivata lungo le curve o derivata totale l'operatore
\begin{equation}\begin{split}
\dfrac{d}{dt} :C^{\infty}(U\times\mathbb{R}^n,\mathbb{R}^K) \rightarrow C^{\infty}(U\times\mathbb{R}^n\times\mathbb{R}^n,\mathbb{R}^k) \\
\dfrac{dg}{dt} (q,\dot{q},\ddot{q}) =\sum_{i=1}^n \dfrac{\partial g}{\partial q_i}(q,\dot{q}) \dot{q}_i + \dfrac{\partial g}{\partial \dot{q}_i}(q,\dot{q})\ddot{q}_i
\end{split}
\end{equation}
che si estenderà a funzioni del tempo 
\begin{equation}\begin{split}
\dfrac{d}{dt} :C^{\infty}(U\times\mathbb{R}^n\times\mathbb{R},\mathbb{R}^K) \rightarrow C^{\infty}(U\times\mathbb{R}^n\times\mathbb{R}^n\times\mathbb{R},\mathbb{R}^k) \\
\dfrac{dg}{dt} (q,\dot{q},\ddot{q},t) =\sum_{i=1}^n \dfrac{\partial g}{\partial q_i}(q,\dot{q},t) \dot{q}_i + \dfrac{\partial g}{\partial \dot{q}_i}(q,\dot{q},t)\ddot{q}_i+\dfrac{\partial g}{\partial t}(q,\dot{q},t)
\end{split}
\end{equation}
\end{enumerate}

Le $\ddot{q}$ si diranno \textbf{accelerazioni generalizzate}.


\subsection{Le equazioni di Lagrange nel caso fisso}\label{sezLagrangeFisso}

Le equazioni di un moto vincolato da un vincolo fisso ideale con varietà vincolare $Q$ e forze attive $F$ sono le restrizioni a $TQ$ di
$M\ddot{X}=F(X,\dot{X})+\Phi_{id}(X,\dot{X})$ dove $\Phi_{id}(X,V)\in (T_XQ)^\perp$ per ogni $(X,V)\in TQ$.

Vogliamo scrivere la restrizione a $TQ$ di queste equazioni usando le coordinate lagrangiane locali $(q,\dot{q})\mapsto(\tilde{X}(q),\tilde{W}(q,\dot{q})$

\begin{definizione}
Chiamiamo componenti Lagrangiane delle forza attive le $n$ funzioni
\[
Q_i : U \times \mathbb{R}^n \rightarrow\mathbb{R}^n  \hspace{0.15\linewidth} i = 1, \dots , n 
\]
definite da
\[
Q_i(q,\dot{q}) := F\left(\tilde{X}(q),\tilde{W}(q,\dot{q})\right) \cdot \dfrac{\partial \tilde{X}}{\partial q_i}(q)
\]
corrispondenti punto per punto alle proiezioni delle forze attive sui vettori della base naturale di $T_{\tilde{X}(q)}Q$.
\end{definizione}

Le $Q_i$ sono anche chiamate \textbf{forze generalizzate} e non sono componenti cartesiane delle $F$.

\begin{proposizione}[Equazioni di Lagrange (fisso)]\label{propLagrangeFisso}
La restrizione delle equazioni di Lagrange di prima specie per vincoli ideali fissi a $TQ$, in coordinate locali $(q,\dot{q})$, è data dal sistema di $n$ equazioni del secondo ordine 
\begin{equation}\label{eqLagrangeFisso}
\left(\dfrac{d}{dt}\dfrac{\partial T}{\partial \dot{q}_i} \right) (q,\dot{q},\ddot{q}) - \dfrac{\partial T}{\partial q_i}(q,\dot{q})=Q_i(q,\dot{q})
\end{equation}
Queste equazioni si possono porre in forma normale.
\end{proposizione}

La $T(q,\dot{q})$ è l'energia cinetica nelle $(q,\dot{q})$ con forma esplicita
\[
T(q,\dot{q})=\dfrac{1}{2} \dot{q} \cdot A(q) \, \dot{q} = \dfrac{1}{2} \sum_{h,k=1}^n A_{hk}(q)\, \dot{q}_h \dot{q}_k
\]
Quindi le derivate sono
\[
\dfrac{\partial T}{\partial q_i}(q,\dot{q}) = \dfrac{1}{2} \sum_{h,k=1}^n \dfrac{\partial A_{hk}(q)}{\partial q_i}\, \dot{q}_h \dot{q}_k
\]
\begin{align*}
\dfrac{\partial T(q,\dot{q})}{\partial \dot{q}_i} &= \dfrac{1}{2} \sum_{h,k=1}^n A_{hk}(q) \left( \dfrac{\partial \dot{q}_h}{\partial \dot{q}_i} \dot{q}_k + \dot{q}_h \dfrac{\partial \dot{q}_k}{\partial \dot{q}_i} \right) 
 = \dfrac{1}{2} \sum_{h,k=1}^n A_{hk}(q) \left( \delta_{h,i} \dot{q}_k + \dot{q}_h \delta_{k,i} \right) = \\
 & =  \sum_{h=1}^n A_{hi}(q) \dot{q}_h =  \sum_{h=1}^n A_{ih}(q) \dot{q}_h  = \left( A(q)\, \dot{q} \right)_i
\end{align*}
che con la derivata totale diventa
\begin{align*}
\dfrac{d}{dt}\dfrac{\partial T(q,\dot{q})}{\partial \dot{q}_i} &= 
\sum_{i,h=1}^n \dfrac{\partial A_{ih}(q)}{\partial q_k} \, \dot{q}_h \dot{q}_k + \sum_{i,h=1} A_{hk}(q) \, \dfrac{\partial \dot{q}_h}{\partial\dot{q}_k} \ddot{q}_i  \\
&= \sum_{h,k=1}^n \dfrac{\partial A_{ih}(q)}{\partial q_k}\, \dot{q}_h \dot{q}_k + \left( A(q)\, \ddot{q} \right)_i
\end{align*}
\begin{importante}
La forma esplicita delle equazioni di Lagrange \ref{eqLagrangeFisso} è
\begin{equation}
\left( A(q)\, \ddot{q} \right)_i +\sum_{h,k=1}^n \left( \dfrac{\partial A_{ih}(q)}{\partial q_k} - \dfrac{1}{2} \dfrac{\partial A_{hk}(q)}{\partial q_i}
\right)\dot{q}_h\dot{q}_k =Q_i(q,\dot{q})
\end{equation}


Dato che la matrice $A$ è invertibile queste equazioni si possono porre come equazioni del secondo ordine in forma normale.
\end{importante}

\begin{dimo}[Dimostrazione proposizione \ref{propLagrangeFisso}]
Se $t\mapsto X(t)=\tilde{X}(q_t)\in Q$ è una curva integrale dell'equazione differenziale allora $\dot{X}(t)=\tilde{W}(q_t,\dot{q_t})=\tilde{X}'(q)\dot{q}$ con $\tilde{W}(q,\dot{q})=\tilde{X}'(q)\dot{q}$ e
\[
\ddot{X}(t)= \tilde{X}''(q_t)\dot{q}_t^2 + \tilde{X}'(q_t)\ddot{q}_t\]
che per componenti è 
\[ 
\ddot{X}_s(t)=\sum_{h,k=1}^{n}\dfrac{\partial \tilde{X}_{sk}'(q_t)}{\partial q_h} (\dot{q}_t)_h (\dot{q}_t)_k + \sum_{k=1}^n \tilde{X}_{sk}'(q_t) (\ddot{q}_t)k
\]
che lungo la curva sarà uguale alla derivata totale
\[
\ddot{X}=\dfrac{d}{dt}\tilde{W}(q,\dot{q})=\dfrac{d\tilde{W}}{dt}(q,\dot{q},\ddot{q})=\sum_{h,k=1}^{n}\dfrac{\partial \tilde{W}_{sk}(q,\dot{q})}{\partial q_h} \dot{q}_h \dot{q}_k + \sum_{k=1}^n \tilde{W}_{sk}(q,\dot{q}) \ddot{q}_k
\]
Quindi $t\mapsto q_t$ soddisfa
\[
M\dfrac{d\tilde{W}}{dt}(q,\dot{q},\ddot{q})=F\left(\tilde{X}(q),\tilde{W}(q)\right)+\Phi_{id}\left(\tilde{X}(q),\tilde{W}(q,\dot{q}\right)
\]
e anche le equazioni
\begin{equation}\label{eqPrincipioDAlembert}
\dfrac{\partial \tilde{X}(q)}{\partial q_i} \cdot \left[ M\dfrac{d\tilde{W}}{dt}(q,\dot{q},\ddot{q})-F\left(\tilde{X}(q),\tilde{W}(q)\right)-\Phi_{id}\left(\tilde{X}(q),\tilde{W}(q,\dot{q}\right)
\right]=0
\end{equation}
I vettori $\dfrac{\partial\tilde{X}(q)}{\partial q_i}$ formano la base naturale di $T_{\tilde{X}(q)}Q$ e quindi $\dfrac{\partial\tilde{X}(q)}{\partial q_i}\cdot\Phi_{id}(\tilde{X}(q),\tilde{W}(q,\dot{q})$ per la condizione di idealità.
L'equazione quindi dà
\[
\dfrac{\partial \tilde{X}(q)}{\partial q_i} \cdot  M\dfrac{d\tilde{W}}{dt}(q,\dot{q},\ddot{q})
=Q_i(q,\dot{q})=\dfrac{\partial\tilde{X}(q)}{\partial q_i} \cdot F\left(\tilde{X}(q),\tilde{W}(q)\right)
\]
Dovremo quindi provare che il primo membro di questa equazione sia uguale a quello delle equazioni di Lagrange per dimostrare la proposizione.

Per il resto della dimostrazione omettiamo gli argomenti $q,\dot{q},\ddot{q}$ per alleggerire la notazione.

Dato che $\tilde{W}_{s}=\sum_k^n \tilde{X}'_{sk} \dot{q}_k$ con la derivata totale
\[
\dfrac{d\tilde{W}_s}{dt}=\sum_{h,k}^n \dfrac{\partial \tilde{X}'_{sk}}{\partial q_h}\dot{q}_h\dot{q}_k + \sum_{k}^n \tilde{X}'_{sk} \ddot{q}_k
\]
Ci ricordiamo che la matrice cinetica $A$ è definita da $A=(\tilde{X}')^T M \tilde{X}'$ e quindi ha componenti
\[
A_{ik}=\sum_{r,s}^n \tilde{X}'_{ri} M_{rs} \tilde{X}'_{sk}
\]
e vediamo che
\[
\left( \dfrac{\partial \tilde{X}}{\partial q_i} \right)_r=\dfrac{\partial \tilde{X}_r}{\partial q_i}=\tilde{X}'_{ri}
\]
Quindi tornando al primo membro dell'equazione esplicitando abbiamo
\begin{align*}
\dfrac{\partial \tilde{X}}{\partial q_i} \cdot M\dfrac{d\tilde{W}}{dt}&=
\sum_{r,s}^n \dfrac{\partial \tilde{X}_r}{\partial q_i} M_{rs} \dfrac{d\tilde{W}_s}{dt}= \sum_{r,s}^n \tilde{X}'_{ri} M_{rs} \left( \sum_{h,k}^n \dfrac{\partial \tilde{X}'_{sk}}{\partial q_h}\dot{q}_h\dot{q}_k + \sum_{k}^n \tilde{X}'_{sk} \ddot{q}_k \right)=\\
&=\sum_{r,s,h,k}^n \tilde{X}'_{ri}M_{rs}\dfrac{\partial \tilde{X}'_{sk}}{\partial q_h} \, \dot{q}_h\dot{q}_k +
\sum_{r,s,k}^n \tilde{X}'_{ri}M_{rs}\tilde{X}'_{sk} \, \ddot{q}_k = \\
&= \sum_{r,s,h,k}^n \tilde{X}'_{ri}M_{rs}\dfrac{\partial\tilde{X}'_{sk}}{\partial q_h} \, \dot{q}_h\dot{q}_k + \left( A \ddot{q} \right)_i
\end{align*}
Usando quanto visto prima sulle derivate dell'energia cinetica 

\[
\dfrac{d}{dt}\dfrac{\partial T}{\partial \dot{q}_i} - \dfrac{\partial T}{\partial q_i} = \left( A\, \ddot{q} \right)_i +\sum_{h,k=1}^n \left( \dfrac{\partial A_{ih}}{\partial q_k}  - \dfrac{1}{2} \dfrac{\partial A_{hk})}{\partial q_i} 
\right)\dot{q}_h\dot{q}_k 
\]
Quindi per provare la proposizione dovremo mostrare
\[
\sum_{h,k=1}^n \left( \dfrac{\partial A_{ih}}{\partial q_k}  - \dfrac{1}{2} \dfrac{\partial A_{hk}}{\partial q_i} 
\right)\dot{q}_h\dot{q}_k  =
\sum_{r,s,h,k}^n \tilde{X}'_{ri}M_{rs}\dfrac{\partial\tilde{X}'_{sk}}{\partial q_h} \, \dot{q}_h\dot{q}_k 
\]
Vediamo che
\begin{align*}
\dfrac{\partial A_{ih}}{\partial q_k}  - \dfrac{1}{2} \dfrac{\partial A_{hk}}{\partial q_i} &= 
\dfrac{\partial A_{ih}}{\partial q_k}  - \dfrac{1}{2} \dfrac{\partial A_{kh}}{\partial q_i} =\\
&=\dfrac{\partial}{\partial q_k} \left(
\sum_{r,s}^n \tilde{X}'_{ri} M_{rs} \tilde{X}'_{sh}
 \right) -\dfrac{1}{2}\dfrac{\partial}{\partial q_i} \left( 
\sum_{r,s}^n \tilde{X}'_{rk} M_{rs} \tilde{X}'_{sh}
 \right) \\
 & = \sum_{r,s}^n \left[\dfrac{\partial \tilde{X}'_{ri}}{\partial q_k} 
  M_{rs} \tilde{X}'_{sh}+ 
 \tilde{X}'_{ri} M_{rs}  \dfrac{\partial\tilde{X}'_{sh}}{\partial q_k} \right]
  -\dfrac{1}{2}\sum_{r,s}^n  \left[
\dfrac{\partial\tilde{X}'_{rk}}{\partial q_i} M_{rs} \tilde{X}'_{sh}
- \tilde{X}'_{rk} M_{rs} \dfrac{\partial\tilde{X}'_{sh}}{\partial q_i}
 \right]\\ 
&= \sum_{r,s}^n \left( \dfrac{\partial \tilde{X}'_{ri}}{\partial q_k} 
 -\dfrac{1}{2}\dfrac{\partial\tilde{X}'_{rk}}{\partial q_i} \right) M_{rs} \tilde{X}'_{sh} + \sum_{r,s}^n
 \tilde{X}'_{ri} M_{rs}  \dfrac{\partial\tilde{X}'_{sh}}{\partial q_k}  -\dfrac{1}{2}\sum_{r,s}^n \tilde{X}'_{rk} M_{rs} \dfrac{\partial\tilde{X}'_{sh}}{\partial q_i}
\end{align*}
Però
\[
\dfrac{\partial \tilde{X}'_{ri}}{\partial q_k}=\dfrac{\partial^2 \tilde{X}_{i}}{\partial q_k \partial q_i}= \dfrac{\partial \tilde{X}'_{rk}}{\partial q_i}
\hspace{0.15\linewidth}
\dfrac{\partial \tilde{X}'_{sh}}{\partial q_i}=\dfrac{\partial^2 \tilde{X}_{s}}{\partial q_i \partial q_h}= \dfrac{\partial \tilde{X}'_{si}}{\partial q_h}
\]
quindi
\[
\dfrac{\partial A_{ih}}{\partial q_k}  - \dfrac{1}{2} \dfrac{\partial A_{hk}}{\partial q_i} =
\sum_{r,s}^n \dfrac{1}{2}\dfrac{\partial \tilde{X}'_{ri}}{\partial q_k} 
 M_{rs} \tilde{X}'_{sh} + 
 \sum_{r,s}^n
 \tilde{X}'_{ri} M_{rs}  \dfrac{\partial\tilde{X}'_{sh}}{\partial q_k}  -\dfrac{1}{2}\sum_{r,s}^n \tilde{X}'_{rk} M_{rs} \dfrac{\partial\tilde{X}'_{si}}{\partial q_h}
\]
Riportando i termini nella sommatoria di indici $h,k$ e scambiando adeguatamente gli indici muti
\begin{align*}
&\sum_{h,k}^n \left(\dfrac{1}{2}\sum_{r,s}^n \dfrac{\partial \tilde{X}'_{ri}}{\partial q_k} 
 M_{rs} \tilde{X}'_{sh} -\dfrac{1}{2}\sum_{r,s}^n \tilde{X}'_{rk} M_{rs} \dfrac{\partial\tilde{X}'_{si}}{\partial q_h} \right) \dot{q}_h \dot{q}_h =\\
&=\sum_{h,k}^n \left(\dfrac{1}{2} \sum_{r,s}^n \dfrac{\partial \tilde{X}'_{si}}{\partial q_h} 
 M_{rs} \tilde{X}'_{rk} -\dfrac{1}{2}\sum_{r,s}^n \tilde{X}'_{rk} M_{rs} \dfrac{\partial\tilde{X}'_{si}}{\partial q_h} \right) \dot{q}_h \dot{q}_h =0
\end{align*}
e quindi  si ha
\[
\sum_{h,k=1}^n \left( \dfrac{\partial A_{ih}}{\partial q_k}  - \dfrac{1}{2} \dfrac{\partial A_{hk})}{\partial q_i} 
\right)\dot{q}_h\dot{q}_k =
\sum_{h,k=1}^n \left(
\sum_{r,s}^n
 \tilde{X}'_{ri} M_{rs}  \dfrac{\partial\tilde{X}'_{sh}}{\partial q_k}
\right)\dot{q}_h\dot{q}_k 
\]
che è quanto dovevamo mostrare. \qedhere
\end{dimo}

In qualunque sistema di coordinate locali le equazioni del moto vincolato da un vincolo olonomo fisso ideale sono nella forma \eqref{eqLagrangeFisso}, per scriverle in un caso particolare dovremo solo ricavare l'energia cinetica $T(q,\dot{q})$ e le componenti lagrangiane delle forze $Q_i$.

Le equazioni di Lagrange sono equivalenti alle condizioni \eqref{eqPrincipioDAlembert}, chiamate \textit{principio di D'Alembert}.

Comunemente le equazioni di Lagrange si scrivono nella forma simbolica
\[
\dfrac{d}{dt}\dfrac{\partial T}{\partial \dot{q}_i} - \dfrac{\partial T}{\partial q_i} = Q_i
\]
o anche
\[
\dfrac{d}{dt}\dfrac{\partial T}{\partial \dot{q}} - \dfrac{\partial T}{\partial q} = Q
\]
questa seconda forma è un po' imprecisa: le $\dfrac{\partial T}{\partial q}, \dfrac{\partial T}{\partial \dot{q}}$ sono matrici $1\times n$ cioè vettori riga mentre le $Q$ sono normali vettori colonna. Questo problema si può ignorare se non sono presenti prodotti righe per colonne ma fa emergere un cambiamento nella definizione di forza: le forze saranno vettori dello spazio duale a $\mathbb{V}^{3N}$ e non semplicemente vettori in $\mathbb{V}^{3N}$, cioè sono funzioni lineari sui vettori o anche \textit{covettori}. 
Quindi perchè l'equazione di Newton $M\ddot{X}=F$ abbia senso dovremo ammettere che la matrice $M$ trasformi vettori in covettori, cioè che sia una \textit{metrica}.


\subsection{Le equazioni di Lagrange per vincoli o coordinate dipendenti dal tempo}

Le equazioni di un moto vincolato da un vincolo mobile ideale con varietà vincolare $\{Q_t:t\in\mathbb{R}\}$ e forze attive $F$ sono le restrizioni agli atti di moto vincolati di $M\ddot{X}=F(X,\dot{X},t)+\Phi_{id}(X,\dot{X},t)$ dove $\Phi_{id}(X,V)\in (T_XQ_t)^\perp$ per ogni $(X,V)\in \mathcal{V}Q_t$ e $\forall t$.

\begin{proposizione}[Equazioni di Lagrange (mobile)]\label{propLagrangeMobile}
La restrizione delle equazioni di Lagrange di prima specie per vincoli ideali mobili allo spazio degli atti di moto vincolati è
\begin{equation}\label{eqLagrangeMobileo}
\left(\dfrac{d}{dt}\dfrac{\partial T}{\partial \dot{q}_i} \right) (q,\dot{q},\ddot{q},t) - \dfrac{\partial T}{\partial q_i}(q,\dot{q},t)=Q_i(q,\dot{q},t)
\end{equation}
Queste equazioni si possono porre in forma normale.
\end{proposizione}
Non diamo una dimostrazione precisa della \ref{propLagrangeMobile} però ne diamo un'idea: come per il caso fisso, partendo dalle equazioni di prima specie e scrivendole mediante la parametrizzazione
\[
M\dfrac{d\tilde{W}}{dt}(q,\dot{q},\ddot{q},t)=F\left(\tilde{X}(q,t),\tilde{W}(q,\dot{q},t),t\right)+\Phi_{id}\left(\tilde{X}(q,t),\tilde{W}(q,\dot{q},t),t\right)
\]
moltiplicando per i vettori della base naturale di $T_{\tilde{X}(q,t)}Q_t$ si trova
\[
\dfrac{\partial \tilde{X}}{\partial q_i}(q,t)\cdot M\dfrac{d\tilde{W}}{dt}(q,\dot{q},\ddot{q},t)=\dfrac{\partial \tilde{X}}{\partial q_i}(q,t)\cdot F\left(\tilde{X}(q,t),\tilde{W}(q,\dot{q},t),t\right)= Q_i\left(\tilde{X}(q,t),\tilde{W}(q,\dot{q},t),t\right)
\]

La dipendenza dal tempo farà comparire termini aggiuntivi e quindi l'espressione esplicita delle equazioni di Lagrange sarà diversa. Nello specifico
\[
\dfrac{d}{dt}\dfrac{\partial T}{\partial \dot{q}_i }= \sum_{j=1}^n \dfrac{\partial^2 T}{\partial \dot{q}_j \partial \dot{q}_i} \ddot{q}_j + 
\sum_{j=1}^n \dfrac{\partial^2 T}{\partial q_j \partial \dot{q}_i} \dot{q}_j +
\dfrac{\partial^2 T}{\partial t \partial \dot{q}_i} 
\]


%Parte 3--------------------------------------------
\cleardoublepage
\addcontentsline{toc}{part}{Meccanica Lagrangiana}
\thispagestyle{plain}
	\begin{center}
		\vspace{.1\textheight}
		{\LARGE\scshape \textsc{Terza parte}\par}
		\vspace{.1\textheight}
		\textsc{}\par
		\vspace{.05\textheight}
		{\Huge \textsc{Meccanica Lagrangiana}\par}
		\vfill

\end{center}
\fancyfoot[C]{\small Meccanica Lagrangiana}
%-----------------------------------------------------------
\section{Lagrangiana}
Se le componenti lagrangiane della forza si potessero esprimere come derivate di una funzione, un potenziale, con una scrittura in qualche modo simile a quella del primo membro delle equazioni di Lagrange potremo semplificare ulteriormente le equazioni del moto.

\begin{definizione}
Consideriamo un sistema olonomo con parametrizzazione, possibilmente dipendente dal tempo, $\tilde{X} : U \times \mathbb{R}\rightarrow Q$ sul quale agisce una
forza generalizzata $Q = (Q_1, . . . , Q_n)$. Si dice che:\begin{enumerate}[i.]
\item $Q$ è posizionale se non dipende da $\dot{q}$ e si potrà quindi riguardare come mappa $U\times\mathbb{R}\rightarrow\mathbb{R}^n$, $(q,t)\mapsto Q(q,t)$.
\item $Q$ ammette energia potenziale se è posizionale ed esiste una funzione $V:U\times\mathbb{R}\rightarrow\mathbb{R}$ $(q,t)\mapsto V(q,t)$ tale che
\[
Q_i (q,t) = -\dfrac{\partial V}{\partial q_i}(q,t) \qquad i=1,\dots n \qquad \forall q,t
\]
Chiameremo $V$ energia potenziale.
\end{enumerate}
Se $Q$ non dipende dal tempo ed ammette energia potenziale la chiameremo conservativa.
\end{definizione}

Naturalmente il potenziale di una data $Q$ non è unico nel senso che due potenziali $V_a$ e $V_b$ daranno le stesse $Q_i$ se la loro differenza è una costante additiva o una funzione non dipendente dalle $q_i$.

\begin{proposizione}
Le equazioni di Lagrange per un sistema olonomo con energia cinetica $T(q,\dot{q},t)$ soggetto a forze generalizzate $Q(q,t)$ posizionali che ammettono energia potenziale $V(q,t)$ si possono scrivere come
\begin{equation}
\dfrac{d}{dt}\dfrac{\partial L}{\partial \dot{q}_i} - \dfrac{\partial L}{\partial q_i}=0 \qquad i=1,\dots n
\end{equation}
dove $L(q,\dot{q},t)=T(q,\dot{q},t)-V(q,t)$ è detta Lagrangiana del sistema.
\end{proposizione}
\begin{dimo}
\[
\dfrac{d}{dt}\dfrac{\partial T}{\partial \dot{q}_i} - \dfrac{\partial T}{\partial q_i}+\dfrac{d}{dt}\dfrac{\partial V}{\partial \dot{q}_i} - \dfrac{\partial V}{\partial q_i}=
\dfrac{d}{dt}\dfrac{\partial T}{\partial \dot{q}_i} - \dfrac{\partial T}{\partial q_i}- \dfrac{\partial V}{\partial q_i}=0 \qedhere
\]
\end{dimo}

Potremo generalizzare il concetto di potenziale alle forze dipendenti dalla velocità, in modo che il potenziale per le forze posizionali sia un caso particolare di questo nuovo potenziale generalizzato.

\begin{definizione}[Potenziale generalizzato]
Si dice che la forza generalizzata $Q:U\times\mathbb{R}^n\times\mathbb{R}\rightarrow\mathbb{R}^n$, $(q,\dot{q},t)\mapsto Q(q,\dot{q},t)$ ammette un potenziale che dipende dalle velocità, o potenziale generalizzato, se esiste una funzione $V:U\times\mathbb{R}^n\times\mathbb{R}\rightarrow\mathbb{R}$ $(q,\dot{q},t)\rightarrow V(q,\dot{q},t)$ tale che
\[
Q_i(q,\dot{q},t)=\dfrac{d}{dt}\dfrac{\partial V}{\partial \dot{q}_i}(q,\dot{q},t) - \dfrac{\partial V}{\partial q_i}(q,\dot{q},t)
\]
\end{definizione}

Allora esplicitamente
\[
Q_i(q,\dot{q},t)=\sum_{j=1}^n \dfrac{\partial^2 V}{\partial \dot{q}_i\partial\dot{q}_j}\ddot{q}_j +\sum_{j=1}^n \dfrac{\partial^2 V}{\partial \dot{q}_i \partial q_j}\dot{q}_j + \dfrac{\partial^2 V}{\partial \dot{q}_i \partial t} - \dfrac{\partial V}{\partial q_i}
\]
però dato che le $Q_i$ non dipendono dalle $\ddot{q}$ dovremo avere
$ \dfrac{\partial^2 V}{\partial \dot{q}_i \partial \dot{q}_j}(q,\dot{q},t) =0 $ per $ i,j=1,\dots n$ $\forall q,\dot{q},t$
che significa che la matrice $\dfrac{\partial^2 V}{\partial \dot{q}\partial \dot{q}}$ dovrà essere nulla. 

\begin{proposizione}
Se $Q:U\times\mathbb{R}^n\times\mathbb{R}\rightarrow\mathbb{R}^n$ ammette un potenziale dipendente dalle velocità $V:U\times\mathbb{R}^n\times\mathbb{R}\rightarrow\mathbb{R}$ allora \begin{enumerate}[i.]
\item $V$ è affine, cioè lineare non omogeneo, nelle velocità
\[
V(q,\dot{q},t)=V_1(q,\dot{q},t)+V_0(q,t) \qquad \forall q,\dot{q},t
\]
con $V_0=(q,t)$ indipendente dalle $\dot{q}$ e $V_1(q,\dot{q},t)$ lineare in $\dot{q}$.
\item Se inoltre $V$ non dipende dal tempo
\[
Q(q,\dot{q})=-\nabla V_0(q) + Z(q)\dot{q} \qquad \forall q
\]
dove $Z(q)$ è una matrice $n\times n$ antisimmetrica.

\end{enumerate}
\end{proposizione}
\begin{dimo}~\begin{enumerate}[i.]
\item Poichè $\dfrac{\partial^2 V}{\partial \dot{q}\partial\dot{q}}=0$ la $\dfrac{\partial V}{\partial \dot{q}}$ è indipendente dalle $\dot{q}$ cioè $\dfrac{\partial V}{\partial \dot{q}}=\beta(q,t)$ per una qualche $\beta:U\times\mathbb{R}\rightarrow\mathbb{R}^n$.
Integrando in $\dot{q}$ si trova quindi
\[
V(q,\dot{q},t)=\beta(q,t)\dot{q}+ V_0(q,t)
\]
con $\beta$ e $V_0$ indipendenti da $\dot{q}$.
\item Se $V(q,\dot{q})=\beta(q)\dot{q}+ V_0(q)=V_0(q) + \sum_{j=1}^n \beta_j(q)\dot{q}_j $ allora 
\begin{align*}
Q_i (q,\dot{q}) & = \dfrac{d}{dt}\dfrac{\partial V}{\partial \dot{q}_i} - \dfrac{\partial V}{\partial q_i}
=\dfrac{d}{dt} \left(\dfrac{\partial}{\partial\dot{q}_i}\sum_{j=1}^n \beta_j\dot{q}_j \right) - \dfrac{\partial}{\partial q_i}\sum_{j=1}^n\beta_j\dot{q}_j -\dfrac{\partial V_0}{\partial q_i} =\\
& = \dfrac{d}{dt} \beta_i - \dfrac{\partial}{\partial q_i}\sum_{j=1}^n\beta_j\dot{q}_j -\dfrac{\partial V_0}{\partial q_i} 
= \sum_{j=1}^n \dfrac{\partial \beta_i}{\partial q_j} \dot{q}_j - \dfrac{\partial}{\partial q_i}\sum_{j=1}^n\beta_j\dot{q}_j -\dfrac{\partial V_0}{\partial q_i} = \\
&= \sum_{j=1}^n \left( \dfrac{\partial \beta_i}{\partial q_j} - \dfrac{\partial\beta_j}{\partial q_i}\right)\dot{q}_j -\dfrac{\partial V_0}{\partial q_i}
\end{align*}
con $\dfrac{\partial \beta_i}{\partial q_j}-\dfrac{\partial \beta_j}{\partial q_i}=Z_{ij}$. \qedhere
\end{enumerate}
\end{dimo}

Una forza, o una componente di una forza, $Z(q)\dot{q}$ con $Z(q)$ antisimmetrica $\forall q$ si chiama \textbf{forza girostatica}. Vedremo più avanti che le forze girostatiche non dissipano energia.

\begin{proposizione}
Dato un sistema olonomo di $N$ punti materiali con varietà vincolare $Q$, anche dipendente dal tempo, di dimensione $n\leq 3N$.
La forza attiva $F=(f_1,\dots f_N):\mathbb{R}^{3N}\times\mathbb{R}^{3N}\times\mathbb{R}\rightarrow\mathbb{R}^{3N}$ ammette un potenziale generalizzato se esiste una $\tilde{V}:\mathbb{R}^{3N}\times\mathbb{R}^{3N}\times\mathbb{R}\rightarrow\mathbb{R}^{3N}$ tale che
\[ f_k(X;\dot{X},t)= \dfrac{d}{dt}\dfrac{\partial \tilde{V}}{\partial \dot{X}_k} (X,\dot{X},t) - \dfrac{\partial \tilde{V}}{\partial X_k}(X;\dot{X},t)
\]
Allora se il sistema è soggetto a vincoli ideali le componenti lagrangiane della forza $Q_i$ ammettono il potenziale generalizzato
\[
V(q,\dot{q},t):=\tilde{V}(\tilde{X}(q,t),\tilde{W}(q,\dot{q},t),t)
\]
\end{proposizione}

\begin{proposizione}
Se le componenti lagrangiane della forza $Q_i(q,\dot{q},t)$ ammettono un potenziale generalizzato $V(q,\dot{q},t)$ allora per ogni funzione $G(q,t)$ ammettono come potenziale generalizzato anche
\[ V(q,\dot{q},t)+\dfrac{dG}{dt}(q,\dot{q},t)\]
\end{proposizione}
\begin{dimo}
Vogliamo avere
\[
\dfrac{d}{dt}\left(\dfrac{\partial}{\partial \dot{q}_i}\dfrac{dG}{dt}(q,\dot{q},t) \right) - \dfrac{\partial}{\partial q_i}\dfrac{dG}{dt}=0
\]
La derivata totale è
\[
\dfrac{d}{dt}G(q,t)=\sum_{i=1}^n \dfrac {\partial G}{q_i}(q,t) \dot{q}_i + \dfrac{\partial G}{\partial t}(q,t) 
\]
da cui $\dfrac{\partial}{\partial \dot{q}_i}\dfrac{d G}{dt}=\dfrac{\partial G}{\partial q_i}$ e quindi
$ \dfrac{d}{dt}\left(\dfrac{\partial}{\partial \dot{q}_i}\dfrac{dG}{dt}(q,\dot{q},t) \right)=
\dfrac{d}{dt}\left( \dfrac{\partial G}{\partial q_i} (q,t) \right)$ e
\[
\dfrac{d}{dt}\left( \dfrac{\partial G}{\partial q_i} (q,t) \right)=
\sum_{j=1}^n \dfrac{\partial^2 G}{\partial q_i\partial q_j}(q,t)\dot{q}_j+\dfrac{\partial^2 G}{\partial q_i\partial t}=\dfrac{\partial}{\partial q_i}\dfrac{dG}{dt}(q,\dot{q},t) \qedhere
\]
\end{dimo}


\begin{esempio}[Pendolo rotante]
Un punto materiale soggetto a forza peso e vincolato ad un cerchio che ruota con velocità angolare $\omega$ costante attorno al suo asse verticale costituisce un pendolo rotante. 
Possiamo considerare due sistemi di riferimento: uno solidale al cerchio in cui la guida circolare è fissa nel piano $(x_1,x_3)$ e centrata nell'origine e un secondo (che diremmo \textit{solidale al laboratorio}) nel quale la guida ruota attorno all'asse verticale. Il primo riferimento avrà velocità angolare $\omega$ costante rispetto al secondo.
\begin{itemize}[label=$\cdot$]\item
Nel riferimento solidale al cerchio il vincolo è fisso $Q=C_R(O)$ ma sul punto materiale agiscono, oltre alla forza peso, la forza centrifuga e la forza di Coriolis. Parametrizziamo la varietà con $\tilde{x}(\theta)=(R\cos{\theta},0,R\sin{\theta})$ e l'energia cinetica del sistema vincolato sarà
\[
T(\theta,\dot{\theta})=\dfrac{1}{2}mR^2\dot{\theta}^2
\]
L'energia potenziale della forza peso è $V(\theta)=\left.mgx_3\right|_{x_3=R\sin{\theta}}=mgR\sin{\theta}$. La forza centripeta e il suo potenziale sono
\[
F_{c}(x,\dot{x})=m \omega^2 x e_1
\qquad
\tilde{V}_{c}(x,\dot{x})=-\dfrac{1}{2}m\omega^2x_1^2\longrightarrow V_{c}(\theta,\dot{\theta})=-\dfrac{1}{2}m\omega^2R^2\cos^2{\theta}
\]
mentre la forza di Coriolis è ortogonale al piano e quindi ha componente lagrangiana nulla e non influisce sul moto lungo la guida (si può calcolare il potenziale generalizzato e poi verificare che è nullo sugli atti di moto vincolati).

Complessivamente il potenziale  è \[V(\theta)=mgR\sin{\theta}-\dfrac{1}{2}m\omega^2R^2\cos^2{\theta}\] e la Lagrangiana del sistema vincolato $L=T_2-V_0$ è
\[
L(\theta,\dot{\theta})=\dfrac{1}{2}mR^2\dot{\theta}^2-mgR\sin{\theta}+\dfrac{1}{2}m\omega^2R^2\cos^2{\theta}
\]
\item
Nel riferimento che vede la guida ruotare l'unica forza attiva agente sul punto materiale è la forza peso ma il vincolo è mobile. 
Scegliamo il riferimento in modo che la guida abbia centro nell'origine e che al tempo $t=0$ sia nel piano $(x,z)$.La parametrizzazione dipendente dal tempo sarà
\[
\tilde{x}(\theta,t)=(R\cos{\theta}\cos{\omega t},R\cos{\theta}\sin{\omega t},R\sin{\theta})
\]
L'energia cinetica sarà
\[
T(\theta,\dot{\theta},t)=\dfrac{1}{2}mR^2\left( \dot{\theta}^2+\omega^2\cos^2{\theta} \right)=T_2+T_0
\]
L'energia potenziale della forza peso è $V(\theta)=mgR\sin{\theta}$ e quindi la Lagrangiana $L(\theta,\dot{\theta},t)=T_2+T_0-V_0$ è
\[
L(\theta,\dot{\theta})=\dfrac{1}{2}mR^2\dot{\theta}^2-mgR\sin{\theta}+\dfrac{1}{2}m\omega^2R^2\cos^2{\theta}
\]
ed è la stessa trovata nell'altro punto ma con una diversa corrispondenza per il termine centrifugo.
\begin{figure}[H]\begin{center}
\includegraphics[width=0.45\linewidth]{img/pendoloRotante.png}
\end{center}\end{figure}
\end{itemize}\end{esempio}


\section{Sistemi Lagrangiani}
\begin{definizione}[Sistema Lagrangiano]
Sia $U\subseteq\mathbb{R}^n$ un aperto e $L:U\times\mathbb{R}^n\times R\rightarrow\mathbb{R}$ una funzione $(q,\dot{q},t)=L(q,\dot{q},t)$. Il sistema di equazioni
\begin{equation}\label{eqLagrangeLagrangiana}
\left(\dfrac{d}{dt}\dfrac{\partial L}{\partial \dot{q}_i}\right) (q,\dot{q},\ddot{q},t)-\dfrac{\partial L}{\partial q_i}(q,\dot{q},t)=0 \qquad i=1,\dots,n
\end{equation}
in $(q,\dot{q})\in U\times\mathbb{R}^n$ si chiama sistema lagrangiano di Lagrangiana $L$ o anche equazione di Lagrange di Lagrangiana $L$. 
La dimensione $n$ di $U$ è il numero di gradi di libertà del sistema.
\end{definizione}

\begin{importante}
Le equazioni del moto di un sistema soggetto a vincoli olonomi ideali con energia cinetica $T(q,\dot{q},t)$ e forze attive che ammettono un potenziale generalizzato $V$ si possono dare con le \eqref{eqLagrangeLagrangiana} in
\begin{equation}\label{eqLagrangianaMecc}
L=T(q,\dot{q},t)-V(q,\dot{q},t)
\end{equation}
ponendo
\begin{equation}\label{eqLagrangianaMeccCinPot}
T(q,\dot{q},t)=T_2(q,\dot{q},t)+T_1(q,\dot{q},t)+T_0(q,t) 
\qquad
V(q,\dot{q},t)=V_1(q,\dot{q},t) + V_0(q,\dot{q},t)
\end{equation}
dove $T_2$ è quadratico nelle $\dot{q}$, $T_1$ e $V_1$ sono lineari nelle $\dot{q}$ e $T_0$ e $V_0$ non dipendono dalle $\dot{q}$.
\end{importante}

Chiameremo \textbf{sistema lagrangiano di tipo meccanico} un sistema in cui la lagrangiana è data dalle \eqref{eqLagrangianaMecc} \eqref{eqLagrangianaMeccCinPot}.
I sistemi lagrangiani \eqref{eqLagrangeLagrangiana} ma non di tipo meccanico sono molto importanti in fisica teorica e ne studieremo alcuni esempi più avanti. Tratteremo il più possibile in modo indipendente dalla forma della lagrangiana.


Le equazioni di Lagrange \eqref{eqLagrangeLagrangiana} si possono esplicitare
\[
\sum_{j=1}^n \dfrac{\partial^2 L}{\partial \dot{q}_i \partial \dot{q_j}}(q,\dot{q},t)\, \ddot{q}_j + \sum_{j=1}^n \dfrac{\partial^2 L}{\partial \dot{q}_i \partial q_j}(q,\dot{q},t)\, \dot{q}_j + \dfrac{\partial^2 L}{\partial \dot{q}_i\partial t}(q,\dot{q},t) - \dfrac{\partial L}{\partial q_i}(q,\dot{q},t)=0
\]
ovvero
\begin{equation}\label{eqLagrangeLagrangianaEsplicite}
\left(\dfrac{\partial^2 L}{\partial \dot{q} \partial\dot{q}}(q,\dot{q},t)\, \ddot{q}\right)_i =
\dfrac{\partial L}{\partial q_i}(q,\dot{q},t) - \sum_{j=1}^n \dfrac{\partial^2 L}{\partial \dot{q}_i \partial q_j}(q,\dot{q},t)\, \dot{q}_j - \dfrac{\partial^2 L}{\partial \dot{q}_i\partial t}(q,\dot{q},t) 
\end{equation}
quindi sono equazioni differenziali del secondo ordine per le $q$ e si potranno porre in forma normale se la matrice $\dfrac{\partial^2 L}{\partial \dot{q} \partial\dot{q}}(q,\dot{q},t)$ è invertibile per ogni $(q,\dot{q},t)$ cioè se il suo determinante è diverso da zero: una Lagrangiana di questo tipo si dirà \textbf{regolare}.
Le Lagrangiane meccaniche sono tutte regolari per costruzione.

In questa sezione tratteremo soltanto le Lagrangiane regolari.

\subsection{Integrale di Jacobi}
\begin{proposizione}
Per un Lagrangiana indipendente dal tempo $L(q,\dot{q})$ la funzione
\[
E(q,\dot{q}):=\sum_{i=1}^n \dot{q}_i\dfrac{\partial L}{\partial \dot{q}_i}(q,\dot{q}) -L(q,\dot{q})
\]
è un integrale primo delle equazioni di Lagrange di Lagrangiana $L$.
\end{proposizione}
\begin{dimo}
Consideriamo una curva integrale $t\mapsto q_t$ delle $\dfrac{d}{dt}\dfrac{\partial L}{\partial\dot{q}_i}-\dfrac{\partial L}{\partial q_i}=0$ e deriviamo la funzione composta $t\mapsto E(q_t,\dot{q}_t)$ rispetto al tempo (derivata lungo una curva uguale alla derivata totale):
\begin{align*}
\dfrac{\partial E(q_t,\dot{q}_t)}{\partial t}&=
\sum_{i=1}^n\left(
\ddot{q}_{i,t}\dfrac{\partial L(q_t,\dot{q}_t)}{\partial \dot{q}_{i,t}}+\dot{q}_{i,t}\dfrac{\partial}{\partial t}\dfrac{\partial L(q_t,\dot{q}_t)}{\partial\dot{q}_{i,t}}
-\dfrac{\partial L(q_t,\dot{q}_t)}{\partial q_{i,t}}\dot{q}_{i,t}-\dfrac{\partial L(q_t,\dot{q}_t)}{\partial \dot{q}_{i,t}}\ddot{q}_{i,t}
\right)=\\
&=\dfrac{d}{dt}E(q,\dot{q})=\sum_{i=1}^n \dot{q}_i\left(\dfrac{d}{dt}\dfrac{\partial L}{\partial \dot{q}_i}(q,\dot{q})-\dfrac{\partial L}{\partial q_i}(q,\dot{q})
\right)=0
\end{align*}
ed essendo costante lungo una soluzione qualsiasi $E$ è un integrale primo.\qedhere
\end{dimo}
Ricordiamo alcuni fatti sulle funzioni omogenee:\begin{itemize}
\item Una funzione $f:\mathbb{R}^n\setminus\{0\}\rightarrow\mathbb{R}$ si dice \textbf{omogenea di grado p} se
\begin{equation}
f(\lambda x)=\lambda^pf(x)
\end{equation}
\item Il \textbf{Teorema di Euler} dice che una $f:\mathbb{R}^n\setminus\{0\}\rightarrow\mathbb{R}$ è omogenea di grado $p$ se e solo se
\begin{equation}
\sum_{i=1}^n x_i\dfrac{\partial f}{\partial x_i}(c)=pf(x) \qquad \forall x
\end{equation}
\end{itemize}
Notiamo che nelle lagrangiane meccaniche $L=T_2+T_1+T_0-V_1-V_0$ i pedici indicano il grado di omogeneità rispetto alla $\dot{q}$ dei termini.\par
Per sistemi meccanici di lagrangiana $L$ potremo allora usare le funzioni omogenee per esplicitare l'integrale di Jacobi
\begin{align*}
E(q,\dot{q})&=\sum_{i=1}^n \dot{q}_i \dfrac{\partial L}{\partial \dot{q}_i}(q,\dot{q}) -L(q,\dot{q})=\\
& = \sum_{i=1}^n \dot{q}_i \dfrac{\partial T_2}{\partial \dot{q}_i}(q,\dot{q}) + \dot{q}_i \dfrac{\partial T_1}{\partial \dot{q}_i}(q,\dot{q}) - \dot{q}_i \dfrac{\partial V_1}{\partial \dot{q}_i}(q,\dot{q}) - \left( T_2+T_1+T_0-V_1-V_0\right) = \\
& = 2T_2 + T_1 - V_1 - \left( T_2+T_1+T_0-V_1-V_0\right) 
\end{align*}
e quindi
\begin{equation}
E_{mecc}(q,\dot{q})=T_2(q,\dot{q})-T_0(q)+V_0(q)
\end{equation}
in particolare se l'energia cinetica è $T=T_2$ e le forze ammettono energia potenziale cioè $V=V_0$, come nel caso di vincoli e parametrizzazioni fissi e forze conservative posizionali, l'integrale di Jacobi sarà $E=T_2+V_0=T+V$ cioè sarà l'\textbf{energia totale del sistema}, de
finita come la somma di energia cinetica e potenziale.

\begin{esempio}[Integrale di Jacobi del pendolo rotante]
Abbiamo già visto che un pendolo rotante si può con una Lagrangiana 
\[L(\theta,\dot{\theta})=\dfrac{1}{2}mR^2\dot{\theta}^2+\dfrac{1}{2}m\omega^2R^2\cos^2{\theta}-mgr\sin{\theta}=T_2-V_0\]
con un riferimento mobile.

Vediamo che $\dfrac{\partial L}{\partial t}=0$, $T=T_2$ e la forza peso è posizionale e conservativa quindi l'integrale di Jacobi sarà
\[
E(\theta,\dot{\theta}=T_2+V_0=\dfrac{1}{2}mR^2\dot{\theta}^2-\dfrac{1}{2}m\omega^2R^2\cos^2{\theta}+mgr\sin{\theta}
\]
\end{esempio}



\subsection{Momenti coniugati}
\begin{definizione}
Dato un sistema lagrangiano di Lagrangiana $L(q,\dot{q},t)$ chiamiamo momenti collegati alle coordinate lagrangiane le funzioni
\[
p_i(q,\dot{q},t):=\dfrac{\partial L}{\partial \dot{q}_i}(q,\dot{q},t) \qquad i=1,\dots, n
\]
\end{definizione}
Data una soluzione $t\mapsto q_t$ delle equazioni lagrangiane di Lagrangiana $L$ vediamo che lungo la soluzione
\[
\dfrac{\partial p_i(q_t,\dot{q}_t,t)}{\partial t}=\sum_{j=1}^n\dfrac{\partial^2 L(q_t,\dot{q}_t,t)}{\partial \dot{q}_{i,t} \partial q_{j,t}}\dot{q}_{j,t}+\dfrac{\partial^2 L(q_t,\dot{q}_t,t)}{\partial \dot{q}_{i,t} \partial \dot{q}_{j,t}}\ddot{q}_{j,t} + \dfrac{\partial^2 L(q_t,\dot{q}_t,t)}{\partial \dot{q}_{i,t} \partial t}
\]
e ricordando l'espressione esplicita \eqref{eqLagrangeLagrangianaEsplicite} ed usando la derivata totale avremo
\[
\dfrac{d}{dt}p_i(q,\dot{q},t)=\dfrac{\partial L}{\partial q_i} (q,\dot{q},t)
\]
\begin{proposizione}
Il momento coniugato alla coordinata $q_i$ è un integrale primo delle equazioni di Lagrange di Lagrangiana $L$ se e solo se \[\dfrac{\partial L}{\partial q_i}(q,\dot{q},t)=0\] cioè se la Lagrangiana non dipende dalla coordinata $q_i$.

In tal caso diremo che $q_i$ è ignorabile o ciclica.
\end{proposizione}

La presenza di coordinate ignorabili dipende dalla scelta del sistema di riferimento e ci permetterà di semplificare l'integrazione del sistema.

\begin{esempio}[Momento di una particella libera]
Per una particella libera soggetta a forze che ammettono un potenziale $V(x)$ la lagrangiana sarà
\[
L(x,\dot{x})=\dfrac{1}{2}m\left|\left|\dot{x}\right|\right|^2-V(x)
\]
quindi il momento sarà
\[
p_i=\dfrac{\partial L(x,\dot{x})}{\partial \dot{x}_i}=m\dot{x}_i
\]
Consideriamo ora un potenziale centrale $V(x)=V(\left|\left|x\right|\right|)$ come quello elettrico o gravitazionale: Non sembra che $p_i$ sia un integrale prima ma studiamo in coordinate sferiche $\tilde{X}(\rho,\theta,\varphi)=(\rho\sin{\theta}\cos{\varphi},\rho\sin{\theta}\sin{\varphi},\rho\cos{\theta})$:
\[
L(\rho,\theta,\varphi,\dot{\rho},\dot{\theta},\dot{\varphi})=\dfrac{1}{2}m\left(\dot{\rho}^2+\dot{\theta}^2\rho^2+\dot{\varphi}^2\sin^2{\theta} \right)-\tilde{V}(\rho)
\]
e quindi dato che $L$ non dipende da $\varphi$ la $p_\varphi$ sarà un integrale primo
\[
p_\varphi=m\rho^2\sin^2{\theta}\,\dot{\varphi}
\]
\end{esempio}




\subsection{Invarianza in forma delle equazioni di Lagrange per cambi di coordinate}
Le equazioni di Lagrange sono equazioni del secondo ordine nella variabile $q$, dato un cambiamento di coordinate $C$ nello spazio delle configurazioni il suo sollevamento tangente $TC$ coniuga le equazioni di Lagrange con altre equazioni del secondo ordine.

\begin{proposizione}
Sia $C : U \rightarrow \tilde{U}$ un diffeomorfismo fra due aperti di $\mathbb{R}^n$.
Il sollevamento tangente $TC$ coniuga le equazioni di Lagrange di Lagrangiana $L$ su $TU = U\times\mathbb{R}^n$ alle equazioni di Lagrange di Lagrangiana $(TC)_* L = L \circ (TC)^{-1}$ su $T \tilde{U} = \tilde{U}\times\mathbb{R}^n$.
\end{proposizione}

\begin{dimo}
\end{dimo}

Per sistemi meccanici questo è abbastanza ovvio:\\
in una qualunque parametrizzazione $\tilde{X} : U\rightarrow\tilde{U}$ della varietà delle configurazioni Q, le equazioni del moto del sistema sono le equazioni di Lagrange di Lagrangiana $L$ che è la restrizione a $T\tilde{U}$ della Lagrangiana $\tilde{L}$ del sistema libero scritta in quella parametrizzazione $L = \tilde{L}\circ  T\tilde{X}$. Un cambio di coordinate $C: U\rightarrow \hat{U}$ produce una nuova parametrizzazione $\hat{X}=\tilde{X}\circ C^{-1}$ per $Q$ ed in essa le equazioni del moto sono le equazioni di Lagrange per la Lagrangiana $\tilde{L}\circ T(\tilde{X}\circ C^{-1})=\tilde{L}\circ T\tilde{X}\circ (TC)^{-1}=L\circ(TC)^-1=(TC)_*L$.

\section{Equilibri nei sistemi lagrangiani}
Le equazioni di Lagrange sono equazioni differenziali del secondo ordine e quindi i loro punti di equilibri sono $(q^*,0)$ dove chiameremo $q^*$ configurazione di equilibrio.

Per $L=T_2+T_1+T_0-V_1-V_0$ possiamo inglobare $T_1$ in $V_1$ e $T_0$ in $V_0$ e intenderla come una $L=T_2-V_1-V_0$.

Chiamiamo \textbf{regolare} una Lagrangiana $L$ per cui $ \det{ \dfrac{\partial^2 L (q,\dot{q})}{\partial \dot{q} \partial \dot{q} }  } \neq 0 \>\> \forall (q,\dot{q})$

\begin{proposizione}
Per una Lagrangiana $L:TU\rightarrow\mathbb{R}$ regolare e indipendente dal tempo. Un punto $q^*$ è una configurazione di equilibrio per le equazioni di Lagrange se e solo se 
\[
\dfrac{\partial L}{\partial q}(q^*,0)=0
\]
\end{proposizione}

\begin{dimo}Se la Lagrangiana è indipendente dal tempo le equazioni di Lagrange diventano
\[
\dfrac{\partial^2 L}{\partial \dot{q}_i\partial \dot{q}_j}\ddot{q}_j + \dfrac{\partial^2 L}{\partial\dot{q}_i\partial q_j}\dot{q}_j - \dfrac{\partial L}{\partial q_i}=0 \qquad\quad \left( \dfrac{\partial^2 L}{\partial \dot{q}\partial \dot{q}}\ddot{q} \right)_i =   \dfrac{\partial L}{\partial q_i} - \dfrac{\partial^2 L}{\partial\dot{q}_i\partial q_j}\dot{q}_j 
\]
La matrice $\dfrac{\partial^2 L}{\partial \dot{q}\partial \dot{q}}\ddot{q}$ è invertibile quindi gli $(\overline{q},0)$ zeri del secondo membro sono gli unici punti di equilibrio.
Chiaramente in $(q^*,0)$ il secondo membro si annulla quindi $q^*$ è una configurazione di equilibrio.
\end{dimo}

\begin{corollario}\label{corEquilibrioLagrangeCriticoV0}
Per Lagrangiane meccaniche indipendenti dal tempo $L=T_2-V_1-V_0$ una configurazione $q^*$ è una configurazione di equilibrio se e solo se è un punto critico di $V_0$.
\end{corollario}
\begin{dimo}
$T_2$ e $V_1$ sono omogenee in $\dot{q}$ e quindi 
\[
\dfrac{\partial L}{\partial q} (q,0) = \dfrac{\partial T_2}{\partial q}(q,0)-\dfrac{\partial V_1}{\partial q}(q,0)-\dfrac{\partial V_0}{\partial q}(q,0)=-\dfrac{\partial V_0}{\partial q}(q,0) \quad  \qedhere
\]
\end{dimo}

\begin{corollario}\label{corEquilibriLagrangianaConfig}
Consideriamo un sistema meccanico costituito da $N$ punti materiali soggetto a vincoli olonomi fissi ideali con forze attive $F(X,\dot{X})$ che ammettono un potenziale $\tilde{V}_0(X)+\tilde{V}_1(X,\dot{X})$. Una configurazione nella varietà vincolare $X^*\in Q$ è una configurazione del sistema se e solo se
\[
\nabla \tilde{V}_0 (X^*) \perp T_{X^*}Q
\]
\end{corollario}

\begin{dimo}

\end{dimo}

\begin{corollario}Nella notazione del corollario \ref{corEquilibriLagrangianaConfig} i punti di minimo e  massimo di $\left.\tilde{V}_0\right|_Q$ sono configurazioni di equilibrio del sistema
\end{corollario}
\begin{dimo}
\end{dimo}

\subsection{Stabilità}
Vogliamo studiare la stabilità alla Lypaunov degli equilibri di Lagrangiane di tipo meccanico indipendenti dal tempo. L'esistenza dell'integrale di Jacobi rende impossibile la stabilità asintotica e la stabilità di un equilibri necessariamente per tutti i tempi. 

\begin{proposizione}[Teorema di Lagrange-Dirichlet]\label{teoLagrangeDirichlet}
Sia $L=T_2-V_1-V_0$ indipendente dal tempo. Ogni punto di minimo stretto di $V_0$ è una configurazione di equilibrio stabile per tutti i tempi.
\end{proposizione}
\begin{dimo}
Dalla proposizione \ref{corEquilibrioLagrangeCriticoV0} sappiamo che un $q^*$ minimo di $V_0$ è una configurazione di equilibrio delle equazioni di Lagrange.
Usiamo il teorema di Lyaounov usando l'integrale di Jacobi $E=T_2+V_0$ come funzione di Lyapunov. Ci basterà mostrare che $E$ ha minimo stretto in $(q^*,0)$. Ma $A(q)$ è definita positiva e $t\mapsto\sfrac{1}{2} \,\dot{q}\cdot A(q)\dot{q}$ ed è nulla per $\dot{q}=0$ mentre per ipotesi $V_0$ ha minimo stretto in $q^*$ quindi $(q^*,0)$ è un minimo stretto di $E$.\qedhere
\end{dimo}

Dobbiamo restringerci alle Lagrangiane $L=T_2-V_0$ per ottenere risultati significativi per l'instabilità.

\begin{proposizione}\label{propInstabilitàEquilibriLagrange}
Sia $L=T_2-V_0$ indipendente dal tempo. Sia $q^*$ un punto critico di $V_0$ allora se l'Hessiana \[\dfrac{\partial^2V_0}{\partial q\partial q}(q^*)\] ha qualche autovalore negativo $q^*$ è una configurazione di equilibrio instabile.
\end{proposizione}
La dimostrazione di questo risultato richiede di linearizzare le equazioni di Lagrange e quindi la vedremo nella prossima sezione.

\section{Linearizzazione delle equazioni di Lagrange}
Consideriamo Lagrangiana $L=T_2-V_0$ con equilibri $(q^*,0)$.

\begin{proposizione}
La linearizzazione delle equazioni di Lagrange di Lagrangiana $L=T_2-V_0$ attorno ad un loro equilibrio $(q^*,0)$ sono equazioni di Lagrange di Lagrangiana quadratizzata $L^*(q,\dot{q})=\dfrac{1}{2}\dot{q}\cdot A(q^*)\dot{q} - \dfrac{1}{2} (q-q^*)\cdot V''(q^*)(q-q^*)$ cioè il sistema
\[
A(q*)\,\dot{q}+V''(q^*)(q-q^*)=0
\]
\end{proposizione}
\begin{dimo}
Le equazioni di Lagrange con $L=T_v-V_0$ sono equazioni del secondo ordine del tipo $\dot{q}=Y(q,\dot{q})$ e quindi possiamo linearizzarle come visto nella sezione \ref{sezLinearizzazione}. 
\[
\dfrac{d}{dt}\dfrac{\partial L}{\partial \dot{q}}-\dfrac{\partial L}{\partial q}=A\,\ddot{q}-b(q,\dot{q})-V_0'(q)=0 \implies \ddot{q}=A^{-1}(q)\left[ b(q,\dot{q}) + V_0'(q)\right]
\]
dove $b(q,\dot{q})$ è il termine derivante dalle derivate in $q$ di $A$ che vediamo essere di secondo grado in $\dot{q}$ e di cui non ci interessa la scrittura precisa. 

Per sviluppare $Y(q,\dot{q})$ in serie di Taylor attorno a $(q^*,q)$ vediamo innanzitutto che
\[
b(q,\dot{q})=\mathcal{O}_2 \qquad V_0'(q)=V_0'(q^*)+V_0''(q^*)(q-q^*)+\dots=V_0''(q^*)(q-q^*)+\mathcal{O}_2
\]
La matrice cinetica è $A(q)=A(q^*)+A'(q^*)(q-q^*)+\dots=A(q^*)+\mathcal{O}_1$ ed essendo invertibile con componenti differenziabili $A^{-1}(q)=A^{-1}(q^*)+\mathcal{O}_1$.
Quindi
\[
Y(q,\dot{q})=-\left(A^{-1}(q^*)+\mathcal{O}_1)\right) \left[ V_0''(q^*)(q-q^*)+\mathcal{O}_2 \right]
\] e le equazioni di Lagrange linearizzate attorno all'equilibrio sono
\begin{equation}
A(q^*)\ddot{q}=V_0''(q^*)(q-q^*)=0 \qedhere
\end{equation}
\end{dimo}

Il sistema linearizzato attorno all'equilibrio per le equazioni di Lagrange sarà
\begin{equation}
\left(\begin{array}{c}
\dot{q}\\
\dot{v}
\end{array}\right)=\left(\begin{array}{cc}
\mathbb{O} & \mathbb{I} \\
-A^{-1}(q^*)\,V_0''(q^*) & \mathbb{O}
\end{array}\right)\left(\begin{array}{c}
q-q^*\\
v
\end{array}\right)
\end{equation}

Indichiamo con $\Lambda$ la matrice a blocchi del sistema linearizzato e per alleggerire la notazione da qua in avanti scriveremo semplicemente $A$ e $V''$ al posto di $A(q^*)$ e $V_0''(q^*)$.

Chiamiamo \textbf{spettro} di una matrice quadrata l'insieme dei suoi autovalori, ciascuno ripetuto tante volte quanto la sua molteplicità.

Ricordiamo per $B\in M_{n}$ con $\text{sp}M=\{\mu_1,\dots,\mu_n\}$ il polinomio caratteristico di $B$ sarà \[\det{B-\mu\mathbb{I}}=(-1)^n\prod_{i=1}^n (\mu-\mu_i)\]


\begin{lemma}\label{lemmaSpettroLambda}
Per $\text{sp}(A^{-1}V'')=\{\lambda_1,\dots,\lambda_n\}$ vale
\[\text{sp}\Lambda=\{\sqrt{-\lambda_1},-\sqrt{-\lambda_1},\dots,\sqrt{-\lambda_n},-\sqrt{-\lambda_n}\}
\]
\end{lemma}
\begin{dimo}
Chiamiamo $B=-A^{-1}V''$ 
\[ \text{sp}B=\left\{-\lambda_1,\dots,-\lambda_n\right\}\qquad \Lambda=\left(\begin{array}{cc}\mathbb{O} & \mathbb{I} \\ B & \mathbb{O} \end{array}\right)
\]

Allora $\det{B-\lambda\mathbb{I}}=(-1)^n\prod_{i=1}^n (\lambda+\lambda_i)$ e sfruttando le proprietà del determinante per le operazioni su matrici
\begin{align*}
\det{\Lambda-\mu\mathbb{I}}&=\det{\left(\begin{array}{cc} -\mu \mathbb{I} & \mathbb{I} \\ B & -\mu\mathbb{I} \end{array}\right)} = \det{\left(\begin{array}{cc} -\mu \mathbb{I}+\mu\mathbb{I} & \mathbb{I} \\ B-\mu^2\mathbb{I} & -\mu\mathbb{I} \end{array}\right)} = (-1)^n\det{\left(\begin{array}{cc} \mathbb{I} & \mathbb{O} \\  -\mu\mathbb{I} & B-\mu^2\mathbb{I} \end{array}\right)} = \\
&= (-1)^n \det{\mathbb{I}}\det{B-\mu^2\mathbb{I}}=\prod_{i=1}^n (\mu^2+\lambda_i)
\end{align*}
e lo spettro è dato dagli zeri quindi $\text{sp}\Lambda=\left\{\pm\sqrt{-\lambda_1},\dots,\pm\sqrt{-\lambda_n}\right\}$\qedhere
\end{dimo}

Per poter dimostrare la \ref{propInstabilitàEquilibriLagrange} dobbiamo ancora trovare un modo per collegare lo spettro di $V''$ allo spettro di $B=-A^{-1}V''$.

Vediamo innanzitutto che $A^{-1}V'' u = \lambda u$ equivale a $V'' u = \lambda A u$ dunque possiamo ricavare gli autovalori di $A^{-1}V''$ come le soluzioni di
\begin{equation}
\det(V''-\lambda A)=0
\end{equation}
Chiameremo, prendendoci un po' di libertà nella notazione, gli autovalori di $A^{-1}V''$ gli \textbf{A-autovalori} di $V''$, allo stesso modo gli autovettori di $A^{-1}V''$ saranno gli \textbf{A-autovettori} di $V''$ cioè i vettori $u$ per cui $V''u_j=\lambda_j Au_j$.\par
La normale equazione agli autovettori di $V''$ è $V'' u=\lambda \mathbb{I}u$ con corrispondenti autovalori le soluzioni di $\det(V''-\lambda \mathbb{I})=0$, queste equazioni sono in tutto simili a quelle per gli $A$-autovettori e $A$-autovalori ed addirittura, essendo $A$ simmetrica e definita positiva, potremmo ripetere la dimostrazione del teorema spettrale per gli $A$-autovalori senza grosse differenze. 

\begin{importante}
Se $A$ è una matrice simmetrica e definita positiva e $V''$ è una matrice simmetrica gli $A$-autovalori di $V''$ sono reali ed esiste una base di $\mathbb{R}^n$ composta da $A$-autovettori linearmente indipendenti ed $A$-ortonormali di $V''$
\[
\{
u_1,\dots,u_n
\}\qquad \quad
u_i\cdot Au_j=\delta_{ij}
\]
\end{importante}

\begin{dimo}[Dimostrazione proposizione \ref{propInstabilitàEquilibriLagrange}]
Per ipotesi $V''$ ha almeno un autovalore negativo, chiamiamolo $k$ con autovettore associato $v$. Allora $v\cdot V''v=k\left|\left|v\right|\right|^2<0$

Consideriamo una base $\{ u_1,\dots,u_n\}$ di $A$-autovettori di $V''$ 
\[
v=\sum_{i=1}^n c_i u_i  \qquad u_j\cdot Au_j=\delta_{ij} \qquad V'' u_j = \lambda_j Au_j
\]
Allora
\begin{align*}
0> v\cdot V''u_j &= v\cdot V''\sum_{i=1}^n c_i u_i = \sum_{i,j=1}^n c_jc_i u_j\cdot V'' u_i \\
&= \sum_{i,j=1}^n c_jc_i u_j\lambda_j A u_i = \sum_{i,j=1}^n c_jc_i\lambda_j u_j A u_i=\\
&=\sum_{i,j=1}^n c_jc_i\lambda_j \delta_{ij}= \sum_{i=1}^n c_i^2\lambda_i<0
\end{align*}
quindi esiste almeno un $\lambda_j<0$. Ma allora nello spettro di $\Lambda$ l'autovalore $-\sqrt{-\lambda}$ è reale negativo.

Per il teorema di Lyapunov allora l'equilibrio è instabile. \qedhere
\end{dimo}

\begin{esempio}[Pendolo doppio]
In questo esempio $L=T_2-V_0$ perchè il vincolo è fisso e la forza attiva è posizionale.
\begin{minipage}[H]{0.7\linewidth}
La varietà vincolare per il pendolo doppio è il prodotto di due cerchi $Q=S^1\times S^1$, la parametrizziamo con gli angoli $\varphi_1,\varphi_2$ tra le direzioni dei pendoli e la verticale.
\[
\tilde{X}(\varphi_1,\varphi_2)=(l_1\sin{\varphi_1},0,-l_1\cos{\varphi_1},l_2\sin{\varphi_2},0,-l_2\sin{\varphi_2})
\]
La forza attiva è la forza peso che come sappiamo ammette
\begin{gather*}
\tilde{V}(X)=mgz_1+mgz_2 \\ V(q)= \left. V(X)\right|_{X=\tilde{X}(q)}=-mg(l_1\cos{\varphi_1}+l_2\cos{\varphi_2})
\end{gather*}
Quindi per cercare gli equilibri del sistema
\[
\nabla V(\varphi_1,\varphi_2)=(mgl_1\sin{\varphi_1},mgl_2\sin{\varphi_2})
\]
e tutti gli equilibri possibili si hanno per $\varphi_{1,2}=0,\pi$. Sottolineiamo che ci sono solo due configurazioni di equilibrio su ciascuna coordinata e non infinite con $k\pi$ perchè solo due equilibri sono indipendenti.
\end{minipage}
\begin{minipage}[H]{0.25\linewidth}
\begin{figure}[H]\begin{center}
\includegraphics[width=\linewidth]{img/pendoloDoppio.png}
\end{center}\end{figure}
\end{minipage}

La matrice Hessiana del potenziale è $V''(\varphi_1,\varphi_2)=mg\left(\begin{array}{cc}
2l_1\cos{\varphi_1} & 0 \\
0 & 2l_2\cos{\varphi_2}
\end{array}\right)$ quindi chiaramente tutti gli equilibri aventi una delle due coordinate in $\pi$ hanno almeno un autovalore negativo e sono instabili.
Invece in $(0,0)$ la matrice Hessiana avrà due autovalori positivi e $(0,0)$ è un minimo stretto del potenziale ed un equilibrio stabile per tutti i tempi.
\end{esempio}

\subsection{Piccole oscillazioni}
Consideriamo  un sistema lagrangiano $L=\dfrac{1}{2}\dot{q}\cdot A(q)\dot{q} - V_0(q)$ tale che l'energia potenziale $V_0$ sia definita positiva.
Come abbiamo visto le equazioni linearizzate attorno ad un equilibrio, che in queste ipotesi sarà stabile, sono
\[
A(q^*)\ddot{q}=V_0''(q^*)(q-q^*)
\]

\begin{lemma}\label{lemmaFrequenze}
Consideriamo il sistema Lagrangiano linearizzato $A(q^*)\ddot{q}=V_0''(q^*)(q-q^*)$ attorno ad un equilibrio. Allora se $V_0''(q^*)$ è definita positiva con $q^*$ minimo stretto gli $A$-autovalori di $V_0''(q*)$ sono positivi.
\end{lemma}
\begin{dimo}
Gli $A$-autovalori di $V_0''$ sono $\{\lambda_1,\dots,\lambda_n\}$ e sappiamo che esistono $A$-autovettori di $V_0''$ soluzioni di $V_0''u_j=\lambda_j A u_j$.
Allora
\[
u_j\cdot V_0'' u_j = \lambda_j u_j \cdot A u_j =\lambda_j
\]
Ma essendo $V_0''$ definita positiva $u_j\cdot V_0'' u_j = k \left|\left|u_j\right|\right|^2>0$ e $\lambda_j>0$ e quindi $A^{-1}(q^*)V_0''(q^*)$ è definita positiva.\qedhere
\end{dimo}

In questa sezione consideriamo soltanto Lagrangiane con Hessiana dell'energia potenziale definita positiva nella configurazione di equilibrio quindi, essendone gli $A$-autovalori $\lambda_j$ tutti positivi possiamo porre $\lambda_j=\omega_j^2$ e quindi scegliere $\omega_j=+\sqrt{\lambda_j}$.

Chiameremo le $\omega_1,\dots,\omega_n$ così definite \textbf{frequenze delle piccole oscillazioni} attorno a $q^*$, esse sono le soluzioni di
\begin{equation}
\det{[V_0''(q^*)-=\omega_j^2A(q^*)]}=0 \qquad V_0''(q^*)\,u_j=\omega_j^2A(q^*)\,u_j
\end{equation}

\begin{proposizione}
Consideriamo il sistema Lagrangiano linearizzato $A(q^*)\ddot{q}=V_0''(q^*)(q-q^*)$ ed assumiamo che $V_0''(q^*)$ sia definita positiva e che $q^*$ ne sia un minimo stretto.
Siano $\omega_1,\dots,\omega_n$ le frequenze delle piccole oscillazioni attorno a $q^*$ e $u_1,\dots,u_n$ un insieme di $A$-autovettori linearmente indipendenti di $V_0''(q^*)$. Allora:\begin{enumerate}[i.]
\item $\text{sp}(\Lambda(q^*))=\{\pm i\omega_1,\dots,\pm i\omega_n\}$
\item Le equazioni di Lagrange linearizzate hanno l'integrale generale
\begin{equation}\label{eqIntGeneralePiccoleOscillazioni}
q(t;a,b)-q^*=\sum_{i=1}^n\left[  a_j\cos{\omega_j t} + b_j\sin{\omega_jt} \right]u_j \qquad a_j,b_j\in\mathbb{R} \text{ per } j=1,\dots, n
\end{equation}
\end{enumerate}
\end{proposizione}
\begin{dimo}
Nelle ipotesi date sappiamo che $\text{sp}(A^{-1}(q^*)V_0''(q^*))=\{\omega_1^2,\dots,\omega_n^2\} $ quindi dal lemma \ref{lemmaSpettroLambda} segue direttamente $\text{sp}(\Lambda(q^*))=\{\pm i\omega_1,\dots,\pm i\omega_n\}$.

Derivando due volte rispetto al tempo la \ref{eqIntGeneralePiccoleOscillazioni} otteniamo
\[
\ddot{q}=\sum_{j=1}^n -\omega_j^2 \left[a_j\cos{\omega_jt}+b_j\sin{\omega_jt}\right]u_j
\]
Sappiamo che $V_0'' u_j = \omega_j^2 A u_j$ e quindi $A\ddot{q}=-\omega_j^2A\sum_j(\dots)u_j=-V_0'' \sum_j(\dots) u_j=V_0''(q-q^*)$ e quindi la \ref{eqIntGeneralePiccoleOscillazioni} è soluzione del sistema linearizzato.\qedhere
\end{dimo}

Notiamo che ciascuno dei termini della sommatoria $\left[a_j\cos{\omega_j t} + b_j\sin{\omega_jt} \right]u_j$ è a sua volta una soluzione delle equazioni di Lagrange linearizzate e sono periodiche di periodo $\sfrac{2\pi}{\omega_j}$. Chiameremo queste soluzioni i \textbf{modi normali di oscillazione} di frequenza $\omega_j$.

I modi normali di oscillazione sono, per ogni $j$, una famiglia di soluzioni dipendente dai due parametri $a_j,b_j\in\mathbb{R}$. All' interno di ciascuna di queste famiglie le componenti di $q-q^*$ oscillano tutte in fase con frequenza $\omega_j$ ed ampiezze relative date dalle componenti di $u_j=(u_j^1,\dots,u_j^n)$.

Possiamo riscrivere in forma equivalente i modi normali di oscillazione come
\begin{equation}
c_j\cos{(\omega_jt+\delta_j)}u_j \qquad c_j=\sqrt{a_j^2+b_j^2}\geq 0 \quad j=1,\dots n
\end{equation}
L'integrale generale per le piccole oscillazioni è dato da una combinazione lineare dei modi normali
\begin{equation}
\left(\begin{array}{c}
q_1(t;c,\delta)-q_1^*\\\vdots\\q_n(t;c,\delta)-q_n^*
\end{array}\right)=\sum_{j=1}^n c_j\cos{(\omega_jt+\delta_j)}\left(\begin{array}{c}
u_j^1\\ \vdots\\u_j^n
\end{array}\right)
\end{equation}
ed in base alle condizioni iniziali in una piccola oscillazione possono essere ``eccitati'' diversi modi normali.

\begin{esempio}[Punto vincolato ad un paraboloide]

 
\begin{minipage}[H]{0.69\linewidth}
Consideriamo un punto materiale vincolato alla superficie interna di un paraboloide.
Q è parametrizzata da $z=\sfrac{1}{2} \left(ax^2by^2\right)$ con $a,b\in\mathbb{R}_+$. L' unica forza attiva agente sul punto è la forza peso, che è posizionale ed ammette un $V_0$, il vincolo è fisso e quindi $L=T_2-V_0$.
\[
V(x,y)=\left. mgz \right|_{z=\tilde{z}(x,y)} = \dfrac{1}{2}mg\left(ax^2+by^2\right)
\]
$\nabla V(x,y)=mg\left(\begin{array}{cc} ax & by \end{array}\right) \rightarrow V'(x,y)\overset{!}{=} (0,0)$ soltanto per $ (x,y)=(0,0)$ mentre
\end{minipage}
\begin{minipage}[H]{0.3\linewidth}\flushright
\includegraphics[width=\linewidth]{img/paraboloide.png}
\end{minipage}
 \[\text{He}V(x,y)=\left(\begin{array}{cc}
a & 0 \\ 0 & b \end{array}\right)\] non dipende dal punto ed è definita positiva quindi $(0,0)$ è una configurazione di equilibrio stabile.
Vediamo che $\dot{z}=ax\dot{x}+by\dot{y}$ quindi \[T(x,\dot{x},y,\dot{y})=\dfrac{1}{2}m\left(\dot{x}^2+\dot{y}^2+(ax\dot{x}+by\dot{y})^2\right)\]  e  dovremo cercare
\[
T(x,\dot{x},y,\dot{y}\overset{!}{=} \dfrac{1}{2}\left(\begin{array}{c}
\dot{x}\\ \dot{y} \end{array}\right) \cdot A(x,y)\left(\begin{array}{c}
\dot{x}\\ \dot{y} \end{array}\right)
\]
potremmo fare il calcolo esplicito ma dato che ci serve solo per $(x,y)=(0,0)$ semplifichiamo il conto a $T(q^*,\dot{q})=\sfrac{1}{2}\dot{q}\cdot A(q^*)\dot{q}$ e quindi
\[
T(0,0,\dot{x},\dot{y})=\dfrac{1}{2}m(\dot{x}^2+\dot{y}^2) \implies A(0,0)=\left(\begin{array}{cc} m & 0 \\ 0 & m \end{array} \right)
\]
Per trovare le piccole oscillazioni devo prima di tutto ricavare le frequenze quindi
\[
\det{V''(q^*)-\lambda A(q^*)}=\det{\left(\begin{array}{cc} mga-m\lambda & 0 \\ 0 & mgb-m\lambda \end{array} \right)} \overset{!}{=} 0 \implies (ga-\lambda)(gb-\lambda)=0 
\]
e quindi le due frequenze dei modi normali sono $\omega_1=\sqrt{ga}$ e $\omega_2=\sqrt{gb}$.
Da $\ker{V''(q^*)-A\lambda_1}=\ker{V''(q^*)-mga}$ scegliamo $u_1=\left(\begin{array}{c}1\\0
\end{array}\right)$ e $u_2=\left(\begin{array}{c}0\\1
\end{array}\right)$ quindi i due modi normali di oscillazione sono
\begin{gather*}
\left( \begin{array}{c} x_1(t) \\y_1(t) \end{array} \right) = \left( \begin{array}{c} c_1\cos{(\sqrt{ga} t +\delta_1)} \\ 0 \end{array}\right) \qquad
\left( \begin{array}{c} x_2(t) \\y_2(t) \end{array} \right) =  \left( \begin{array}{c} 0 \\ c_2\cos{(\sqrt{gb} t +\delta_2)} \end{array}\right) 
\end{gather*}
quindi uno è solo lungo l'asse $x$ e uno solo lungo l'assa $y$. Una piccola oscillazione generica sarà data da una somma dei due modi
\[
t\mapsto \left(\begin{array}{c}
c_1\cos{(\sqrt{ga}t+\delta_1)} \\ c_2\cos{(\sqrt{gb}t+\delta_2)} 
\end{array}\right)
\]
e sarà periodica solo se
\[
\left(\begin{array}{c} x(0)\\ y(0) \\ \dot{x}(0) \\ \dot{y}(0) \end{array}\right)=\left(\begin{array}{c} x(T)\\ y(T) \\ \dot{x}(T) \\ \dot{y}(T) \end{array}\right) \implies
 T_1=k\dfrac{2\pi}{\omega_1} = T_2 = k_2 \dfrac{2\pi}{\omega_2} \text{ per } k_1,k_2\in\mathbb{N}
\]
quindi se e solo se $\dfrac{\omega_1}{\omega_2}=\sqrt{\dfrac{a}{b}}\in\mathbb{Q}$.

Notiamo che se $a=b$ ci sarà una degenerazione: $V''(q^*)=mga\mathbb{I}$ e $A=m\mathbb{I}$ e quindi gli $A$-autovalori di $V''(q^*)$ saranno gli stessi dell'identità e cioè vi sarà un unico autovalore doppio, ogni vettore sarà un autovettore di questo e tutti i moti possibili saranno periodici.
\end{esempio}

\subsubsection{Oscillatori armonici disaccoppiati}

Siano $L=T_2-V_0$ con $V_0''(q^*)=0$ e $\{u_1, \dots, u_n\}$ una base $A$-ortonormale di $A$-autovettori di $V''(q^*)$.

Consideriamo il cambiamento di base $q\mapsto Q=U^{-1}(q-q^*)$ dove $U=\left(u_1,\dots ,u_n\right)$.


\begin{proposizione}
Consideriamo il sistema Lagrangiano di Lagrangiana $L$ avente in $q^*$ un minimo stretto di $V(q)$ con $V''(Q^*)$ definita positiva. Allora nelle coordinate $Q$ la Lagrangiana del sistema linearizzato diventa
\[
L(Q,\dot{Q})=\dfrac{1}{2}\sum_{i=1}^n \left(\dot{Q}^2-Q\omega_i^2Q_i^2\right)=\dfrac{1}{2}\left( \dot{Q}\cdot\dot{Q} - Q\cdot\Omega Q\right)
\]
\end{proposizione}
\begin{dimo}
La lagrangiana del sistema linearizzato nelle coordinate $q$ è
\[ L^*=\dfrac{1}{2}\dot{q}\cdot A(q^*)\dot{q} - \dfrac{1}{2}(q-q^*)\cdot V(q^*)(q-q^*)\]
quindi dato che $q=q^*+UQ$ e $\dot{q}=u\dot{Q}$ avremo


\begin{minipage}[H]{0.6\linewidth}
\begin{align*}
L(Q,\dot{Q})&=L^*(q^*+UQ,U\dot{Q})=\dfrac{1}{2} U\dot{Q}\cdot A U\dot{Q} - \dfrac{1}{2} UQ\cdot V'' UQ  =
\\
&= \dfrac{1}{2} \dot{Q}\cdot U^\dag AU \dot{Q} -\dfrac{1}{2} QU^\dag V'' UQ  
\\
&=\dfrac{1}{2} \dot{Q}\cdot\dot{Q} - \dfrac{1}{2} Q\Omega Q
\end{align*}
\end{minipage}
\begin{minipage}[H]{0.3\linewidth}\flushright
$\left(U^\dag AU\right)_{ij}=u_i \cdot A u_j = \delta_{ij}$  \\
$ \left(U^\dag V'' U \right)_{ij}=u_i\cdot V''u_j=$ $=u_i\cdot\omega_j^2 Au_j=\omega_j^2\delta_{ij} $
\end{minipage}

con $\Omega=\text{diag}\{\omega_1^2,\dots,\omega_n^2\}$ \qedhere
\end{dimo}

Ma una Lagrangiana $L(Q,\dot{Q})=\sum_i\dfrac{\dot{Q}_j^2-\omega_j^2 Q_j^2}{2}$ è la Lagrangiana di un sistema di $n$ oscillatori armonici disaccoppiati con frequenze $\omega_i$.
Ciascuno moto di un singolo oscillatore costituirà un modo normale di oscillazione del sistema.

\subsubsection{Modi normali e sottospazi invarianti}

Abbiamo introdotto i modi normali a partire dall'integrale generale del sistema linearizzato attorno alla configurazione di equilibrio $q^*$
\[
q(t;a,b)=\sum_{j=1}^n \left(a_j\cos{\omega_j t} +b_j\sin{\omega_j t} \right)\,u_j
\]
Se però studiamo le equazioni linearizzate come equazioni lineari
\[
\dot{z}=\Lambda(q^*)z \quad\qquad z=(q,\dot{q} )\in\mathbb{R}^{2n} \qquad \Lambda(q^*)=\left(\begin{array}{cc}
0 & \mathbb{I} \\ -A^-1(q^*)V''(q^*) & 0
\end{array}\right)
\]
dove come sappiamo se $V''(q^*)$ è definita positiva allora $\text{sp}\Lambda=\{\pm i\omega_1,\dots,\pm i\omega_n\}$.
Ma allora dato che $\Lambda$ è diagonalizzabile con tutti i suoi autovalori immaginari potremo decomporre $\mathbb{R}^{2n}$ nella somma diretta degli autospazi bidimensionali $\langle v_j,w_j \rangle$, dove $v_j\pm iw_j$ è l'autovettore di $\Lambda$ relativo all'autovalore $\pm i\omega_j$.

Ciascuno di questi autospazi sarà invariante e la dinamica ristretta su di esso sarà una rotazione e $\mathbb{R}^{2n}=E^c$ cioè sarà presente solo il sottospazio centrale.

\begin{lemma}
Per ogni $j=1,\dots ,n$ i $(u_j,\pm i\omega_j u_j)$ sono autovettori di $\Lambda$ relativi a $\pm i\omega_j$.
\end{lemma}
\begin{dimo}
\begin{align*}
\Lambda \left(\begin{array}{c}
u_j \\ i\omega_j u_j 
\end{array}\right) &=
\left(\begin{array}{cc} 0 & \mathbb{I} \\ -A^-1(q^*)V''(q^*) & 0 \end{array}\right) \left(\begin{array}{c}
u_j \\ i\omega_j u_j \end{array}\right) = \left( \begin{array}{c}
i\omega_j u_j \\ -A^{-1}(q^*)V''(q^*) u_j
\end{array} \right) = \\
&= \left( \begin{array}{c}
i\omega_j u_j \\ -\omega_j^2 u_j
\end{array}\right) =i\omega_j \left(\begin{array}{c}
u_j \\ i\omega_j u_j
\end{array} \right) \qedhere
\end{align*}
\end{dimo}

Quindi gli autovettori $v_j,w_j$ della decomposizione di $\mathbb{R}^{2n}$ sono $v_j=\left(\begin{array}{c}u_j \\ 0 \end{array}\right)$ e $w_j=\left(\begin{array}{c}0 \\ \omega_j u_j \end{array}\right)$
\begin{equation}
\mathbb{R}^{2n}=\left\langle \left(\begin{array}{c}u_1 \\ 0 \end{array}\right) , \left(\begin{array}{c}0 \\ \omega_1 u_1 \end{array}\right)  \right\rangle \bigoplus \dots \bigoplus \left\langle \left(\begin{array}{c}u_n \\ 0 \end{array}\right) , \left(\begin{array}{c}0 \\ \omega_n u_n \end{array}\right)  \right\rangle 
\end{equation}

Se consideriamo una soluzione $t\mapsto (q_t-q^*,\dot{q}_t)\in \langle v_j,w_j\rangle$ allora scrivendola nella base di $\{v_j,w_j\}$ questa diventa una $t\mapsto(x_t,y_t)$ e quindi
\[
\left(\begin{array}{c} q_t - q^* \\ \dot{q}_t \end{array}\right) = x_t \left(\begin{array}{c} u_j \\ 0 \end{array}\right) + y_t \left(\begin{array}{c} 0 \\ \omega_j u_j \end{array}\right)=\left(\begin{array}{c}
x_t \, u_j \\ y_t \, \omega_j u_j
\end{array}\right)
\]
da quanto visto nella sezione \ref{sezMatriciDiagonalizzabili} la restrizione del moto è una rotazione
\[
\left(\begin{array}{c}
x_t \\ y_t \end{array}\right) = R(\omega_j t)\left( \begin{array}{c}
x_0 \\ y_0 \end{array}\right)=\left(\begin{array}{cc}
\cos{\omega_j  t} & \sin{\omega_j t} \\ -sin{\omega_j t} & \cos{\omega_j t} \end{array} \right) \left(\begin{array}{c}
x_0 \\ y_0 \end{array}\right)
\]
e quindi una soluzione contenuta nel sottospazio $j$-esimo e è
\[
\left(\begin{array}{c}
q_t \\ \dot{q}_t \end{array}\right) = \left(\begin{array}{c}
\left[ \cos{\omega_j t} x_0 + \sin{\omega_j t} y_0 \right] u_1 \\
\left[ -\sin{\omega_j t}x_0 + \cos{\omega_j t} y_0 \right] \omega_j u_j 
\end{array}\right)
\]
la prima componente è il modo normale di frequenza $\omega_j$ e la seconda è la sua velocità.

Quindi la decomposizione in sottospazi della linearizzazione permette di restringere la dinamica ad essi e queste restrizioni nella base degli $A$-autovettori rotazioni e nella base normale i modi normali di oscillazione.

L'integrale generale per le piccole oscillazioni sarà dato dalla somma delle componenti sui sottospazi cioè da una combinazione dei modi normali.


\section{Simmetrie, integrali primi e riduzione}

\begin{definizione}[Azione di gruppi ad un parametro]
Definiamo azione di un gruppo $G$, con $G=\mathbb{R}$ o $G=\mathcal{S}^1$, su un aperto $U\subseteq \mathbb{R}^n$ una mappa differenziabile 
\[
\varphi: G\times U \rightarrow U \quad \qquad
(\lambda,q)\mapsto \varphi (\lambda,q)
\]
tale che le mappe $\varphi_\lambda :U\rightarrow U$ definite da $\varphi_\lambda(q)=\varphi(\lambda,q)$ abbiano le proprietà seguenti:\begin{enumerate}[i.]
\item $\varphi_0:U\rightarrow U$ è l'identità di $U$ 
\[\varphi_0=\mathbb{I}_U\]
\item Per ogni $\lambda, \mu\in G$ vale
\[
\varphi_{\lambda+\mu}=\varphi_\lambda\circ\varphi_\mu
\]
\item Per ogni $\lambda\in G$ la mappa $\varphi_\lambda$ è un diffeomorfismo
\end{enumerate}
\end{definizione}

Le proprietà ci dicono che le mappe $\varphi_\lambda$ con $\lambda\in G$ formano esse stesse un gruppo (commutativo) rispetto alla composizione.

Per un $q\in U$ chiamiamo orbita di $q$ sotto l'azione di $G$ $\varphi_\lambda$ l'insieme $\{\varphi_\lambda(q) : \lambda\in G\}$

\begin{esempio}[Flusso]
Per $U=\mathbb{R}^n$ il flusso di un'equazione differenziale $\dot{q}=\xi(q), q\in\mathbb{R}^n$ è un' azione di $\mathbb{R}$ su $\mathbb{R}^n$
\[
\varphi_\lambda(q)=\Phi_\lambda^\xi
\]
e le orbite dell'azione sono le orbite dell'equazione differenziale.
\end{esempio}


Un flusso di un campo vettoriale è un' azione di $\mathbb{R}$ ma non solo: ogni azione di $\mathbb{R}$ o $S^1$ su $\mathbb{R}^n$ è un flusso di un campo vettoriale su $\mathbb{R}^n$

\begin{definizione}
Il generatore infinitesimo di un azione $\varphi:G\times U\rightarrow U$, $(\lambda,q)\mapsto\varphi(\lambda,q)$ è il campo vettoriale definito da
\begin{equation}
\xi(q):=\dfrac{\partial\varphi}{\partial \lambda} (0,q)
\end{equation}
\end{definizione}

Il generatore infinitesimo associerà ad ogni $q\in U$ un vettore tangente ad $U$ in $q$.

\begin{proposizione}
Se $\xi$ è il generatore infinitesimo dell'azione $\varphi$ di $G=\mathbb{R},S^1$ allora
\[
\varphi=\Phi^\xi
\]
\end{proposizione}
\begin{dimo}
Il flusso di $\xi$ è la mappa $\Phi^\xi:\mathbb{R}\times U\rightarrow U$ tale che $\Phi^\xi(0,q)=0$ e che $\forall t,q$
\[
\dfrac{\partial\Phi^\xi}{\partial t}(t,q)=\xi(\Phi^\xi(t,q))
\]

Vogliamo mostrare che anche l'azione $\varphi$ ha queste due proprietà. 

La prima è banale $\varphi_0=\mathbb{I}_U$ quindi $\varphi_0(q)=q\equiv\Phi_0^\xi$.

Dalla definizione di generatore infinitesimo $\xi(q)=\left. \partial_t\varphi_t(q) \right|_{t=0}$ troviamo che per qualsiasi $\overline{t}\in\mathbb{R}$
\[
\left.\dfrac{d}{dt}\varphi_t(q)\right|_{t=\overline{t}}=
\left.\dfrac{d}{ds}\varphi_{s+\overline{t}}(q)\right|_{s=0}=\left.\dfrac{d}{ds}\varphi_s(\varphi_{\overline{t}})\right|_{s=0}=
\xi(\varphi_{\overline{t}}) \qedhere
\]
\end{dimo}


Anche se le azioni di gruppi di un parametro sembrano poco utili di per sè, dato che sono flussi, esustibi azioni di gruppi qualunque che non hanno una traduzione in flusso.

\begin{esempio}[Traslazioni e rotazioni]
\begin{itemize}
\item Le traslazioni in $U=\mathbb{R}^n$ lungo un $e\in\mathbb{R}^n\setminus\{0\}$ sono azioni 
$\varphi_\lambda(q)=q+\lambda \, e$ con generatore infinitesimo $\xi(q) = e$ $\forall q$
\item Le rotazioni in $\mathbb{R}^3$ attorno ad un asse  $e\in\mathbb{R}^3\setminus\{0\}$ sono azioni $\varphi_\lambda(q)=R_\lambda q $ con
\[
\left.\dfrac{d}{d\lambda}\varphi_\lambda(q)\right|_{\lambda=0}=\left.\dfrac{d}{d\lambda}R_\lambda q \right|_{\lambda=0}
\]
e 
\[
\dfrac{d}{d\lambda} R_\lambda q = \dot{R}_\lambda q = \dot{R}_\lambda R_\lambda^\dag R_\lambda q = (\dot{R}_\lambda R_\lambda^\dag)R_\lambda q
\]
quindi il generatore infinitesimo è
\[
\xi(q)=\left. (\dot{R}_\lambda R_\lambda^\dag )R_\lambda q \right|_{\lambda=0} = \left. a\times  R_\lambda q\right|_{\lambda=0} =a\times q
\]
\end{itemize}
\end{esempio}

\begin{definizione}
Una funzione $f:U\rightarrow\mathbb{R}$ si dice invariante sotto un'azione $\varphi$ di gruppo $G$ su $U$ se 
\[
f\circ\varphi_\lambda=f \qquad \forall \lambda \in G
\]

\end{definizione}
$f$ è invariante per $\varphi_\lambda$ se è costante lungo le sue orbite ed equivalentemente è integrale primo del flusso del suo generatore infinitesimo.

\begin{definizione}
Il sollevamento tangente di un'azione $\varphi$ di $G$ su $Q\subseteq\mathbb{R}^n$ è l'azione di $G$ su $TQ$ 
\[
\varphi^{TQ}_\lambda :G\times TQ \rightarrow TQ \qquad \quad \varphi_\lambda^{TQ}(q,\dot{q})=T\varphi_\lambda \, (q,\dot{q})=(\varphi_\lambda(q),\varphi_\lambda'(q)\,\dot{q})
\]
\end{definizione}
\begin{esempio}
Per una traslazione in $\mathbb{R}^n$ di azione $\varphi_\lambda(q)=q+\lambda\, e$ il sollevamento tangente è $T\varphi_\lambda (q,\dot{q})=(q+e,\dot{q})$.

Per una rotazione in $\mathbb{R}^3$ attorno all'asse $a$ di azione $\varphi_\lambda(q)=R_\lambda q$ il sollevamento tangente è $T\varphi_\lambda(q,\dot{q})=(R_\lambda q, R_\lambda \dot{q})$.
\end{esempio}
\subsection{Teorema di Noether}
Consideriamo una lagrangiana $L:TQ\rightarrow\mathbb{R}$, non necessariamente meccanica e possibilmente anche dipendente dal tempo, allora per un'azione $\varphi$ di $G$ su $Q$ diremo che $L$ è invariante sotto $\varphi$, o sotto l'azione sollevata $\varphi^{TQ}$, se 
\begin{equation}
L(\varphi_\lambda(q),\varphi_\lambda'(q)\,\dot{q})=L(q,\dot{q}) \qquad \quad \forall  \lambda\in G \> \forall (q,\dot{q})\in TQ
\end{equation}

\begin{proposizione}
Se la Lagrangiana $L$ è invariante sotto l'azione $\varphi$ allora le equazioni di Lagrange per $L$ hanno l'integrale primo
\begin{equation}
I(q,\dot{q}):=\xi\cdot p(q,\dot{q})
\end{equation}
dove $\xi$ è il generatore infinitesimo di $\varphi$ e $p=\dfrac{\partial L}{\partial\dot{q}}$ è il vettore dei momenti coniugati alle $q$.
\end{proposizione}
\begin{dimo}
Sia $t\mapsto q_t$ una soluzione delle equazioni di Lagrange. Per ipotesi 
\[
L(q_t,\dot{q}_t)=L(\varphi_\lambda(q_t),\varphi_\lambda'(q_t)\,\dot{q}_t) \qquad \quad \forall \lambda,t
\]
Allora
\begin{align*}
0=\dfrac{\partial}{\partial\lambda}L(q_t,\dot{q}_t)&=L(\varphi_\lambda(q_t),\varphi_\lambda'(q_t)\,\dot{q}_t) = \dfrac{\partial}{\partial \lambda} L(\varphi(\lambda, q_t),\dfrac{d}{d t}\varphi(\lambda, q_t) ) =\\
&=\sum_{j=1}^n\left[\dfrac{\partial L}{\partial q_j}\dfrac{\partial \varphi^j}{\partial \lambda}(\lambda,q_t) + \dfrac{\partial L}{\partial \dot{q}_j} \dfrac{\partial}{\partial \lambda}\dfrac{d}{dt}\varphi^j(\lambda,q_t) \right]\footnotemark[1]
\end{align*}
Possiamo scambiare l'ordine delle derivate $\partial_\lambda, dt$ di $\varphi(\lambda,q_t)$ quindi
\[
0=\sum_{j=1}^n\left[\dfrac{\partial L}{\partial q_j}\dfrac{\partial \varphi^j}{\partial \lambda}(\lambda,q_t) + \dfrac{\partial L}{\partial \dot{q}_j} \dfrac{d}{dt}\dfrac{\partial}{\partial \lambda}\varphi^j(\lambda,q_t) \right] \qquad\quad \forall \lambda,t
\]
Calcoliamo in $\lambda=0$ ricordando che $\varphi(0,q_t)=q_t$ e usando la definizione di generatore infinitesimo  
\begin{align*}
0&=\sum_{j=1}^n\left[ \left. \dfrac{\partial L(\varphi_\lambda(q_t),\varphi_\lambda'(q_t)\dot{q}_t)}{\partial q_j}\dfrac{\partial \varphi^j}{\partial \lambda}(\lambda,q_t)\right|_{\lambda=0} + \left. \dfrac{\partial L(\varphi_\lambda(q_t),\varphi_\lambda'(q_t)\dot{q}_t)}{\partial \dot{q}_j} \dfrac{d}{dt}\dfrac{\partial}{\partial \lambda}\varphi^j(\lambda,q_t)\right|_{\lambda=0} \right] =\\
&=\sum_{j=1}^n\left[ \dfrac{\partial L(q_t,\dot{q}_t)}{\partial q_j}\xi_j(q_t)  +  \dfrac{\partial L(q_t,\dot{q}_t)}{\partial \dot{q}_j} \dfrac{d}{dt}\xi_j(q_t) \right] 
\end{align*}
e dato che $t\mapsto q_t$ è una soluzione $\dfrac{\partial}{\partial q_j}L(q_t,\dot{q}_t)=\dfrac{d}{dt}\dfrac{\partial}{\partial\dot{q}_j}L(q_t,\dot{q}_t)$ quindi
\[
0=\sum_{j=1}^n \dfrac{d}{dt} \left( \dfrac{\partial L(q_t,\dot{q}_t)}{\partial \dot{q}_j}\right)\xi_j(q_t)  + \sum_{j=1}^n \dfrac{\partial L(q_t,\dot{q}_t)}{\partial \dot{q}_j} \dfrac{d}{dt}\xi_j(q_t) =\sum_{j=1}^n \dfrac{d}{dt} \left( \dfrac{\partial L(q_t,\dot{q}_t)}{\partial\dot{q}_j} \xi_j(q_t) \right)=0
\]
Lungo una soluzione la quantità $\dfrac{\partial L}{\partial \dot{q_j}}\xi_j=\xi\cdot p$ è costante e quindi è un integrale primo.\qedhere
\end{dimo}

\footnotetext[1]{Omettiamo di scrivere le dipendenze nelle derivate della lagrangiana sottointendendo $L(\varphi_\lambda(q_t),\varphi_\lambda'(q_t)\dot{q}_t)$.}

L'integrale primo $I=\xi\cdot p$ si chiama momento dell'azione $\varphi$.
Il teorema di Noether\footnotemark[2] permette di trvoare integrali primi a partire da un'invarianza o simmetria del sistema.

\footnotetext[2]{Qua il teorema è presentato in una versione abbastanza specializzata, nel nostro corso di studio è trattato anche in \textit{Fisica Teorica} e, nella sua versione generale, in \textit{Campi elettromagnetici}.}

\begin{esempio}[Punto materiale non vincolato soggetto a forze attive conservative]
Per un punto materiale soggetto solo a forze attive conservative la Lagrangiana espressa in coordinate cartesiane è
\[
L(x,dot{x})=\dfrac{1}{2}m||\dot{x}||^2-V(q)
\] 
Se consideriamo per esempio l'azione di una traslazione $\varphi_\lambda(q)=q+\lambda e$ vediamo che $(L\circ\varphi_\lambda)(q)=\dfrac{1}{2}m||\dot{x}||^2-V(q+\lambda e)$.

Quindi la Lagrangiana sarà invariante per la traslazione se e solo se l'energia potenziale è invariante per quella traslazione.
In tal caso l'integrale primo $\xi\cdot = e\cdot m\dot{q}$è la componente della quantità di moto nella direzione della traslazione.

Se l'energia potenziale è invariante per traslazioni lungo direzioni qualsiasi allora si conservano tutte le componenti della quantità di moto.
\par ~\par\noindent
\textbf{Coordinate ignorabili:}\par\indent
Per una $L(q_1,\dots,q_n,\dot{q}_1,\dots,\dot{q}_n)$ sia una coordinata qualsiasi, diciamo $q_n$, ignorabile $\dfrac{\partial L}{\partial q_n}=0$. Allora chiaramente $L$ sarà invariante per tralsazioni su $q_n$ e si avrà un $\xi(q)=(0,\dots,0,1)$ e un integrale primo $I=p_n=\dfrac{\partial L}{\partial \dot{q}_n}$
\par ~\par\noindent
\textbf{Punto materiale in un campo di forze centrali:}\par\indent
Per un punto materiale soggetto a sole forze centrali la lagrangiana è
\[
L(q,dot{q})=\dfrac{1}{2}m||\dot{q}||^2-V(||q||)
\] 
L'azione di una rotazione attorno ad un asse $e$ $\varphi_\lambda(q)=R_\lambda q$ cambia la Lagrangiana in $(L\circ\varphi_\lambda)(q)=\dfrac{1}{2}m ||R_\lambda q||^2 -V(||R_\lambda q||)=\dfrac{1}{2}m ||q||^2 -V(||q||)$ quindi la Lagrangiana è invariante per rotazioni.

L'integrale primo dal teorema di Noether sarà $\xi\cdot p=(e\times q)\cdot p=e\cdot(mq\times \dot{q})$ che è la componente del momento angolare lungo l'asse $e$ di rotazione. 
Ma essendo la Lagrangiana invariante per qualsiasi rotazione tutte le componenti del momento angolare si conservano se le forze attive sono centrali. 
\par ~\par\noindent
\textbf{Problema dei due corpi:}\par\indent
Consideriamo due corpi di masse $m_1$ e $m_2$ che interagiscono fra loro con una forza centrale conservativa. Usando le coordinate cartesiane $(r_1,r_2)\in\mathbb{R}^3\times\mathbb{R}^3$ la Lagrangiana è
\[
L(r_1,r_2,\dot{r}_1,\dot{r}_2)=\dfrac{1}{2}\left( m_1||\dot{r}_1||^2 + m_2 ||\dot{r}_2||^2 \right) - V(||r_1-r_2||)
\]
e se consideriamo un'azione $\varphi_\lambda(r_1,r_2)=(r_1+\lambda e,r_2+ \lambda e)$ che trasla entrambi i punti nello stesso modo chiaramente la Lagrangiana è invariante e quindi l'integrale primo è
\[
I(r_1,r_2,\dot{r}_1,\dot{r}_2)=\xi(r_1,r_2)\cdot p(r_1,r_2,\dot{r}_1,\dot{r}_2) = \left(\begin{array}{c} e\\e \end{array}\right)\left(\begin{array}{c}
m_1\dot{r}_1 \\ m_2\dot{r}_2 \end{array}\right) = e \left( m_1\dot{r}_1 + m_2\dot{r}_2\right)
\]
che è la componente lungo $e$ del momento totale, questo vale per $e$ qualsiasi quindi tutte le componenti del momento totale sono conservate.
\end{esempio}







Consideriamo un sistema lagrangiano con $n$ gradi di libertà che abbia almeno una coordinata ignorabile.

Da $\dfrac{\partial L}{\partial q_i}=0$ segue che la Lagrangiana è invariante per traslazioni lungo $q_i$ e che $p_i$ è un integrale primo.

Questo significa che lo spazio degli atti di moto $\mathbb{R}^{2n}$ è foliato da $p_i$ nel senso visto nella sezione \ref{sezFoliazione} e quindi possiamo abbasare la dimensione delle equazioni di Lagrange da $2n$ a $2n-1$.

Nel metodo di Riduzione alla Routh riusciamo a fare di più, sostanzialmente riusciremo ad esprimere le equazioni con $n-1$ gradi di libertà, cioè in $\mathbb{R}^{2n-1}$.

Dividiamo le coordinate $q$ in $q=(Q,\xi)=(Q_1,\dots,Q_{n-1},xi)$ dove $\xi$ è la coordinata ignorabile e le $Q$ sono le altre.
Chiamiamo $\pi(Q,\dot{Q},\dot{\xi})$ il momento coniugato alla coordinata $\xi$ 
\begin{equation}
\pi(Q,\dot{Q},\dot{\xi})=\dfrac{\partial}{\partial \dot{\xi}} L(Q,\dot{Q},\dot{\xi})
\end{equation}
Chiamiamo $\Sigma_c$ gli insiemi di livello $c$ di $\pi$, che come sappiamo è un integrale primo.

Sostanzialmente nella riduzione alla Routh vorremo usare il sistema di coordinate $(Q,\dot{Q},\xi,\pi)$ al posto di quello normale contenente $\dot{\xi}$. Per farlo però servirà che la mappa $C:(Q,\dot{Q},\xi,\dot{\xi})\mapsto(Q,\dot{Q},\xi,\pi)$ sia un diffeomorfismo almeno locale. 

Perchè questo avvenga serve che la matrice $\dfrac{\partial^2 L}{\partial \dot{\xi}\partial\dot{\xi}}$ sia invertibile, cioè che abbia determinante non nullo in ogni punto.

Come sappiamo la matrice $\dfrac{\partial^2 L}{\partial \dot{\xi}\partial\dot{\xi}}$ è un ``pezzo'' della matrice $\dfrac{\partial^2 L}{\partial \dot{q}\partial\dot{q}}$

\[
\dfrac{\partial^2 L}{\partial \dot{q}\partial\dot{q}}=
\left(\begin{array}{c|c} 
\dfrac{\partial^2 L}{\partial \dot{Q}\partial\dot{Q}} & 
\dfrac{\partial^2 L}{\partial \dot{Q}\partial\dot{\xi}} \\ \hline
\dfrac{\partial^2 L}{\partial \dot{\xi}\partial\dot{Q}} & \dfrac{\partial^2 L}{\partial \dot{\xi}\partial\dot{\xi}}
\end{array}\right)
\]
 Quindi se $\dfrac{\partial^2 L}{\partial \dot{q}\partial\dot{q}}$ è definita positiva anche $\dfrac{\partial^2 L}{\partial \dot{\xi}\partial\dot{\xi}}$ dovrà essere definita positiva e $C$.

\begin{lemma}\label{lemmaPerRouth}
Consideriamo un sistema di Lagrangiana $L$ con $n$ gradi di libertà e una coordinata $q_n=\xi$ ignorabile. 
Supponendo che la matrice $\dfrac{\partial^2 L}{\partial \dot{q}\partial\dot{q}}$ sia definita positiva: $\forall c\in\mathbb{R}$\begin{enumerate}[i.]
\item L'insieme $\Sigma_c=\{\pi(Q,\dot{Q},\dot{\xi}\}$ è invariante per traslazioni lungo $\xi$, cioè \[
(Q_1,\xi_1,\dot{Q}_1,\dot{\xi}_1)\in\Sigma_c\implies(Q_1,\xi_2,\dot{Q}_1,\dot{\xi}_1)\in\Sigma_c \> \forall \xi_1,\xi_2
\]
\item $\Sigma_c$ è una sottovarietà di dimensione $2n-1$ che può essere rappresentata, almeno localmente, come grafico  di una $U_c(Q,\xi,\dot{Q})$ che è in realtà indipendente da $\xi$. Quindi \[
\dot{\xi}=U_c(Q,\dot{Q})
\]
\end{enumerate}
\end{lemma}
\begin{dimo}
La prima proprietà è chiara dato che $\pi$ non dipende da $\xi$.

Per mostrare che $\Sigma_c$ è una sottovarietà (ipersuperficie) basta verificare che il gradiente di $\pi$ non si annulli mai.

Sappiamo che $\pi=\dfrac{\partial L}{\partial \dot{\xi}}$ e quindi $\dfrac{\partial \pi}{\partial \dot{\xi}}=\dfrac{\partial^2 L}{\partial \dot{\xi}\partial\dot{\xi}}$ ma abbiamo assunto che $\dfrac{\partial^2 L}{\partial\dot{q}\partial\dot{q}}$ sia definita positiva quindi $\dfrac{\partial \pi}{\partial \dot{\xi}}$, che ne è un minore principale, non può annullarsi.

Quindi il gradiente di $\pi$ non potrà mai annullarsi.

Dato che $\dfrac{\partial \pi}{\partial \dot{\xi}}$ non si annulla mai il teorema della funzione implicita assicura che $\pi(Q,\dot{Q},\dot{\xi})=c$ possa essere esplicitata in una $U_c(Q,\dot{q})=\dot{\xi}$ $\forall (Q,\dot{Q},\dot{\xi})\in\Sigma_c$ per ogni $c$, dove l'invarianza di $U_c$ da $\xi$ discende dall'invarianza di $\Sigma_c$ per traslazioni lungo $\xi$.\qedhere
\end{dimo}


Dal lemma sappiamo che per ogni $c$ fissata la $(Q,\xi,\dot{Q})\mapsto(Q,\xi,\dot{Q},U_c(Q,\dot{Q})$ è una parametrizzazione locale di $\Sigma_c$.
Ma allora la restrizione delle equazioni di Lagrange a $\Sigma_c$ si potrà vedere come un'equazione nelle $(Q,\xi,\dot{Q})$.

Fissato un valore $c$ del momento conservato $\pi$ una soluzione $t\mapsto(Q_t,\xi_t,\dot{Q}_t,\dot{\xi}_t)$ delle equazioni di Lagrange sarà contenuta in $\Sigma_c$ per ogni $t$ e quindi 
\begin{equation}
\dot{\xi}_t=U_c(Q_t,\dot{Q}_t)
\end{equation}

Se è nota la componente $t\mapsto Q_t$ di una soluzione possiamo determinare la $\xi_t$ semplicemente tramite
\begin{equation}
\xi_t=\xi_0+\int_0^t U_c(Q_s,\dot{Q}_s) ds
\end{equation}
Inoltre dato che le equazioni di Lagrange sono equazioni del secondo ordine che si possono mettere in forma normale la $t\mapsto Q_t$ è soluzione di una $\ddot{Q}_t=Y(Q_t,\dot{Q}_t,\dot{\xi}_t)$, non dipendente da $\xi$ perchè è ignorabile. Ma allora 
\[
\ddot{Q}_t=Y(Q_t,\dot{Q},U_c(Q_t,\dot{Q}_t)
\]
e quindi per $X_c(Q_t,\dot{Q}_t):=Y(Q_t,\dot{Q}_t,U_c(Q_t,\dot{Q}_t)$ le componenti nelle $Q$ delle soluzioni, $t\mapsto Q_t$, saranno soluzioni di un sistema ridotto di $n-1$ equazioni del secondo ordine $\ddot{Q}_t=X_c(Q_t,\dot{Q}_t$. Possiamo quindi ricavare la $t\mapsto Q_t$ a prescindere dalla $\xi_t$ e usarla per ricostruire il moto tramite la $U_c$.

\begin{proposizione}Sia una Lagrangiana $L$ con $\xi$ coordinata ignorabile e $\sfrac{\partial^2 L}{\partial\dot{q}\partial\dot{q}}$ definita positiva. Data $t\mapsto (Q_t,\xi_t)$ soluzione delle equazioni di Lagrange per $L$ con dato iniziale $(Q_0,\xi_0,\dot{Q}_0,\dot{\xi}_0)$ e $\pi(Q_0,\xi_0,\dot{Q}_0,\dot{\xi}_0)=c$ si ha:\begin{enumerate}[i.]
\item $t\mapsto Q_t$ è soluzione delle equazioni di Lagrange di Lagrangiana $L^R_c$ (\textit{Lagrangiana ridotta})
\begin{equation}
L^R_c(Q,\dot{Q}):=L(Q,\dot{Q},U_c(Q,\dot{Q}) - c\,U_c(Q,\dot{Q})
\end{equation}
con la $U_c$ definita nel lemma \ref{lemmaPerRouth}.

\item $t\mapsto\xi_t$ è data da $\xi_t=\xi_0+\int_0^t U_c(Q_s,\dot{Q}_s) ds$.
\end{enumerate}
\end{proposizione}
\begin{dimo}

Fissato un valore $c=\pi(q_0,\dot{q}_0)$ una soluzione $t\mapsto(Q_t,\xi_t,\dot{Q}_t,\dot{\xi}_t)$ delle equazioni di Lagrange sarà contenuta in $\Sigma_c$ per ogni $t$ quindi dal lemma avremo $\dot{\xi}_t=U_c(Q_t,\dot{Q}_t)$. Nota la $t\mapsto Q_t$ possiamo usare $\xi_t=\xi_0+\int_0^t U_c(Q_s,\dot{Q}_s) ds $. 

Derivando con la regola della catena
\begin{align*} 
\dfrac{\partial }{\partial Q_i}L^R_c(Q,\dot{Q}) & =\dfrac{\partial}{\partial Q_i} \left( L(Q,\dot{Q},U_c(Q,\dot{Q}) - c\,U_c(Q,\dot{Q}) \right)=\\
&=\dfrac{\partial L}{\partial Q_i}(Q,\dot{Q},U_c(Q,\dot{Q}) + \dfrac{\partial L}{\partial Q_i}(Q,\dot{Q},U_c(Q,\dot{Q})\dfrac{\partial U_c}{\partial Q_i}(Q,\dot{Q}) -c\dfrac{\partial U_c}{\partial Q_i}(Q,\dot{Q})
\end{align*}
Però lungo le soluzioni $\pi(Q_t,\dot{Q}_t,U_c)=c$ $\forall t$ quindi lungo una soluzione
\[
\dfrac{\partial }{\partial Q_i}L^R_c(Q,\dot{Q}) =\dfrac{\partial L}{\partial Q_i}(Q_t,\dot{Q}_t,U_c(Q_t,\dot{Q}_t)
\] 
Con lo stesso procedimento si trova
\[
\dfrac{\partial }{\partial \dot{Q}_i}L^R_c(Q,\dot{Q}) =\dfrac{\partial L}{\partial \dot{Q}_i}(Q_t,\dot{Q}_t,U_c(Q_t,\dot{Q}_t)
\] 

Sapendo che $\xi_t = U_c(Q_t,\dot{Q}_t)$ si ha che le soluzioni $t\mapsto(Q_t,\xi_t)$ soddisfano
\begin{align*}
\dfrac{d}{dt}\dfrac{\partial L_c^R}{\partial \dot{Q}_{i}}(Q_t,\dot{Q}_t) &= \dfrac{d}{dt} \dfrac{\partial L}{\partial \dot{Q}_i}(Q_t,\dot{Q}_t,U_c(Q_t,\dot{Q}_t)=  \dfrac{d}{dt} \dfrac{\partial L}{\partial \dot{Q}_i}(Q_t,\dot{Q}_t,\xi_t)=\\
\text{eq. Lagrange } L  &=  \dfrac{\partial L}{\partial \dot{Q}_i}(Q_t,\dot{Q}_t,\xi_t)=\\
&=\dfrac{\partial L^R_c}{\partial Q_i}(Q_t,\dot{Q_t})
\end{align*}
e quindi le $t\mapsto Q_t$ delle soluzioni delle equazioni di Lagrange di Lagrangiana $L$ sono soluzioni delle equazioni di Lagrange di Lagrangiana $L^R_c$. \qedhere
\end{dimo}

\begin{esempio}[Punto vincolato ad un piano e soggetto a forze centrali]
La Lagrangiana di un punto soggetto a forze centrali è $L(x,\dot{x}=\sfrac{1}{2}m||\dot{x}||^2-V(||x||)$ e come abbiamo visto è invariante per rotazioni attorno a qualsiasi asse. Quindi se il punto è anche vincolato ad un piano
\[
L(r,\theta,\dot{r},\dot{\theta})=\dfrac{1}{2}m(\dot{r}^2+r^2\dot{\theta}^2)-V(r)
\]
e $\theta$ è ignorabile. Quindi $p_\theta=mr^2\dot{\theta}$ è un integrale primo e per $p_\theta=c$ fissato avremo $\dot{\theta}=\dfrac{c}{mr^2}=U_c(r)$.

La Lagrangiana ridotta sarà
\begin{align*}
L_c^R(r,\dot{r})&=L(r,\dot{r},U_c(r)-cU_c(r)=\\
&= \dfrac{m}{2}(\dot{r}^2+r^2U_c^2(r))-V(r)-cU(r)=\\
&=\dfrac{m}{2}\dot{r}^2 -\dfrac{c^2}{2mr^2}-V(r)
\end{align*}
In questo caso possiamo interpretare la Lagrangiana ridotta come un $\dfrac{m\dot{r}^2}{2}-W_c(r)$ di un sistema unidimensionale con \textit{potenziale efficace} $W_c{r}=V(r)+\dfrac{c^2}{2mr^2}$.

L'equazioni ridotta del sistema è quindi
\[
m\ddot{r}=W'_c(r))=-V'(r)+\dfrac{c^2}{2mr^3}
\]
\end{esempio}

Per calcolare la Lagrangiana ridotta di una Lagrangiana $L$ meccanica, o che comunque abbia dipendenza da $\dot{\xi}$ solo quadratica e lineare, possiamo usare un procedimento più veloce.
Scriviamo $L=\hat{L}_2+\hat{L}_1+\hat{L}_0$ dove il pedice indica il grado di omogeneità in $\dot{\xi}$.
La Lagrangiana ridotta è $L_c^R=\left[ L-\dot{\xi}\dfrac{\partial L}{\partial \dot{\xi}}\right|_{\dot{\xi}=U_c(Q,\dot{Q})}$ e dal teorema di Euler sulle funzioni omogenee sappiamo che
\[
\dot{\xi}\dfrac{\partial \hat{L}_s}{\partial\dot{\xi}}=s\hat{L}_s \qquad s=0,1,2
\]
quindi \begin{align*}
L-\dot{\xi}\dfrac{\partial L}{\partial \dot{\xi}}&=L-\dot{\xi}\dfrac{\partial \hat{L}_2}{\partial\dot{\xi}}-\dot{\xi}\dfrac{\partial \hat{L}_1}{\partial\dot{\xi}}-\dot{\xi}\dfrac{\partial \hat{L}_0}{\partial\dot{\xi}}=\\
&=\hat{L}_2-2\hat{L}_2+\hat{L}_1-\hat{L}_1+\hat{L}_0
\end{align*}
Allora
\begin{equation}
L_c^R(Q,\dot{Q})=\hat{L}_0(Q,\dot{Q})-\left.\hat{L}_2(Q,\dot{Q},\dot{\xi})\right|_{\dot{\xi}=U_c(Q,\dot{Q})}
\end{equation}

I moti del sistema completo che si proiettano sugli equilibri del sistema ridotto sono chiamati \textbf{equilibri relativi}. 

Sia $\overline{Q}$ una configurazione di equilibrio del sistema ridotto, allora i moti del sistema che si proiettano su di essa avranno $Q_t=\overline{Q}$ e $\xi_t=\xi_0+\int_0^t U_c(\overline{Q},0) ds=\xi_0+tU_c(\overline{Q},0)$.

Quindi negli equilibri relativi la coordinata ignorabile si muove con velocità costante dipendente solo da $c$. 

Nel caso particolare in cui la coordinata ignorabile è un angolo allora gli equilibri relativi saranno moti periodici di periodo $\sfrac{2\pi}{U_c(\overline{Q},0)}$.


\begin{esempio}[Pendolo sferico]
Consideriamo un punto vincolato alla superficie di una sfera e soggetto alla forza peso. Il sistema sarà chiaramente invariane per rotazioni attorno all'asse verticale e possiamo usare il procedimento di riudzioni per integrarlo.

Studiamo il sistema in coordinate polari sferiche, quindi i moti passanti per i poli andrebbero studiati a parte.

La lagrangiana è
\[
L(\theta,\varphi,\dot{\theta},\dot{\varphi})=\dfrac{mR^2}{2}(\dot{\theta}^2+\dot{\varphi}^2\sin^2{\theta})-mgR\cos{\theta}
\]
è conveniente dividere tutto per $mR^2$ ottenendo la nuova Lagrangiana
\[
L(\theta,\varphi,\dot{\theta},\dot{\varphi})=\dfrac{1}{2}(\dot{\theta}^2+\dot{\varphi}^2\sin^2{\theta})-k\cos{\theta}\qquad \quad k=\dfrac{g}{R}
\]

Come avevamo previsto l'angolo $\varphi$ è ignorabile e quindi il suo momento  $p_\varphi=\dot{\varphi}\sin^2{\theta}$ si conserva (fattore $mR^2$ o meno). Allora fissando $p_\varphi=J$
\[
\dot{\varphi}=\dfrac{J}{\sin^2{\theta}}
\]
e la Lagrangiana ridotta è
\begin{align*}
L^R_c=\hat{L}_0-\left.\hat{L}_2\right|_{\dot{\varphi}=\dfrac{J}{\sin^2{\theta}}}&=\dfrac{1}{2}\dot{\theta}^2-k\cos{\theta}-\left.\dfrac{1}{2}\dot{\varphi}^2\sin^2{\theta}\right|_{ \dot{\varphi}=\dfrac{J}{\sin^2{\theta}} }=\\
&=\dfrac{1}{2}\dot{\theta}^2-k\cos{\theta}-\dfrac{1}{2}\dfrac{J^2}{\sin^2{\theta}}=\dfrac{1}{2}-W_J{\theta}
\end{align*}
Questa è la Lagrangiana di un sistema ad un solo grado di libertà con spazio delle configurazioni l'intervallo $(0,\pi)$ e soggetto a forze conservative con energia potenziale $W_j(\theta)$.

Per $J=0$ i moti passano per i poli e le coordinate non funzionano quindi consideriamo solo $J\neq0$.


\begin{minipage}[H]{0.65\linewidth}
Deriviamo $W_j$ per verificare se il sistema ridotto ha equilibri:
\[
W_j'(\theta)=-k\sin{\theta}-\dfrac{J^2\cos{\theta}}{\sin^3{\theta}}=-\dfrac{1}{\sin^3{\theta} }\left( k\sin^4{\theta}+J^2\cos{\theta} \right)
\]
Ricerchiamo gli zeri di $W_j'(\theta)$ graficamente e troviamo un unico zero $\theta_m$ compreso tra $\sfrac{\pi}{2}$ e $\pi$.
\end{minipage}
\begin{minipage}[H]{0.3\linewidth}
\centering 
\includegraphics[width=\linewidth]{img/pendoloSfericoPotenziale.png}
\end{minipage}
Quindi il sistema ridotto avrà un unico equilibri e tutti i suoi altri moti saranno periodici, nei moti periodici la coordinate $\theta$ varierà tra due estremi determinati dall'energia e da $J$.

Ricostruiamo il sistema completo: lungo una soluzione del sistema ridotto $t\mapsto\theta(t)$
\[
\varphi(t)=\varphi(0)=\int_0^t \dfrac{J^2}{\sin^2{\theta(s)} } ds
\]
Negli equilibri relativi $\theta(t)=\theta_m(J)$ e $\varphi(t)=\varphi(0)+\dfrac{J^2}{\sin^2{\theta_m} }t$, cioò il pendolo ruota con velocità costante lungo un parallelo (cioè a $\theta=\theta_m$ fisso).

\begin{minipage}[H]{0.65\linewidth}
Negli altri moti $\theta$ si muove periodicamente tra due estremi e $\varphi$ e data dall formula precedente quindi è in generale monotòna. Quindi negli altri moti la coordinata $\theta$ oscilla mentre $\varphi$ avanza formando delle traiettorie tipo quelle in figura
Potremmo chiederci se questi moti sono periodici, cioè se si chiudono per qualche $T$ $(\theta(t),\varphi(t))=(\theta(t+T),\varphi(t+T))$ ma la risposta è difficile da raggiungere.
\end{minipage}
\begin{minipage}[H]{0.3\linewidth}\flushright
\includegraphics[width=\linewidth]{img/pendoloSfericoPeriodici.png}
\end{minipage}
\end{esempio}

\section{Problema di Kepler}

Il problema di Kepler (o sistema) è un sistema meccanico costituito da un punto materiale in un campo di forza gravitazionale, quindi centrale, con centro fisso.

Vogliamo descrivere i moti del sistema che come troveremo seguono traiettorie coniche, e individuare la periodicità dei moti limitati.

\subsection{Campo di forze centrale e moti piani}


Consideriamo un punto materiale in un generico campo di forze conservativo e scegliamo un sistema di coodrinate cartesiane con origine nel centro del campo di forza e assumiamo che il campo sia definito ovunque eccetto il centro, come nel caso della forza gravitazionale (per campi definiti anche nel centro l'adattazione è immediata).

La Lagrangiana del sistema è
\[
L(x,\dot{x})=\dfrac{1}{2}m||\dot{x}||^2 - V(||x||) \qquad \quad x\in\mathbb{R}^3\setminus\{0\}
\]

come sappiamo questa Lagrangiana è invariante per rotazioni attorno ad un qualsiasi asse e quindi il momento angolare $M=mx\times\dot{x}$ è conservato. Quindi il sistema avrà $3$ integrali primi oltre a quello dell'energia.

Dai corsi di fisica è noto che la conservazione del momento angolare implica che i moti siano piani e che la ``velocità areoale'' sia costante, dimostraiamo ora questi fatti.

\begin{proposizione}
In un campo di forze centrali:\begin{enumerate}[i.]
\item Ogni moto $t\mapsto x_t$ con momento angolare $M\neq 0$ si svolge nel piano ortogonale a $M$ contenente l'origine e non transita per l'origine.
\item Un moto $t\mapsto x_t$ con momento angolare nullo è radiale o costante.
\end{enumerate}
\end{proposizione}


\begin{dimo}
$ $\\
\begin{minipage}[H]{0.73\linewidth}
\begin{enumerate}[i.]
\item Abbiamo $mx_t\times\dot{x}_t=M$ $\forall t$ quindi se $M\neq 0$ il vettore $x_t$ (e anche il vettore $\dot{x}_t$) è ortogonale al vettore $M$ ed è pertanto contenuto nel piano ortogonale. $M\neq0\implies x_t,\dot{x}_t\neq0$.
\item Ogni moto radiale o costante avrà chiaramente momento angolare nullo, non dimostriamo il viceversa. \qedhere
\end{enumerate}
\end{minipage}
\begin{minipage}[H]{0.25\linewidth}
\flushright
\includegraphics[width=\linewidth]{img/motoCentralePiano.png}
\end{minipage}
\end{dimo}

\begin{corollario}
$\forall a\in\mathbb{R}^3\setminus\{0\}$ il sottoinsieme $\Pi_a$ dello spazio degli atti di moto
\[
\Pi_a=\{ (x,\dot{x})\in\mathbb{R}^2\setminus\{0\}\times\mathbb{R}^3 : x\perp a , \dot{x}\perp a \}
\]
è invariante.
\end{corollario}
\begin{dimo}
Per $(x,\dot{x})\in\Pi_a$ si ha $mx\times\dot{x}=M\parallel a$ ma $M$ è invariante quindi se in un punto di una soluzione $M\parallel a$ allora rimarrà parallelo lungo tutta la soluzione e quindi $\Pi_a$ è invariante. \qedhere 
\end{dimo}

L'insieme $\Pi_a$ è una sottovarietà di dimenione $4$ dello spazio degli atti di moto, diffeomorfa a $\mathbb{R}^2\times\mathbb{R}^2$, e quindi possiamo restringere le equazioni di Lagrange a $\Pi_a$.

Posso scegliere un riferimento in cui $a=e_3,$\footnote[1]{Le equazioni nello spazio degli atti di moto usando il riferimento $a=e_3$ sono $\dot{x}_i=v_i$ $\dot{v}_i=-\dfrac{1}{m}V'(||x||)\dfrac{x_i}{||x||}$ $i=1,2,3$ dove però $x_3=v_3=0$.} la restrizione delle equazioni di Lagrange di Lagrangiana $L$ a $\Pi_a$ sarà un sistema di equazioni di Lagrange nella stessa forma di quelle di partenza ma a $2$ gradi di libertà 

\[
\dot{x}_i=v_i\qquad \dot{v}_i=-\dfrac{1}{m}V'(||x||)\dfrac{x_i}{||x||} \qquad i=1,2
\]

\subsection{Riduzione del sistema}

Il sistema ristretto a $\Pi_a$ è ancora invariante per rotazioni e quindi possiamo applicare la riduzione alla Routh e ridurci allo studio di un sistema con un solo grado di libertà.

\begin{itemize}
\item Usando le coordinate polari per il piano $(x_1,x_2)$ la Lagrangiana diventa
\[
L(r,\theta,\dot{r},\dot{\theta})=\dfrac{m}{2}(\dot{r}^2+r^2\dot{\theta}^2)-V(r)\qquad r\in\mathbb{R}_+\quad \theta\in S^1
\]
Le coordinate polari sono giustificate dato che abbiamo visto che i moti con momento angolare non nullo non passano per l'origine
\item Fissiamo un valore $J\neq0$ di $p_\theta=\dfrac{\partial L}{\partial \dot{\theta}}=mr^2\dot{\theta}$ e quindi
\[
\dot{\theta}=\dfrac{J}{mr^2}
\]
\item La Lagrangiana ridotto è quindi
\[
L_J^R(r,\dot{r})=\dfrac{m}{2}\dot{r}^2-W_J(r) \qquad \quad W_J(r)=V(r)+\dfrac{J^2}{2mr^2}
\]
L'energia dei moti del sistema ridotto sarà $E(r,\dot{r})=\dfrac{m}{2}\dot{r}^2+V(r)$ e coinciderà con l'energia dei moti completo.

Dato che il sistema ridotto ha un solo grado di libertà possiamo costruirne il ritratto in fase a partire dal grafico dell'energia potenziale efficace.
\end{itemize}


\subsection{Sistema di Kepler ridotto}
Nel sistema di Kepler il potenziale è quello gravitazionale e quindi
\[
V(r)=-\dfrac{k}{r} \qquad W_J(r)=-\dfrac{k}{r}+\dfrac{J^2}{2mr^2} \qquad k>0, J\neq0
\]

Per $r$ grandi domina il primo termine mentre per $r$ piccoli domina il secondo.
La derivata prima del potenziale effice
\[
W_J'(r)=\dfrac{k}{r^2}-\dfrac{J^2}{mr^3}
\]
si annulla in un unico punto $r=r_J^c$ dato da
\begin{equation}
r_J^c := \dfrac{J}{mk}
\end{equation}

\begin{figure}[H]
\centering
\includegraphics[width=0.4\linewidth]{img/potenzialeEfficaceKepler.png}

$W_J(r)$
\end{figure}

In $r=r_J^c$ il potenziale assume il suo valore minimo $W_j^c:=W_j(r_j^c)=-\dfrac{k}{2r_j^c}$.

Dal ritratto in fase capiamo che vi saranno $3$ tipi di orbite del sistema ridotto:\begin{itemize}
\item se $E=W_J^c$ avremo un equilibrio in $r=r_j^c$.
\item Se $W_j<E<0$ avremo orbite chiuse con moti periodici che oscillano fra $r_-(E,J)$ e $r_+(E,J)$ soluzioni di $E=W_J(r)$.
\item se $E>0$ avremo orbite illimitate che vanno a $r\rightarrow+\infty$ per $t\rightarrow\pm\infty$ con una $r$ minima $r_-(E,J)$ soluzione di $E=W_J(r)$. 
\end{itemize}


\subsection{Ricostruzione del sistema di Kepler completo}

Per $J\neq0$ qualunque al moto del sistema ridotto $t\mapsto(r_t,\dot{r}_t)$ corrispondono i moti del sistema completo $t\mapsto(r_r,\theta_t,\dot{r}_t,\dot{\theta}_t)$ dove $\dot{\theta}_t=\dfrac{J}{mr^2_t}$ e $\theta_t=\theta_0+\int_0^t \dot{\theta}_s ds$.

Quindi avremo $3$ tipi di moti

\paragraph{Moti circolari:} Gli equilibri relativi $r(t)=r_J^c$ del sistema ridotto sono moti circolari di raggio $r_J^c$ percorsi con velocità angolare costante $\dot{\theta}=\dfrac{J}{m \left.r_J^{c}\right.^2}$.
Il senso di percorrenza dell'orbita dipendera dal segno di $J$.

\paragraph{Moti limitati:} I moti periodici del sistema ridotto corrisponderanno a moti del sistema completo con $r_-(E,J) \leq r \leq r_+(E,J)$. Il raggio oscillerà in modo periodico tra il minimo e il massimo e $\theta$ varia secondo una $\dot{\theta}_t=\dfrac{J}{mr^2}$ monotòna, con segno dato dal segno di $J$.

Non abbiamo ancora verificato però se il moto $t\mapsto(r_t,\theta_t)$ sia periodico e in caso dopo quanti giri le orbite si chiudano.

\paragraph{Moti illimitati:}I moti illimitati del sistema ridotto corrispondono a moti illimitati del sistema completo. Per energie $E\geq0$ il raggio raggiunge un minimo per $E=W_J(r)$ e va ad infinito sia nei tempi positivi sia in quelli negativi.

Non sappiamo però se il punto faccia più giri attorno al centro prima di andare ad infinito o quale sia l'angolo fra le direzioni asintotiche.

\subsection{Periodicità dei moti limitati del problema di Kepler}

Nei moti con $E\leq0$ e $J\neq0$ il moto della coordinata $r$ è periodica ma questo non implica che il moto complessivo del punto lo sia.

Chiamiamo $\Theta(E,J)$ la quantità di cui varia la coordinata $\theta$ in un periodo $T(E,J)$ della coordinata $r$, cioè se $r(t)=r(t+T)$ definiamo $\Theta(E,J):=\theta(t+T)-\theta(t)$.
Il moto complessivo è periodico soltanto se $\Theta(E,J)$ è commensurabile a $2\pi$, cioè se $n_1\Theta(E,J)=n_2 2\pi$ per $n_1,n_2\in\mathbb{N}$ primi fra loro e in quel caso il periodo sarà $n_1T$.

Ripetiamo quindi il ragionamento visto nella sezione \ref{subsezNewton1} per determinare il periodo dei moti del sistema ristretto.

Gli estremi del moto lungo $r$ sono
\[
r_{\pm}(E,J)=\dfrac{r_J^c}{1\mp\sqrt{1-\dfrac{E}{W_J^c}} }
\]
e questi sono anche i punti di inversione del moto lungo $r$.

Sapendo che $E$ è fissata lungo un moto e che per il sistema ridotto $E=\dfrac{1}{2}\dot{r}-W_J(r)$ troviamo che lungo una soluzione
\begin{equation}
\dot{r}_t=\sqrt{\dfrac{2}{m}\left(E-W_J(r)\right)}
\end{equation}
fino a quando $r$ non raggiune un punto di inversione, in cui cambia la monotònia di $t\mapsto r_t$.
Nell'intervallo in cui $t\mapsto r_t$ è invertibile la sua inversa soddisfa 
\[
\dfrac{dt_r}{dr}=\sqrt{\sfrac{\dfrac{m}{2}}{E-W_J(r_t)}}
\]
e quindi integrandola
\[
t_x=t_0+\sqrt{\dfrac{m}{2}} \int_{r_0}^{r} \dfrac{dr}{\sqrt{E-W_J(r)}}
\]
Quindi il periodo del moto periodico sarà il doppio del tempo necessario per arrivare ad un punto di inversione partendo dall'altro e
\begin{equation}
T(E,J)=\sqrt{2m} \int_{r_-}^{r_+} \dfrac{dr}{\sqrt{E-W_j(r)} } \qquad \forall E<0, J\neq0
\end{equation}
e risolvendo l'integrale\footnotemark[1]

\footnotetext[1]{Si risolve con tecniche standard ma il calcolo esplicito è un po' macchinoso quindi lo tralasciamo} si trova

\[
T(E,J) \sim \dfrac{\text{cost}}{|E^{\sfrac{3}{2}}}=:T(E)
\]

Possiamo quindi determinare l'incremento $\Theta(E,J)$ di $\theta$ in un periodo
\[
\Theta(E,J)=\int_0^{T(E)} \dfrac{J}{mr^2(t)} dt \rightarrow \Theta(E,J)=2\pi \qquad \forall E<0,J\neq0
\]
e quindi concludiamo che tutti del problema di Kepler limitati sono periodici, con periodo che dipende solo dall'energia.

\subsection{Traiettorie Kepleriane}

Adesso vogliamo trovare la forma delle traiettorie dei moti Kepleriani. Possiamo farlo poichè essendo $\dot{\theta}_t$ non nulla in tutti i moti con $J\neq0$ la $t\mapsto\theta_t$ sarà monotòna e quindi invertibile. Allora usando la sua inversa $t(\theta)$ potremo scrivere
\[
\theta\mapsto r(t(\theta))=:\tilde{r}(\theta)
\]
che ci darà la parametrizzazione $\theta\mapsto(\tilde{r}(\theta),\theta)$ delle traiettorie.

\begin{lemma}[Formula di Binet]
La parametrizzazione $\theta\mapsto\tilde{r}(\theta)$ delle traiettorie del problema di Kepler con momento angolare $J\neq0$ soddisfa
\[
\dfrac{d^2}{d\theta^2}\left(\dfrac{1}{\tilde{r}(\theta)}\right) +\dfrac{1}{\tilde{r}(\theta)}=\dfrac{1}{r_J^c}
\]
\end{lemma}
\begin{dimo}
La $t\mapsto r_t$ di un moto è soluzione dell'equazione di Lagrange di Lagrangiana $L_J^R$, esplicitamente
\[
m\ddot{r}=-\dfrac{k}{r^2}+\dfrac{J^2}{mr^3}
\]
Usando $r(t(\theta))=\tilde{r}(\theta)$ 
\begin{gather*}
\dot{r}=\dot{\theta}\dfrac{dr}{d\theta}=-\dfrac{J}{mr^2}\dfrac{d\tilde{r}}{d\theta}=-\dfrac{J}{m}\dfrac{d}{d\theta}\left(\dfrac{1}{\tilde{r}}\right)
\\
\ddot{r}=\dfrac{d}{dt}\left[ -\dfrac{J}{m}\dfrac{d}{d\theta}\left(\dfrac{1}{\tilde{r}}\right) \right]= -\dfrac{J}{m}\dot{\theta}\dfrac{d^2}{d\theta^2}\left(\dfrac{1}{\tilde{r} }\right)=-\dfrac{J^2}{m^2\tilde{r}^2}\dfrac{d^2}{d\theta^2}\left(\dfrac{1}{\tilde{r}}\right)
\end{gather*}
Inserendole in $m\ddot{r}=-\dfrac{k}{r^2}+\dfrac{J^2}{mr^3}$ troviamo
\[
-\dfrac{J^2}{m\tilde{r}^2}\dfrac{d^2}{d\theta^2}\left(\dfrac{1}{\tilde{r}}\right) = -\dfrac{k}{\tilde{r}}+\dfrac{J^2}{m\tilde{r}^2} \quad \implies \quad
\dfrac{d^2}{d\theta^2}\left(\dfrac{1}{\tilde{r}(\theta)}\right) +\dfrac{1}{\tilde{r}(\theta)}=\dfrac{1}{r_J^c}  \qedhere
\]
\end{dimo}

La formula di Binet è un'equazione differenziale non omogenea per $\dfrac{1}{\tilde{r}(\theta)}$ quindi il suo integrale generale sarà dato dalla somma dell'integrale generale dell'omogenea associata e di una soluzione particolare.

L'integrale dell'omogenea lo scriviamo come \[\dfrac{1}{\tilde{r}(\theta)} =A\cos{(\theta-\theta_p)} \qquad A\geq0\quad\theta_p\in S^1\] e sappiamo che $\dfrac{1}{r_J^c}$ è una soluzione particolare quindi le soluzioni dell'equazione della formula di Binet sono
\[
\dfrac{1}{\tilde{r}(\theta)}=A\cos{(\theta-\theta_p)}
\]
e, ponendo $A=\dfrac{e}{r_J^c}$ per una $e\in\mathbb{R}_+$, le traiettorie saranno quindi
\begin{equation}
\tilde{r}(\theta)=\dfrac{r_J^c}{1+e\cos{(\theta-\theta_p)}} \qquad\quad e\geq 0, \theta_p\in S^1
\end{equation}
Questa è l'equazione in coordinate polari delle coniche aventi un fuoco nell'origine ed eccentricità $e$. Per $e=0$ avremo cerchi con $r=r_J^c$, per $0<e<1$ ellissi, per $e=1$ parabole e per $e>1$ iperboli.

\begin{proposizione}
L'eccentricità delle traiettorie Kepleriane di energia $E$ e momento angolare $j\neq0$ è 
\begin{equation}
e(E,J)=\sqrt{1-\dfrac{E}{W_J^c}}
\end{equation}
La traiettoria sarà un cerchio se $E=W_J^c$, un'ellisse se $W_J^c<E<0$, un parabola se $E=0$ e un'iperbole se $E>0$.

Inoltre il semiasse maggiore delle traiettorie ellittiche sarà $a=\dfrac{k}{2|E|}$.
\end{proposizione}
\begin{dimo}
Da $\tilde{r}(\theta)=\dfrac{r_J^c}{1+e\cos{(\theta-\theta_p)}}$ $e\geq 0, \theta_p\in S^1$ la distanza minima di una traiettoria dall'origine è $\dfrac{r_J^c}{1+e}$.
Uguagliando la distanza minima a $r_-(E,J)$ troviamo l'equazione data per l'eccentricità. 
Per una traiettoria ellittica il doppio del semiasse maggiore sarà $r_-(E,J)+r_+(E,J)$ da cui la formula. \qedhere
\end{dimo}






























\end{document}
